% PM 트랙 Day 2 슬라이드 본문 — 손으로 쓴 정본(한 프레임 = 한 쪽). Day1_슬라이드_body.tex 와 같은 형식.
% 래퍼: Day2_슬라이드.tex(슬라이드판) / Day2_슬라이드_노트.tex(노트판, 우측에 \note).
% 화면에는 결론만, 설명·근거·예외는 각 프레임의 \note{}에. Day2_슬라이드.md 는 이 파일의 텍스트 미러다.
% 그림은 ../그림/PM_Day2/*.pdf (벡터). 교재 그림 번호와 파일명의 절 번호가 다를 수 있다 — 캡션은 교재 번호를 쓴다.

\newcolumntype{L}[1]{>{\raggedright\arraybackslash}p{#1}}
\newcommand{\ra}{$\rightarrow$}
\newcommand{\til}{\textasciitilde}
\newcommand{\keymsg}[1]{\par\vfill{\color{mDarkTeal}\bfseries #1}\par}
\newcommand{\figcap}[1]{\par\smallskip{\footnotesize\color{black!60}#1}\par}
\newcommand{\fig}[2][0.66]{\centering\includegraphics[height=#1\textheight,width=\textwidth,keepaspectratio]{#2}\par}
\newcommand{\subhead}[1]{{\color{mDarkTeal}\bfseries #1}\par}
\newenvironment{tbl}[2][\small]{\begin{center}\usebeamercolor[fg]{normal text}\setlength{\tabcolsep}{4pt}\renewcommand{\arraystretch}{1.15}#1\begin{tabular}{#2}\toprule}{\bottomrule\end{tabular}\end{center}}

% =====================================================================
% 도입
% =====================================================================
\begin{frame}{PM 트랙 Day 2 — 어제 결정한 것을 오늘 분해하고 숫자로 만든다}
\begin{itemize}
\item 프로젝트 매니지먼트 4일 특강 · 2일차 · 계획(Planning)
\item 사례는 어제와 같다 — 온담F\&B홀딩스 ONE ONDAM \textbf{P1} · 승인 168.0 / 집행 한도 156.0 / 계획 원가 145.8억
\item 하이브리드 전제 — 2주 스프린트 36회(설계 6 / 구현 20 / 통합·시험 10) + 게이트 G0\til{}G5 + FP 12,400점 고정가
\item Day1 다섯 문서(BC·헌장·등록부·RACI·프리모템)가 오늘 다섯 산출물(⑥\til{}⑩)의 \textbf{입력}이다
\end{itemize}
\keymsg{오늘의 한 문장 — 계획은 예측이 아니라 합의이며, 합의의 형태는 측정 가능한 기준선이다}
\note{표지. 어제 워크북 완성 예시본 5종(헌장 §4·§5·§6·§7·§8·§9·§12·§13, 리스크 13+2건, RACI 19·20행)을 읽어 왔는지 확인한다 — 읽지 않은 수강생은 워크북 §0.3 한 페이지 요약과 §0.4 "Day1에서 가져오는 것" 표를 4교시 전까지 읽도록 안내한다. Day1이 "결정되지 않은 것을 결정하는 문서"였다면 Day2는 "결정된 것을 분해하고 숫자로 만드는 문서"이며, 그래서 오늘의 워크북에는 계산(CPM·PV·EMV)이 있다. 교재 서문 참조. 예상 소요 3분.}
\end{frame}

\begin{frame}{Day1 착수 문서 5종 \ra{} Day2 계획 산출물 5종}
\fig[0.62]{fig_0_1_day1_to_day2}
\figcap{그림 0-1. 실선은 값이 그대로 옮겨지는 입력, 회색은 규칙·구조가 옮겨지는 입력, 파선은 "오늘 만든 것이 헌장 v1.0을 고친다" (G1 2026-04-30에 v1.1 확정)}
\keymsg{Day1 문서는 오늘 고쳐진다 — 그것은 오류가 아니라 계획 단계가 하는 일이다}
\note{교재 서문의 표와 그림 0-1. ⑥ 범위 기술서·WBS는 헌장 §5(포함 5개 = WBS 레벨 2, 제외 7개, 경계 조건 4개)·§6, ⑦ RTM은 헌장 §4 상위 요구 8개(요구 제공자 실명)와 §13 우선순위, ⑧ 일정은 헌장 §7 마일스톤(외부 고정 4개, 컷오버 창)·§9 제약과 R-01·R-10 근접성, ⑨ 원가는 BC §5 내역 8개와 3층 예산·RACI 19·20행, ⑩ 리스크는 13+2건·허용범위 초안·손익분기 66\%. 세 가지를 미리 말한다 — Day1 문서는 오늘 고쳐진다(WBS에서 헌장에 없던 산출물이 나오고 EMV 합이 8억을 넘는다) / 하이브리드 전제를 잊지 말라(릴리스 계획·스프린트 달력·게이트·계약 기준선의 네 층) / 숫자에는 값·단위·출처·분모. 예상 소요 4분.}
\end{frame}

\begin{frame}{오늘의 흐름}
\begin{columns}[T]
\begin{column}{0.47\textwidth}
\subhead{오전 — 이론 (제1\til{}3장)}
\begin{itemize}
\item \textbf{1교시} 90분 · 계획론 · 범위 기술서 · WBS · RTM · FP 기준선 — 제1\til{}2장
\item \textbf{2교시} 90분 · 일정 — CPM · float · 압축 · 몬테카를로 · 스프린트×게이트 — 제3장
\item \textbf{3교시} 90분 · 원가와 조달 — AACE · 3층 · PV 곡선 · 계약 유형 — 제4장
\end{itemize}
\end{column}
\begin{column}{0.51\textwidth}
\subhead{오후 — 실습과 리스크}
\begin{itemize}
\item \textbf{4교시} 90분 · \textbf{실습 ①} 범위 기술서·WBS · RTM 골격 — 워크북 §1\til{}2
\item \textbf{5교시} 75분 · 리스크 — 정량화 · 대응 · 허용범위 — 제5장
\item \textbf{6교시} 75분 · \textbf{실습 ②} 일정 기준선 · 원가 기준선 · 리스크 등록부 — 워크북 §3\til{}5
\item \textbf{마무리} 30분 · 학술 배경 · 읽을 자료 · Day3 예고 — 제6장
\end{itemize}
\end{column}
\end{columns}
\keymsg{오늘 만든 다섯 기준선이 Day3 실행·통제(EVM·변경 통제)의 자(尺)가 된다}
\note{오전 세 교시가 이론(제1\til{}4장 중 리스크 제외), 오후는 실습 ①(범위·RTM) \ra{} 리스크 이론 \ra{} 실습 ②(일정·원가·리스크)의 순서다. 점심은 2교시 뒤가 아니라 3교시 뒤임을 명시한다(Day1과 다르다 — 원가 이론을 오전에 끝내야 오후 실습 ②에서 원가 기준선을 쓸 수 있다). 리스크 이론(5교시)이 실습 ②(6교시) 바로 앞에 오는 것도 같은 이유 — 리스크 정량화의 잔여 EMV가 원가 기준선의 컨틴전시 배정표에 들어간다. 예상 소요 3분.}
\end{frame}

\begin{frame}{학습 목표 — 오늘이 끝나면}
\begin{columns}[T]
\begin{column}{0.5\textwidth}
\subhead{오전 이론}
\begin{enumerate}
\item 계획의 세 기능을 분리해 "틀렸으니 소용없다"가 왜 오류인지 \textbf{설명한다}
\item 하이브리드에서 \textbf{어느 층을 고정하고 어느 층을 여는가}를 P1으로 답한다
\item 산출물 지향 WBS · 양방향 RTM · FP 기준선의 관계를 말한다
\item CPM을 \textbf{손으로} 계산하고 float 소유권을 계약 문장으로 쓴다
\item 예산 3층에서 \textbf{누구의 돈인가}를 구분하고 PV를 작업 시점으로 배분한다
\end{enumerate}
\end{column}
\begin{column}{0.48\textwidth}
\subhead{오후 실습}
\begin{enumerate}\setcounter{enumi}{5}
\item 리스크를 확률·영향·EMV로 정량화하고 허용범위와 \textbf{대조}한다
\item 다섯 기준선을 온담 P1으로 \textbf{직접 작성}하고 서로 \textbf{검산}한다
\end{enumerate}
\end{column}
\end{columns}
\keymsg{정직하게 적을 셋 — 157.8\ra{}145.8의 산수 / 잔여 EMV 6.78 vs 8억 / 컨틴전시 10.2 ≈ P80}
\note{일곱 목표 중 앞의 다섯은 오전 이론(1\til{}3교시), 6은 5교시, 7은 두 실습이다. 오늘의 결론을 미리 말해 둔다 — 계획 단계에서 정직하게 적어야 할 세 가지(교재 "Day2를 마치며"): 승인 내역 일곱 항목 합 157.8에서 계획 원가 145.8로 가는 12.0의 산수, 대응 전 위협 EMV 17.44가 허용범위 8억의 두 배이고 대응 후 잔여 6.78도 여유가 1.22(R-01 하나의 잔여)뿐이라는 사실, 컨틴전시 10.2가 잔여 노출의 약 P80이며 그 너머의 꼬리는 돈이 아니라 구조(분리 컷오버·폴백·선행 릴리스)가 자른다는 사실. 이 셋은 나쁜 소식이 아니라 계획이 한 일이며 스폰서에게 결정을 요청하는 근거다. "계획이 좋은 소식만 담고 있다면 그것은 계획이 아니라 희망이다." 예상 소요 3분.}
\end{frame}

% =====================================================================
\section{1교시 (90분) — 계획론 · 범위와 요구사항}
% =====================================================================

\begin{frame}{이 교시에서 배우는 것}
\begin{itemize}
\item 계획의 세 기능과 기준선의 세 속성 — 왜 P1은 G1까지 4개월을 잠정 기준선으로 일하는가 [§1.1]
\item 하이브리드 계획의 층 — 케이던스는 고정하고 내용은 연다 · rolling wave · 릴리스 R0\til{}R5 [§1.2]
\item 세 표준(PMBOK 8판 4 도메인 · PRINCE2 PBP · ISO 21502 §7)과 분야별 정밀도 [§1.3\til{}1.4]
\item 범위 기술서 · WBS 규칙 · WBS 사전 · PBS와의 결합 [§2.1\til{}2.2]
\item 요구사항 3층과 RTM · MoSCoW · 범위 팽창과 한국 SI의 FP 기준선 · 온담 P1의 WBS와 214건 [§2.3\til{}2.5]
\end{itemize}
\keymsg{범위는 오늘 다섯 산출물 중 나머지 넷의 뼈대다 — WBS 코드가 일정·원가·형상에 같은 뿌리로 박혀야 한다}
\note{제1장과 제2장을 90분에 다룬다. 시간 배분은 §1.1 12분, §1.2 12분, §1.3\til{}1.4 10분, §2.1\til{}2.2 15분, §2.3 12분, §2.4 10분, §2.5 12분, 정리 5분. 제2장 후반(§2.4 FP·변경 대가)은 한국 SI 실무이므로 수강생 회사의 계약서 경험을 물으며 진행한다. §2.5 온담 WBS·214건은 4교시 실습 ①의 직접 입력이다. 예상 소요 2분.}
\end{frame}

\begin{frame}{계획은 예측이 아니라 합의다 — 세 기능과 세 기준선}
\fig[0.55]{fig_1_1_three_functions_pmb}
\figcap{그림 1-1. 합의 · 통제 기준 · 의사소통, 그리고 세 기준선이 같은 코드 체계(작업패키지 \ra{} 활동 \ra{} 시간대별 원가)로 묶여 PMB가 되는 구조}
\keymsg{기준선은 틀리기 위해 있다 — 틀린 정도를 재는 것이 통제다. 바뀌지 않는 것은 자의 존재이지 눈금이 아니다}
\note{§1.1 개념. 17개월 뒤 컷오버 날짜를 착수 시점에 맞힐 수 있는 사람은 없고, 145.8억이 정확히 145.8억으로 끝날 확률은 0에 가깝다 — 예측으로서의 계획은 첫 달에 틀리고 서랍에 들어간다. 계획의 세 기능: 합의(정미란·오정근·박세연·이재현이 같은 문장에 서명 — 가치는 정확성이 아니라 서명자의 수와 무게), 통제 기준(기준선 = 실적을 비교하는 자; 범위·일정·원가 세 기준선이 PMB로 통합), 의사소통(82명이 같은 그림 — Day1 임시조직론의 fragmentation for commitment). 그림 위의 두 오류 문장 — "계획이 틀렸으니 소용없다"는 합의와 통제 기준을 혼동한 것, "계획을 바꾸면 안 된다"는 통제 기준과 의사소통을 혼동한 것. 예상 소요 6분.}
\end{frame}

\begin{frame}{기준선의 세 속성 — 승인 · 버전 · 변경 통제}
\begin{itemize}
\item \textbf{승인된 것} — 초안은 기준선이 아니다. 동결(매월 마지막 수요일) 전은 백로그, 후는 기준선
\item \textbf{버전이 있다} — v1.0(G1) \ra{} v1.1(변경 승인 후) \ra{} v2.0(G2 재기준선). 비교 버전 없는 진척 보고는 보고가 아니다
\item \textbf{변경 통제 아래} — 허용범위 안은 PM의 것, 넘으면 예외 보고. 작업 계획(매주 바뀜) ≠ 기준선
\item \textbf{재기준선}은 이전 편차를 흡수해 지운다 — 남발하면 편차를 없애는 게 아니라 안 보이게 하는 것
\end{itemize}
\keymsg{세 오류 — 기준선 없이 진척 보고("무엇의 42\%인가") / 한 번 세우고 끝까지 고수 / 세 기준선을 따로 만들기}
\note{§1.1 "기준선의 세 가지 속성"과 "실무에서 어떻게 틀리는가". 기준선이 없으면 진척률은 담당자의 느낌이고, 느낌은 종료 3개월 전까지 항상 80\%다. 변경이 승인되었는데 기준선을 갱신하지 않으면 실적은 승인된 범위를, 기준선은 옛 범위를 가리켜 둘의 차이는 편차가 아니라 소음이 된다 — 반대 오류는 재기준선 남발. 어떤 변경이 기준선을 갱신하고 어떤 변경이 편차로 남는가의 규칙이 Day3 변경 관리의 핵심. 세 번째 오류(범위는 사업팀 엑셀, 일정은 PMO 도구, 원가는 재경 회계 코드)의 기술적 방어가 §2.2 코드 체계 — WBS 코드가 일정의 활동 코드와 원가의 계정 코드에 같은 형태로 박혀야 한다. 예상 소요 5분.}
\end{frame}

\begin{frame}{온담으로 읽기 — G1 통과 조건은 세 기준선이고, 그 전 4개월은 잠정 기준선이다}
\begin{itemize}
\item G1(2026-04-30) 통과 조건: \textbf{FP 기준선 확인서}(1,180건 ↔ 12,400점, 3인 서명) · \textbf{원가·일정 기준선} · \textbf{조달 전략 확정}
\item 착수 01-05 \ra{} G1까지 \textbf{거의 4개월을 기준선 없이} 일한다 — 설계 스프린트 S1\til{}S6(\til{}03-27)
\item 결함이 아니라 하이브리드의 정상 형태 — 1,180건 확정 전에 원가 기준선을 서명하면 서명과 동시에 틀린다
\item 그래도 "무엇에 비해" 진척을 말해야 한다 \ra{} \textbf{잠정 기준선}: 헌장 §7 마일스톤 + §8 예산 3층
\item 측정 단위: 요구사항 확정 건수(214/396/570), 인터페이스 문서화 본수(13본)
\end{itemize}
\keymsg{워크북 §3의 심화 과제 — G1 이전 4개월의 잠정 기준선을 무엇으로 재고 누가 승인하는가}
\note{§1.1 "온담의 사례로 읽기". 조달 전략 확정이 통과 조건에 있는 이유 — 상용SW 23.4·인프라 23.9·데이터 전환 11.4 같은 별도 계약이 언제 어떤 유형으로 체결되는가가 기준선의 시간대별 배분에 들어가야 하기 때문. 설계 스프린트 날짜(2026-01-05\til{}03-27)는 [추가 설정]. 생각해 볼 질문 3 — "P1의 설계 스프린트 6회 동안 스폰서에게 진척을 보고해야 한다. 무엇의 몇 \%라고 말할 것인가"를 수강생에게 던진다. 예상 소요 4분.}
\end{frame}

\begin{frame}{하이브리드 계획의 층 — 케이던스는 고정하고 내용은 열어 둔다}
\fig[0.56]{fig_1_2_hybrid_layers}
\figcap{그림 1-2. 위 네 층(게이트 · 계약 범위 기준선 · 스프린트 케이던스 · 원가 기준선)은 착수 시점에 고정, 아래 두 층(릴리스 구성 · 스프린트 내용)은 정보가 생길 때 결정}
\keymsg{"섞는다"는 말은 아무것도 설명하지 않는다 — 어느 층을 예측형으로, 어느 층을 적응형으로 계획하는가에 답하라}
\note{§1.2 개념. 되돌림 비용(Day1 2.6절)이 계획의 형태를 결정한다 — 되돌릴 수 없는 결정 앞에서는 상세하게 미리, 되돌릴 수 있는 것은 가까운 것만 상세히. 예측형의 대가는 앞단의 시간(설계 성숙도 낮은 상태의 예측형 계획은 "정밀하게 틀린다"), 적응형의 대가는 예측 가능성. P1의 층: 게이트·마일스톤(날짜 단위 고정), 범위 기준선(FP 단위, 월 1회 동결), 릴리스 계획(분기 단위, 예측+적응), 스프린트 케이던스(36회 시작·종료일 고정), 스프린트 내용(백로그, 스프린트 계획 회의), 원가 기준선(월별 PV). 게이트는 스프린트 경계에 놓여야 한다 — G2(11-27)는 S24 중간에 떨어지는데, 이것이 2교시 3.5절의 문제다. 예상 소요 6분.}
\end{frame}

\begin{frame}{Rolling wave와 릴리스 계획 — 36 스프린트 안에서 무엇을 미리 정하는가}
\begin{tbl}[\footnotesize]{L{0.1\textwidth}L{0.17\textwidth}L{0.24\textwidth}L{0.4\textwidth}}
\textbf{구간} & \textbf{스프린트} & \textbf{게이트·마일스톤} & \textbf{미리 정하는 것} \\ \midrule
설계 & S1\til{}S6 (\til{}03-27) & G1 04-30 & 요구 확정 순서, 인터페이스 13본 전진 배치(R-06·O-01) \\
구현 & S7\til{}S26 (\til{}01-03) & T0 · 부속서 · 규제 ① · G2 11-27 & 릴리스 R1\til{}R4(포털 선행 R3), 동결 12회, 규제 R2 S13\til{}S18, 마스터 동결(R-10) \\
통합·시험 & S27\til{}S36 (\til{}05-23) & 리허설 · MC · G3 · 컷오버 · G4 & 리허설 3회 일자, UAT 종료 기준, 컷오버·롤백, 안정화 S31\til{}S34 · 이관 S35\til{}S36 \\
\end{tbl}
\begin{itemize}
\item Rolling wave — 다음 게이트까지는 활동 수준, 그 뒤는 \textbf{계획 패키지}(산출물·총 예산·기간만)
\item 릴리스 R0\til{}R5 — R0 설계\ra{}G1 / R2\ra{}규제 ① / R3 포털 선행(2026-Q4, O-02)\ra{}G2 / R4 = 컷오버(02-26\til{}28)\ra{}G3 / R5 컷오버 후(옴니 고도화·BI, §13)\ra{}G4
\end{itemize}
\keymsg{세 오류 — 스프린트를 폭포수의 단위로 / 게이트를 스프린트 옆에(G2 준비는 S23\til{}S24 백로그에) / FP 고정가와 적응형 범위의 충돌을 안 다루기}
\note{§1.2 "Rolling wave와 릴리스 계획", "온담의 사례로 읽기", "실무에서 어떻게 틀리는가". 스프린트 달력은 [추가 설정 — 정본의 설계 6/구현 20/통합·시험 10과 마일스톤에 정합]. 정본은 통합·시험을 "2027-01\til{}02"로 적지만 10회 × 2주 = 20주이므로 컷오버 후 안정화·이관(S31\til{}S36)까지 포함해 읽는다. G0 시점의 상세 계획은 S6까지, 구현 20회는 FP 약 7,400점 [추가 설정, 12,400 × 60\%]의 계획 패키지. 릴리스 계획이 없으면 헌장 §13의 우선순위는 문장으로만 남고 범위를 조정해야 하는 순간 무엇을 빼야 하는지 아무도 모른다. FP 고정가와 적응형 범위의 충돌 장치 = 월 1회 동결 + ±3\% 허용범위 + FP 재산정 트리거(§2.4). 예상 소요 6분.}
\end{frame}

\begin{frame}{표준에서는 — PMBOK 8판 네 도메인 · PRINCE2 7 PBP 네 단계 · ISO 21502 §7}
\begin{tbl}[\footnotesize]{L{0.17\textwidth}L{0.36\textwidth}L{0.4\textwidth}}
\textbf{표준} & \textbf{계획의 구조} & \textbf{Day2에 주는 것} \\ \midrule
\textbf{PMBOK 8판} & Planning Focus Area · Scope / Schedule / \textbf{Finance} / Risk 성과영역 & QC가 Scope로 \ra{} 인수 기준은 범위의 일부. Cost \ra{} Finance(자금 조달·현금흐름). 조달은 부록 X4로 \\
\textbf{PRINCE2 7} & Plans practice · \textbf{PPD \ra{} PBS \ra{} PD \ra{} PFD}, 계획 3수준, 단계 경계 승인 & 활동이 아니라 산출물(명사)에서 시작. PBS = 관리 산출물 포함. 게이트 G1·G2·G3 = 단계 경계 \\
\textbf{ISO 21502:2020} & §7.2 planning · §7.4 scope · §7.6 schedule · §7.7 cost · §7.8 risk & 통합된 기준선, 변경 통제, 컨틴전시와 관리예비비의 \textbf{구분}. 계약서가 인용하는 조항 \\
\end{tbl}
\begin{itemize}
\item 실무적 결합 — \textbf{PBS를 먼저 그리고 WBS의 상위 레벨로 쓰고, 작업패키지 수준에서 활동을 붙인다}
\end{itemize}
\keymsg{어느 표준도 컨틴전시의 크기를 정해 주지 않는다 — "표준에 따라 10\%"는 없는 말이다}
\note{§1.3. PMBOK 8판 Planning Focus Area의 프로세스 개수는 2차 출처 간 19 또는 20으로 갈리므로 판서하지 말고 인쇄본을 보라. 조달이 부록으로 내려갔다는 것은 4.4절의 계약 유형·변경 대가·견적 검산을 PMBOK에서 배울 수 없다는 뜻 — 출처는 계약 실무(FIDIC·NEC4·국가계약법·SW사업 대가산정 가이드). PRINCE2 PBP 4단계: Project Product Description(P1은 "가동 중인 플랫폼과 운영 인수 상태"), Product Breakdown Structure(명사, 관리/전문 산출물 구분), Product Description(WBS 사전 + RTM 인수 기준에 가까움), Product Flow Diagram(활동 네트워크의 뼈대). ISO 21502 조항 번호는 인쇄 전 원문 대조. 흔한 오류 넷: 프로세스 개수 판서 / PBS를 "WBS의 영국식 이름"으로 / ISO 21500으로 잘못 부르기 / "표준에 따라 10\%". 예상 소요 6분.}
\end{frame}

\begin{frame}{분야별 대조 — 무엇을 언제 얼마나 정밀하게 미리 정하는가}
\begin{tbl}[\scriptsize]{L{0.14\textwidth}L{0.2\textwidth}L{0.2\textwidth}L{0.17\textwidth}L{0.17\textwidth}}
\textbf{축} & \textbf{SI} & \textbf{건설·EPC} & \textbf{제조 NPD} & \textbf{제약·규제} \\ \midrule
범위 기준선 & FP 기준선 + 요구 목록. \textbf{계약 문서}가 정본 & 도면·시방서·BOQ. AACE Class가 계약 유형 결정 & PRD + BOM. 게이트마다 동결 상승 & 임상 프로토콜. 규제 승인 시 고정 \\
일정 정밀도 & 릴리스·스프린트. 활동 CPM은 \textbf{통합·전환 구간만} & 활동 단위 CPM 전체 & 금형·부품 리드타임 역산 & 규제 제출·심사 시계 역산 \\
컨틴전시 결정 & 요구 불안정, 인터페이스, 데이터 품질 & 설계 성숙도, 지반, 물가 & 수율, 부품 단산, 설계 변경 & 추가자료 라운드 수 \\
가장 자주 틀리는 곳 & 인터페이스·데이터 전환을 "우리 일"로 안 잡음 & Class 5 견적으로 LSTK & 게이트 사이 설계 변경 & 심사 라운드 0회 가정 \\
\end{tbl}
\keymsg{공유하는 규율 — 되돌릴 수 없는 결정 앞에 게이트를 두고, 거기까지는 상세히, 그 뒤는 굵게. 정밀도를 정하는 것은 방법론이 아니라 계약이다}
\note{§1.4. SI의 원가 기준선은 두 얼굴 — 발주사는 계약 금액 68.2와 지급 일정, 수행사는 자기 인건비; 발주사 PM은 수행사 원가를 못 보므로 발주사 EVM은 FP 진척(범위)과 지급 마일스톤(원가)으로 구성된다. EPC의 규율은 "Class 2 이전에 총액을 확정하지 않는다"(Merrow 2023 — 설계 성숙도 부족 + 경쟁입찰 압박에서 LSTK 성과가 나빠진다) — P2 이천 물류센터(194억)가 바로 이 분야이며 T0(2026-04-24) 후 선급금 29.6억과 함께 되돌릴 수 없다. NPD의 원가 기준선 절반은 양산 단가 목표(target cost). 제약은 P50과 P80의 차이가 다른 분야보다 크고 "라운드 수 시나리오"로 표현 — 온담의 규제 기한과 P4 규제 버퍼 22영업일은 이 논리의 축소판. 그림 1-4(정밀도 곡선)는 시간이 남으면. 예상 소요 4분.}
\end{frame}

\begin{frame}{범위 기술서 — 헌장 §5를 "어디까지가 완료인가"로 상세화한다}
\begin{itemize}
\item 다섯 항목: \textbf{산출물 범위 기술} · \textbf{인수 기준}(측정 가능 + 판정자) · \textbf{제외}(담당 프로젝트·담당자 붙여서) · \textbf{가정과 제약}(검증 상태 갱신) + 온담의 \textbf{경계 조건}
\item 인수 기준이 없으면 헌장의 복사본이다 — "안정적으로 동작" 대신 "UAT N건 통과, 0.4건/FP 이하, 피크 5,100 케이스/h"
\item 제외 항목에 담당이 없으면 제외가 아니라 누락 — 누구도 하지 않는 일은 결국 프로젝트로 돌아온다
\item 통합 재고의 인수 기준은 \textbf{조건부}: 매장·공장 15분 이내 / 센터는 3PL 계약 변경 시 15분, 미변경 시 야간 배치 2회(R-03)
\item 가맹 POS SW 배포의 경계: 호환 기종 기준, 비호환 기종은 배포 제외 + 수기 폴백, 전수 조사(2026-04) 결과가 v1.1 입력(R-13)
\end{itemize}
\keymsg{이 문장이 없으면 컷오버 다음 주 월요일에 90개 매장이 "P1이 우리 단말을 망가뜨렸다"고 말한다}
\note{§2.1. PMBOK 8판이 QC를 Scope 도메인에 넣었으므로 인수 기준은 선택 항목이 아니라 범위 정의의 일부. PRINCE2 PPD는 품질 기대와 인수 책임자를 더 강하게 요구하며, 헌장 §6이 산출물마다 인수자·인수 시점을 적은 것이 그 영향. 경계 조건 넷: P2 설비 인터페이스(규격은 P1, 구현·시험은 P2), 3PL 연동(명세와 P1측 API는 P1, 협상은 물류본부), 규제 아키텍처(반영은 P1, 해석은 P4), 운영 문서(기술 문서는 P1, 절차 문서는 P5). 야간 배치 시각(06:00·18:00)과 호환 기종 사양은 [추가 설정]. 흔한 오류 다섯: 헌장 복사본 / "고객 만족" / 제외에 담당 없음 / 가정 검증 상태 미갱신 / "기타 협의된 사항"(§2.4). 예상 소요 6분.}
\end{frame}

\begin{frame}{WBS — 100\% 규칙 · 산출물 지향 · 8/80 · 통제 계정 · 코드 체계}
\fig[0.56]{fig_2_2_wbs_rules}
\figcap{그림 2-1. 왼쪽 산출물 지향 트리는 "POS 결제 모듈이 끝났다"를 말할 수 있고, 오른쪽 활동 지향 트리는 "시험" 상자 하나에 모든 산출물이 섞여 무엇이 끝났는지 말할 수 없다}
\keymsg{프로젝트 관리 자체도 범위다 — PMO 9.6억(1.6)이 어느 상자에도 없으면 원가 기준선은 145.8이 되지 않는다}
\note{§2.2 개념. 100\% 규칙(ISO 21511:2018, MIL-STD-881F — 원문은 규칙을 강제하기보다 WBS의 목적에서 도출한다)은 두 방향: WBS에 없는 일은 범위 밖(크립 방어), 범위 안의 일은 반드시 어딘가에(누락 방어). 8/80은 절대 규칙이 아니라 "보고 주기의 1\til{}2배" — 2주 스프린트의 P1에서는 "작업패키지 하나가 스프린트 하나 안에 끝난다". 통제 계정(레벨 2\til{}3)에 CAM이 지정되어 PV·EV·AC에 책임 — Day3 EVM의 뼈대. 코드 체계: 1.3.3의 활동은 1.3.3-01, 원가 계정도 형상 항목도 같은 뿌리. 흔한 오류 여섯: 조직도로 그리기 / 활동 지향 / PM을 WBS에서 빼기 / 사전을 안 쓰기 / 코드 불일치 / 수행사 WBS를 그대로 받기(발주사 통제 단위 — 데이터 전환·상용SW·인프라·감리·PMO — 가 빠져 있다). 예상 소요 6분.}
\end{frame}

\begin{frame}{WBS 사전과 PBS \ra{} WBS 결합 — 활동은 마지막에 온다}
\begin{columns}[T]
\begin{column}{0.5\textwidth}
\subhead{WBS 사전 = RTM과 원가 기준선을 잇는 고리}
\begin{itemize}
\item 항목: 코드·명칭·설명·\textbf{인수 기준}·담당(CAM)·선행·후행·자원·기간·\textbf{원가 추정}·가정·품질·\textbf{관련 요구사항 ID}
\item PRINCE2 Product Description과 같은 자리
\item 사전 없는 WBS = 상자 이름의 목록. 상자 안은 담당자 머릿속에
\end{itemize}
\end{column}
\begin{column}{0.48\textwidth}
\subhead{결합 절차 다섯 단계}
\begin{enumerate}
\item 헌장 §6 산출물(명사)로 PBS 레벨 1\til{}2. 관리 산출물은 별도 가지
\item 구성 산출물로 분해(레벨 3) — 여기까지 PBS
\item 레벨 3 이하에서 작업패키지 정의(8/80)
\item 작업패키지에 WBS 사전
\item 산출물 사이의 의존(PFD) — 제3장 네트워크의 뼈대
\end{enumerate}
\end{column}
\end{columns}
\keymsg{활동에서 시작하면 WBS는 "누가 무엇을 한다"의 조직도가 되고, 산출물에서 시작하면 "무엇이 나온다"의 인도 구조가 된다}
\note{§2.2 "WBS 사전"과 "PBS와 WBS의 관계". 작업패키지의 "관련 요구사항 ID" 열이 RTM의 요구 \ra{} 작업패키지 추적을, "원가 추정" 열이 원가 기준선의 시간대별 배분을 가능하게 한다. WBS(작업의 분해, 산출물 지향 권장·활동 지향 허용)와 PBS(정의상 산출물만)의 차이를 짚고, 실무 결합은 PBS를 상위 레벨로. 워크북 §1 항목 10의 우선순위 규칙 — 43개 카드를 다 쓰려다 아무것도 안 쓰지 말고 외부 의존·임계경로·발주사 담당부터. 예상 소요 4분.}
\end{frame}

\begin{frame}{요구사항 3층(29148)과 RTM의 양방향 추적}
\fig[0.56]{fig_2_3_req_hierarchy_rtm}
\figcap{그림 2-2. 헌장 §4 8개 \ra{} StRS 214 \ra{} SyRS 396 \ra{} SRS 570과 RTM의 열. 앞으로(요구 \ra{} 시험)가 커버리지 100\%, 뒤로가 gold plating 탐지기}
\keymsg{한국 SI의 "요구사항정의서"는 어느 층인지 불분명하다 — 그래서 검수 때 "RFP의 이 요구가 어느 화면에서 충족되었는가"에 답하지 못한다}
\note{§2.3. StRS(이해관계자의 언어) / SyRS(시스템 경계와 인터페이스) / SRS(기능·데이터·화면·비기능). 제안요청서는 StRS와 SyRS가 섞이고, 수행사 요구사항정의서는 SRS에 가까우며, 둘 사이의 대응은 아무도 만들지 않는다 — 감리(정하늘)가 지적하는 RTM 커버리지 미달이 정확히 이것(프리모템 원문 15번). 29148의 좋은 요구사항 특성 중 가장 자주 무너지는 것은 검증 가능성("편리하게")과 단일성("발주·마감·출고·배송을 한 화면에서" = 네 요구). RTM 최소 열 11개(요구 ID·상위 요구·출처 실명·MoSCoW·릴리스·WBS·설계 요소·FP·시험 케이스·검수 확인·상태). 두 번째 기능은 변경 영향 분석 — Day3 변경요청서는 RTM 없이는 감으로 한다. StRS 배분 건수는 [추가 설정]. 예상 소요 6분.}
\end{frame}

\begin{frame}{MoSCoW · Kano · 비기능 요구 — FP에 없는 절반}
\begin{itemize}
\item \textbf{MoSCoW} — Must 비율 제한(R-07 완화: \textbf{Must ≤ 60\%}, 90\%면 조정할 것이 없다) + \textbf{Won't 명시}(안 적으면 Could로 돌아온다). 헌장 §13 이연 가능 = Should/Could 경계
\item \textbf{Kano} — 당연(결제 속도) / 일원(발주 리드타임) / 매력(픽업 예약) 품질. MoSCoW의 근거이자 릴리스 순서의 근거(매력은 시간이 지나면 당연이 된다)
\item \textbf{비기능(NFR)} — 헌장 §4의 8개 중 5개가 비기능. FP로 측정되지 않으므로 고정가 범위에서 보이지 않는다 — 12,400점에 "가용성 99.5\%"는 없다
\item 시험이 늦게 만들어진다(부하 3.1억·보안 2.4억은 통합 단계, 기준은 설계 단계에) / 다대다 추적
\item ISO/IEC 25010 품질 특성(2023 개정, 안전성 추가) + fit criterion(측정 방법·합격선)이 RTM의 비기능 행
\end{itemize}
\keymsg{흔한 오류 — Must 90\% / Won't 없음 / "성능은 충분해야 함" / 제공자 열이 부서명 / 검수 직전에 역으로 만든 RTM(추적이 아니라 사후 짝짓기)}
\note{§2.3 "우선순위"와 "비기능 요구의 위치". Must ≤ 60\%는 DSDM의 권고를 P1 계약 구조에 맞춘 것(워크북 §2.2). 매장 직원에게 결제 속도는 당연 품질, 가맹점주에게 발주 리드타임(38h\ra{}12h)은 일원 품질, 앱 사용자에게 픽업 예약은 매력 품질. Kano는 Product 트랙(Day5 이후)에서 다시. 비기능 5개: 성수기 배포 창, 가용성 99.5\%, 규제 요건, 자체 운영 아키텍처, SAP 마이그레이션 고려 구조. RTM을 엑셀로 만들어 변경마다 수작업 갱신하다 세 달 뒤 포기하는 것도 흔한 오류. 워크북 §2.6 해설 — Must가 61\%인 상태를 그대로 적고 "G1까지 오정근이 조정"이라고 쓰는 것이 §13을 작동시키는 방법(Senior User가 Should로 내릴 항목을 스스로 고른다). 예상 소요 5분.}
\end{frame}

\begin{frame}{범위 팽창의 세 경로와 한국 민간 SI의 범위 통제 구조}
\begin{columns}[T]
\begin{column}{0.5\textwidth}
\begin{itemize}
\item \textbf{크립} — 원인은 \textbf{문턱이 너무 높은 것}. 작은 변경도 기록: 동결 + ±35건 \ra{} FP 재산정
\item \textbf{Gold plating} — RTM 역추적이 탐지기
\item \textbf{"기타 협의된 사항"} — 무한 범위 = 무한 리스크. 한정이 발주사에게도 이롭다
\end{itemize}
\end{column}
\begin{column}{0.48\textwidth}
\begin{itemize}
\item \textbf{FP 기준선 확인서} — 1,180건 ↔ 12,400점 ↔ 68.2억. FP는 \textbf{기능만} 잰다 \ra{} 99.8억은 별도 계약
\item 계약 문서 우선순위(계약서 / 특수조건 / 과업지시서 / RFP / 제안서)는 \textbf{조항이 있을 때만}
\item 변경 대가 두 층 — \textbf{변경협의체}(대가) + \textbf{CCB}(예산). 민간엔 과업심의위원회 없음
\end{itemize}
\end{column}
\end{columns}
\keymsg{CCB가 승인해도 협의체가 대가에 합의 안 하면 수행사는 안 하고, 협의체가 합의해도 CCB가 예산을 안 주면 발주사는 지급 못 한다}
\note{§2.4. 그림 2-3(fig\_2\_4\_scope\_control\_si)은 시간이 남으면 띄운다. 입찰 시점 12,400점은 RFP 기준 개략 계수, G1 확인서가 정밀 계수 — 차이 ±3\% 초과 시 계약 금액 조정(정본 "분기 1회 또는 FP 누적 변동 ±3\% 초과 시"). 55만 원/FP가 어떤 보정을 거친 값인지 계약서 부속서에서 확인. 이재현 PM이 "발주사 미결정 대기와 요구사항 추가는 변경대가 대상"이라고 밝힌 것은 두 층 모두에 기록을 남기겠다는 뜻이고, 발주사 PM은 대응하는 자기 기록(결정 요청 응답 SLA 5영업일, RACI)을 가져야 한다. 수강생 회사의 마지막 계약서에 "기타 협의된 사항"이 있었는지, 그 문장으로 무엇이 들어왔는지 묻는다(생각해 볼 질문 1). 예상 소요 6분.}
\end{frame}

\begin{frame}{온담 P1 WBS — 다섯 산출물이 레벨 2가 되고 관리·통합 검증이 1.6이 된다}
\begin{tbl}[\scriptsize]{L{0.04\textwidth}L{0.22\textwidth}ccL{0.56\textwidth}}
\textbf{L2} & \textbf{산출물} & \textbf{FP} & \textbf{억} & \textbf{L3 작업패키지 (43개)} \\ \midrule
1.1 & POS·매장 운영 시스템 & 3,100 & 25.65 & 8 — 설계서 · 거래·운영 모듈 · PG 7본 · 단말 · 가맹 조사 · 파일럿 · 배포 \\
1.2 & 주문 허브 & 2,900 & 15.95 & 5 — 원장 모델 · 코어 · 앱·배달 · B2B EDI · 옴니(기본 R4/고도화 R5) \\
1.3 & 통합 재고 관리 & 2,300 & 12.65 & 5 — \textbf{1.3.1} 규격·협상 지원 · 코어 · \textbf{1.3.3} 3PL·폴백 · \textbf{1.3.4} P2 시험 · 공장 \\
1.4 & 가맹 발주 포털 & 1,500 & 8.25 & 4 — 설계 · 모듈 · 프로토타입·선행 릴리스 · 388점 온보딩 \\
1.5 & 통합 계층·전환·\textbf{공통 기반} & 2,600 & 64.00 & 10 — API GW · MDM · SAP 애드온 · 47본 · \textbf{인프라} · 매핑 · 정본(G2) · 리허설 · 전환 \\
1.6 & 프로젝트 관리·통합 검증 & — & 19.30 & 11 — 기준선 · PMO · RTM · 시험 3종 · UAT · 감리 · \textbf{컷오버(분리)} · 이관·인수 · 종료 \\
\textbf{합} & & \textbf{12,400} & \textbf{145.80} & 43 = 8+5+5+4+10+11 · 발주 담당 10 \\
\end{tbl}
\keymsg{1.3.1이 없으면 한동석은 "자료를 안 줘서 협상을 못 했다"고 말한다 — WBS가 리스크 대응의 형태를 정한다}
\note{§2.5 "WBS 레벨 1\til{}2"와 "레벨 3과 경계 조건의 분리", 워크북 §1.5 항목 7\til{}9 [추가 설정 — 구조와 코드는 워크북 예시본]. 인프라를 별도 L2가 아니라 1.5에 둔 이유: 헌장 §5의 다섯 산출물이 그대로 L2가 되어야 헌장·범위 기술서·WBS·FP 기준선이 같은 상위 구조를 갖고, 인프라는 API GW·MDM·SAP 애드온과 같은 층(다섯 산출물이 공유하는 기반)이다. 1.6은 100\% 규칙이 요구하는 관리·횡단 인도물 — 통합·성능시험 9.7억, UAT, 감리, 컷오버, PMO 9.6억, 2023 감사 지적에 대한 답인 소스·형상 인수 — 이며 각 산출물에 분산하면 G3 준비의 소유자가 사라진다. 1.3.1(재고 원장 규격, P2·3PL 협상 지원 자료)은 T0 2026-04-24 전에 넘겨야 하는 가장 이른 외부 인도물. R-01 완화(매장·주문 1차 / 물류 2차 분리)도 WBS에 나타난다 — 1.3.4(P2 수신·시험)와 1.6.9(컷오버, 분리 옵션)를 따로 두어 물류 부분만 떼어 낼 수 있게 했다. 경계 조건은 범위 밖을 적는 것이 아니라 범위 안의 마지막 조각(1.3.1 규격, 1.3.3 명세·P1측 API·폴백)을 적는 것이다. 예상 소요 6분.}
\end{frame}

\begin{frame}{상위 요구 8개 \ra{} StRS 214건 — 절반이 비기능이고, 요구 제공자가 곧 인수자다}
\begin{tbl}[\scriptsize]{L{0.42\textwidth}L{0.17\textwidth}cL{0.11\textwidth}L{0.1\textwidth}}
\textbf{상위 요구 (헌장 §4)} & \textbf{제공자} & \textbf{StRS} & \textbf{성격} & \textbf{WBS} \\ \midrule
1 단일 재고 원장 15분 · 2 단일 주문 원장·옴니 & 박세연·한동석 · 오정근 & 31 · 46 & 기능 & 1.3 · 1.2 \\
3 가맹 포털 한 화면(발주·마감·출고·배송) & 서정필, 오정근 & 38 & 기능 & 1.4 \\
4 성수기 배포 창 · 5 가용성 99.5\%·롤백 & 오정근 · 박준영 & 9 · 22 & 비기능 & 1.1·1.6 · 1.5·1.6 \\
6 개인정보 142개·망분리·전자금융 · 7 자체 운영·문서 47본 · 8 SAP 분리 구조 & 최영주 · 박세연 · 박세연 & 27 · 24 · 17 & 비기능·규제 & 1.5·1.6, 1.5·1.6, 1.5 \\
\end{tbl}
\begin{itemize}
\item 합계 214 = 기능 115 / 비기능 99 [추가 설정] — FP 12,400점은 기능을 주로 잰다. 서정필은 1.4 검수 참여자, 박준영은 SAC 판정권자
\end{itemize}
\keymsg{워크북에서 결정할 것 — 인프라를 1.5 안에 둘지 별도 L2로 뺄지, 통합 검증을 1.6에 모을지 분산할지, 214건 배분, Must 비율}
\note{§2.5 "상위 요구 8개 \ra{} StRS 214건" [추가 설정 — 건수는 예시]. 214 \ra{} 396 \ra{} 570(1 : 1.85 : 2.66)은 정본 값이며 설계 스프린트 6회의 산출물이 이 분해다. G1의 FP 기준선 확인서는 SRS 570건(과 SyRS)을 계수한 결과가 12,400점과 ±3\% 안인지 확인하는 문서. 범위 기준선에서 FP만 보면 절반이 안 보인다 — 생각해 볼 질문 4 "비기능 99건은 FP 12,400점의 어디에 있는가. 없다면 그것의 범위 기준선은 무엇인가". 4교시 실습 ①(워크북 §2)에서 상위 요구 3(가맹 포털)의 StRS를 예시로 채운다. 예상 소요 4분.}
\end{frame}

\begin{frame}{1교시 핵심 정리}
\begin{enumerate}
\item 계획은 예측이 아니라 \textbf{합의} — 세 기능(합의·통제 기준·의사소통)을 분리하면 두 오류가 보인다. 기준선은 승인·버전·변경 통제의 세 속성
\item 하이브리드는 "어느 층을 고정하고 어느 층을 여는가"에 답한다 — 케이던스·게이트·계약 기준선은 고정, 스프린트 내용·릴리스 구성은 열림. P1은 G1까지 4개월을 \textbf{잠정 기준선}으로
\item PMBOK 8판 Scope·Schedule·Finance·Risk(조달은 부록), PRINCE2 PBP(PPD\ra{}PBS\ra{}PD\ra{}PFD), ISO 21502 §7(통합·변경 통제·예비비 구분). 정밀도는 계약이 정한다
\item 범위 기술서의 핵심은 \textbf{인수 기준}과 \textbf{담당 붙은 제외}. WBS는 100\%·산출물 지향·8/80·통제 계정·같은 뿌리의 코드. WBS 사전이 RTM과 원가를 잇는다
\item 요구사항 3층·양방향 RTM·Must ≤ 60\%·Won't 명시·NFR은 FP에 없다. 범위 팽창 세 경로. 민간 SI의 변경 대가 = 변경협의체 + CCB
\end{enumerate}
\keymsg{2교시 — WBS 작업패키지가 활동이 되고, 활동의 네트워크가 임계경로와 여유를 만든다}
\note{제1장·제2장 핵심 정리. 온담 P1의 WBS 레벨 2 여섯(1.1\til{}1.5 + 1.6, WP 43), 상위 요구 8 \ra{} 214건의 절반이 비기능, 요구 제공자가 곧 인수자라는 점을 한 번 더. 2교시는 제3장 일정 — WBS의 작업패키지에서 활동을 뽑고, PDM으로 잇고, CPM을 손으로 돌린다. 예상 소요 3분.}
\end{frame}

% =====================================================================
\section{2교시 (90분) — 일정: CPM · float · 압축 · 몬테카를로 · 스프린트 × 게이트}
% =====================================================================

\begin{frame}{이 교시에서 배우는 것}
\begin{itemize}
\item 활동 · PDM 4종 의존 · lead/lag · 의존의 근거 · 유사·모수·3점 추정 [§3.1]
\item CPM 전진·후진 계산을 온담 요약 네트워크의 G2\ra{}G3\ra{}컷오버 구간(A11\til{}A17)으로 \textbf{손으로} 한다 — 총여유와 자유여유 [§3.2]
\item float는 누구의 것인가 — SCL Protocol 2판과 한국 SI 계약 문안 [§3.3]
\item 크리티컬 체인 · 압축의 한계 · DCMA 14-Point · 몬테카를로 P50/P80 · 제약 일자 [§3.4]
\item 하이브리드 달력 — 36 스프린트 위에 게이트·외부 고정일·성수기 창을 겹치면 보이는 것 · 임계경로 세 구간 [§3.5\til{}3.6]
\end{itemize}
\keymsg{P1의 임계경로 절반은 P1 PM이 통제하지 못하는 활동(P2 MC)이 결정한다 — 그것을 아는 것이 일정 기준선의 첫 번째 용도다}
\note{제3장 전체를 90분에. 시간 배분: §3.1 12분, §3.2 20분(계산표를 실제로 따라간다), §3.3 10분, §3.4 22분, §3.5\til{}3.6 18분, 정리 5분, 오류 3분. §3.2의 계산표는 6교시 실습 ②(워크북 §3 항목 5)의 예습이므로 수강생이 ES·EF를 한두 활동 직접 계산하게 한다. 예상 소요 2분.}
\end{frame}

\begin{frame}{활동 정의 — 하이브리드에서 어느 구간을 활동 수준 CPM으로 돌리는가}
\begin{itemize}
\item 활동은 작업패키지의 분해 — 동사로 적고, 담당 하나·기간 하나. 크기는 스프린트 안에 끝나는 것
\item \textbf{구현 스프린트 20회의 활동은 스프린트 백로그} — 착수 시점에 20회분을 미리 정의하면 rolling wave 위반
\item \textbf{통합·전환 구간(G2 이후)은 착수 시점에 활동 수준까지} — 되돌릴 수 없는 결정(G3) 앞 + 외부 고정일(컷오버 창)
\item 전환 프로그램 최종화 · 리허설 3회 · UAT · 롤백 검증 · 컷오버 절차가 그 활동들이다
\item 1.4절 대조표의 "SI는 활동 단위 CPM을 통합·전환 구간에서만"이 이 뜻
\end{itemize}
\keymsg{흔한 오류 — 활동을 명사로 적어 산출물과 구분이 안 되는 것 / 대기를 lag로 숨겨 소유자를 없애는 것}
\note{§3.1 "활동 정의". 작업패키지는 "무엇이 나오는가", 활동은 "그것을 만들기 위해 누가 무엇을 하는가". 워크북 §3 예시본은 착수\ra{}G4 전체를 요약 활동 20개(A01\til{}A20)로 놓고, 상세 일정(스프린트 단위 활동 약 180개)은 rolling wave로 남긴다 — 요약 네트워크와 상세 일정의 관계를 미리 말해 둔다. 예상 소요 4분.}
\end{frame}

\begin{frame}{PDM 4종 의존 · lead / lag · 의존의 근거 네 가지}
\begin{tbl}[\footnotesize]{L{0.06\textwidth}L{0.3\textwidth}L{0.52\textwidth}}
\textbf{유형} & \textbf{의미} & \textbf{온담 P1의 예} \\ \midrule
\textbf{FS} & 선행이 끝나야 후행 시작 (90\% 이상) & 리허설 1회차 종료 \ra{} 결함 수정 시작 \\
\textbf{SS} & 선행이 시작해야 후행 시작 & UAT 시작 \ra{} UAT 결함 트리아지 시작 \\
\textbf{FF} & 선행이 끝나야 후행 종료 & 전환 검증 6종 완료 \ra{} 리허설 3회차 보고서 완료 \\
\textbf{SF} & 선행이 시작해야 후행 종료 (거의 안 씀) & 새 POS 가동 시작 \ra{} OD-POS 병행 운영 종료 (컷오버의 정의) \\
\end{tbl}
\begin{itemize}
\item \textbf{Lead}(FS$-$5d)는 사실상 fast tracking, \textbf{lag}(FS+3d)는 물리적 대기에만. "회신 대기 10일"은 lag가 아니라 \textbf{활동}
\item 의존의 근거 — \textbf{필수}(정본 확정 후 전환 최종화) / \textbf{임의}(리허설 3회 순차) / \textbf{외부}(P2 MC) / \textbf{내부}(발주사 인력)
\end{itemize}
\keymsg{근거를 적어야 압축이 필요할 때 무엇을 풀 수 있는지(임의)와 없는지(필수·외부)가 보인다}
\note{§3.1 "순서". SS·FF는 병행을 표현하지만 남발하면 네트워크가 임계경로를 잃는다. SF는 "옛것의 종료가 새것의 시작에 묶이는" 전환 상황뿐. 그림 3-1(fig\_3\_1\_pdm)은 시간이 남으면. 기간 추정 세 방법은 다음 슬라이드. 예상 소요 5분.}
\end{frame}

\begin{frame}{기간 추정 — 유사 · 모수 · 3점}
\begin{itemize}
\item \textbf{유사} — 과거 실적(2014년 SAP 도입의 데이터 전환 실적 — 사례 문서에 없으므로 워크북에서 가정). RCF는 유사 추정의 통계적 정식화
\item \textbf{모수} — FP 12,400 ÷ 팀 생산성(FP/인월), 1,842 테이블 × 테이블당 매핑·검증 시간. 함정은 \textbf{생산성 계수의 출처}
\item 하이브리드의 정직한 방법 — 설계 스프린트 6회 동안 실측한 velocity가 G1 이후 구현 구간의 모수
\item \textbf{3점} — PERT (O + 4M + P)/6, σ = (P $-$ O)/6. 가치는 공식이 아니라 \textbf{비관값을 묻는 행위} — 담당자가 리스크를 말하기 시작한다(제5장 식별 입력)
\item 함정 둘 — 비관값이 충분히 비관적이지 않다(planning fallacy) / 활동별 3점을 그대로 몬테카를로에 넣으면 \textbf{상관 무시}로 P80 과소평가
\end{itemize}
\keymsg{흔한 오류 — 비관값을 최빈값 × 1.2로 기계적으로 만드는 것 / 업계 평균 생산성을 이 팀의 값으로 쓰는 것}
\note{§3.1 "기간 추정". 업계 평균은 이 팀의 값이 아니고, 이 팀의 지난 프로젝트 값은 이 프로젝트의 값이 아니다. 삼각분포 (O+M+P)/3도 쓴다. 3점 추정의 두 번째 함정(Hulett의 상관)은 §3.4 몬테카를로에서 다시. 6.6절 Jørgensen — "최대 X"라고 물으면 사람들은 P70\til{}P80을 말하면서 P100이라고 믿는다; 처방은 "이 값을 넘을 확률이 10\%인 값"처럼 확률을 명시해 묻는 것. 예상 소요 4분.}
\end{frame}

\begin{frame}{CPM — 전진·후진 계산, 총여유와 자유여유}
\begin{columns}[T]
\begin{column}{0.5\textwidth}
\subhead{두 번 지나간다}
\begin{itemize}
\item \textbf{전진} — ES, EF = ES + 기간. 후행의 ES = 모든 선행 EF의 최대
\item \textbf{후진} — LF, LS = LF $-$ 기간. 선행의 LF = 모든 후행 LS의 최소
\item \textbf{총여유 TF} = LS $-$ ES = LF $-$ EF. TF = 0의 연쇄가 임계경로
\item \textbf{자유여유 FF} = 직후 후행 ES $-$ 이 활동 EF (≤ TF)
\end{itemize}
\end{column}
\begin{column}{0.48\textwidth}
\subhead{두 여유의 차이가 실무에서 중요하다}
\begin{itemize}
\item 총여유는 \textbf{경로의 것}, 자유여유는 \textbf{활동의 것}
\item 어느 활동이 TF 10일을 다 쓰면 같은 경로의 뒤 활동들은 여유가 0 — 총여유는 \textbf{경로 위 활동들이 공유하는 자원}
\item 그래서 "float는 누구의 것인가"(§3.3)가 나온다
\end{itemize}
\end{column}
\end{columns}
\keymsg{CPM을 도구가 계산해 주는 것으로만 알고 손으로 한 번도 안 해 본 PM은 도구가 왜 그 경로를 임계라 하는지 설명하지 못한다}
\note{§3.2 개념. 후진 계산의 출발점은 고정된 날짜(워크북 §3.1: P1은 G4·G3·규제 ①에서 시작) — 그 결과 임계경로는 앞이 아니라 뒤에서 결정된다. 흔한 오류: 총여유·자유여유 미구분 / 임계경로가 하나라고 믿기(P1은 둘, 근사 임계 포함 셋) / 외부 마일스톤(P2 MC)을 네트워크 밖에 두기 / 여유를 "남는 시간"으로 알고 첫 활동이 다 쓰기. 다음 슬라이드에서 온담 G2\ra{}G3 구간을 실제로 계산한다. 예상 소요 5분.}
\end{frame}

\begin{frame}{온담 P1 G2 \ra{} G3 \ra{} 컷오버 — 손으로 계산해 보기}
\fig[0.57]{fig_3_6_critical_path_full}
\figcap{그림 3-6(구간 3). 위 줄 ES·EF, 아래 줄 LS·LF(작업일 번호, 0 = 2026-01-05), 붉은 상자 = TF 0. 임계경로는 UAT 사슬(A14\ra{}A16\ra{}MG3)과 P2 사슬(A15\ra{}MG3), 리허설 사슬(A11\ra{}A13)은 TF 5 [워크북 §3 값]}
\keymsg{A13 리허설 1\til{}3회차는 01-22에 끝나고 제약(FNLT)은 01-29다 — 여유 5일은 R-10(정본 후 마스터 변경)이 소비할 몫. 여유를 어디에 둘지는 여러분이 결정한다}
\note{§3.2 "온담 P1의 통합·전환 구간으로 계산해 보기" — 워크북 §3 항목 5 계산표의 구간 3 [추가 설정 — 영업일 기준, 공휴일 미반영]. 활동: A11 마스터 정본 판정 \ra{} G2 20(10-26\til{}11-20, G2 11-27 앞 5일) \ra{} A13 전환 리허설 1\til{}3회차 45(11-23\til{}01-22, FNLT 01-29); A10 구현 R3 종료 12-04 \ra{} A14 통합시험·성능·보안 30(12-07\til{}01-15) \ra{} A16 UAT·컷오버 리허설·롤백 24(01-18\til{}02-18), A12 구현 R4 잔여 20(12-07\til{}01-01, TF 10); A15 P2 설비 인터페이스 시험 5(SNET 02-12, 02-12\til{}02-18); MG3 G3 02-19(선행 A16·A15·A13·A07) \ra{} A17 컷오버 1(SNET 02-26) \ra{} A18 안정화 40. 수강생에게 전진 계산 A16 ES = max(A14 EF 270, A12 EF 260) = 270, EF 294 / MG3 ES = max(294, 294, 275, 180) = 294 = 02-19와 후진 계산 A15 LF = MG3 LS 294, LS = 289 = ES \ra{} TF 0 / A13 LF = 280(FNLT 01-29), LS = 235, TF = 235 $-$ 230 = 5를 직접 구하게 한다. A17 ES = max(MG3 294, SNET 299) = 299 — 02-22\til{}25의 4일은 여유가 아니라 제약 사이의 간격. 예상 소요 10분.}
\end{frame}

\begin{frame}{이 계산이 말하는 세 가지}
\begin{itemize}
\item \textbf{임계경로는 둘} — UAT 사슬(A14\ra{}A16)과 P2 사슬(A15). 리허설 사슬의 TF 5일은 세 조건 위에 있다: G2가 11-27에 판정(R-10 — 정본 후 마스터 변경이면 610 테이블 재매핑 6주가 들어와 즉시 임계) / 통합시험이 12-07 시작(구현 R3가 12-04 종료) / MC가 02-12(R-01 — 41주면 04-02, A15는 G3 뒤로)
\item \textbf{G3에서 컷오버까지 4영업일(02-22\til{}25)은 여유가 아니라 제약 사이의 간격} — 그 뒤 합류 버퍼 15영업일(\ra{}03-12)은 P1 것이 아니라 프로그램의 것. 창을 놓치면 다음은 3월 중순, 그 다음은 2027년 하반기
\item \textbf{P2 사슬이 임계경로에 있다} = P1 PM이 통제하지 못하는 활동이 P1의 임계경로를 결정한다
\item R-01 완화(물류 분리 컷오버)는 A15를 MG3의 선행에서 떼어 "2차 컷오버"를 별도 경로로 두는 것(1.3.4·1.6.9) — 결정은 프로그램 이사회가 G1까지
\end{itemize}
\keymsg{생각해 볼 질문 — P2 MC가 04-02로 통지되면 임계경로는 어떻게 바뀌고 2차 컷오버는 언제 가능한가(성수기 창을 고려하라)}
\note{§3.2 "이 계산이 말하는 것 세 가지". MC 04-02는 [추가 설정: 납품 02-05 + 설치·시운전 8주]. 워크북 §3.6 해설 — "P2 인터페이스 시험(A15)이 5일뿐인 이유: 충분하지 않다는 것을 기준선이 보여 주는 것이 이 활동의 존재 이유"; 프리모템 원문 3번("물류 인터페이스 시험을 못 한 채 컷오버를 갈지 말지를 G3 전날까지 아무도 결정하지 못했다")은 이 5일이 네트워크에 없을 때 일어난다. 6교시 실습 ②의 기본 과제가 정확히 이 시나리오(T0에서 41주 확정, MC 04-02)다. 예상 소요 5분.}
\end{frame}

\begin{frame}{float는 누구의 것인가 — SCL Delay and Disruption Protocol 2판}
\fig[0.56]{fig_3_3_float_ownership}
\figcap{그림 3-3. 수행사·발주사·프로그램의 세 주장은 각각 정당해 보인다. Core Principle 8 — 별도 약정이 없으면 float는 프로젝트의 공용 자원이며 먼저 쓰는 쪽이 쓴다. 여유(계산 결과)와 버퍼(소유자 있는 의도적 시간)는 다르다}
\keymsg{그것이 공정해 보이지 않는다면, 그것이 바로 계약 문안으로 달리 정할 이유다}
\note{§3.3. 데이터 경로(A02\ra{}A09\ra{}A11\ra{}A13) TF 5일 — 이재현 PM("여유가 있으니 인력을 다른 사업에"), 발주사 PM("발주사 미결정 대기를 흡수하기 위한 것"), 프로그램 PMO장("P2가 늦으면 P2에 준다"). 결정적인 이유는 지연 클레임 — 발주사 사유로 3일 늦었는데 여유 5일이 있었다면 EOT를 청구할 수 있는가. SCL 2판(2017) CP8: first come, first served — 발주사 사유로 여유가 소진되어도 종료일이 안 늦으면 시공자는 EOT를 못 받고, 그 뒤 시공자 사유 지연은 여유가 없으므로 시공자 책임. 2판의 다른 둘: 지연 분석 6종(TIA 단일 선호 폐기, 기록과 사실관계에 따라 선택), 동시 지연을 좁게 정의(진정한 동시성만 — "P2도 늦고 P1도 늦었다"의 판단 근거). AACE 29R-03은 미국 쪽 참조(9개 MIP). 예상 소요 6분.}
\end{frame}

\begin{frame}{한국 SI 계약에는 float 조항이 없다 — 발주사 PM이 할 수 있는 셋}
\begin{itemize}
\item SI 계약의 이분법 — "발주사 사유 지연은 공기 연장, 그 밖은 지체상금". 여유의 존재를 무시하므로 분쟁에서 두 주장이 충돌한다
\item ① \textbf{여유의 위치를 기준선에서 명시} — "데이터 경로(A02\ra{}A13) 여유 5영업일은 R-10과 발주사 데이터 정제 지연을 흡수하기 위한 것"
\item ② \textbf{버퍼와 여유를 구분} — 합류 버퍼 15영업일(프로그램 PMO장), 규제 버퍼 22영업일(P4)은 소유자가 명시된 버퍼
\item ③ \textbf{계약 변경 시 float 조항} — 변경협의체 의결서에 "이 변경은 여유 N일을 소진하며, 소진 후 지연은 어느 쪽 사유인가"
\item 워크북 §3 항목 8의 표준 문안 — "여유는 어느 쪽의 소유도 아니며 먼저 쓰는 쪽이 쓴다 … 완료일이 밀리지 않으면 공기 연장은 없다" + 예외 2개
\end{itemize}
\keymsg{흔한 오류 — 여유와 버퍼를 같은 말로 / 동시 지연을 "둘 다 늦었으니 반반"으로 / 지연 분석 방법을 하나만 알고 표준이라 믿기}
\note{§3.3 "한국 SI 계약에서의 함의"와 워크북 §3 항목 8. 계약 본문을 바꿀 수 없어도 변경협의체 의결서에 그 문장이 있으면 다음 분쟁의 근거가 된다. 생각해 볼 질문 1 — "데이터 경로의 총여유 5영업일을 누구에게 줄 것인가. 기준선 문서에 어떻게 적을 것인가"를 조별로 한 문장씩. 예상 소요 4분.}
\end{frame}

\begin{frame}{크리티컬 체인 — 방법론이 아니라 두 개념을 가져온다}
\begin{itemize}
\item Goldratt(1997)의 관찰 — 담당자는 안전여유를 넣고(50\% 추정의 두 배 가까이), 그 여유는 \textbf{학생 증후군·파킨슨 법칙·병합 편향}으로 낭비된다
\item 처방 — 50\% 추정으로 계획하고 걷은 여유를 \textbf{프로젝트 버퍼}(임계 체인 끝)와 \textbf{피딩 버퍼}(합류 지점)로 모아 소진율로 관리
\item 가져올 것 ① \textbf{병합 편향} — G3에 여섯 경로가 합류. 각각 P50이면 모두 제때 올 확률은 훨씬 낮다 \ra{} 합류 버퍼 15영업일의 논리
\item 가져올 것 ② \textbf{버퍼 소진율 관리} — "며칠 늦었다"가 아니라 "버퍼의 몇 \%를 썼는가, 진척 대비 빠른가"(Day3 보고 지표)
\item SI·EPC에서 방법론으로 실패하는 이유 — \textbf{하도급·용역 계약이 활동 단위 기한으로 묶여 여유를 회수할 법적 근거가 없다}
\end{itemize}
\keymsg{이재현 PM에게 "리허설 기간의 안전여유를 걷어 프로젝트 버퍼로 넘기라"고 할 수 있는 계약 조항은 없다}
\note{§3.4 "크리티컬 체인". 임계 체인 = 자원 제약을 반영한 임계경로(OD-POS를 아는 한 사람의 시간이 곧 임계 체인 — §4.4 자원 평준화에서 다시). Raz, Barnes \& Dvir(2003)의 학술적 비판은 마무리 블록 6.3절. 오독 둘 — CCPM을 "여유를 없애는 방법"으로 읽기(위치를 옮기고 소유자를 정하는 것) / 버퍼 소진율 관리를 도입하면서 버퍼 소유자를 안 정하기(float 소유권 문제가 이름만 바뀌어 돌아온다). 예상 소요 5분.}
\end{frame}

\begin{frame}{일정 압축 — crashing과 fast tracking의 한계, 그리고 정직한 답}
\begin{columns}[T]
\begin{column}{0.5\textwidth}
\subhead{Crashing — 자원을 더 넣는다}
\begin{itemize}
\item 비용은 \textbf{비선형} — 첫 하루는 싸고 마지막 하루는 비싸다
\item \textbf{Brooks의 법칙} — 늦은 SW 프로젝트에 인력을 더하면 더 늦는다. 경로 n(n$-$1)/2, 82명이면 3,321개
\item 고정가 계약에서 crashing은 수행사 원가 = \textbf{변경 대가}
\end{itemize}
\end{column}
\begin{column}{0.48\textwidth}
\subhead{Fast tracking — 순차를 병행한다}
\begin{itemize}
\item 비용은 안 들지만 \textbf{재작업 리스크} — Williams(2002)의 악순환: 압축 \ra{} 병행 \ra{} 재작업 \ra{} 추가 압축
\item \textbf{임의 의존에서만} 가능, 필수 의존에서는 불가 — 근거를 적어 둔 이유
\end{itemize}
\end{column}
\end{columns}
\begin{itemize}
\item P1 — 리허설 3\ra{}2회는 가능하나 R-10 완화 약화 / A16 UAT·리허설·롤백 24\ra{}19일은 0.4건/FP 위협 / P2 MC와 규제 기한은 \textbf{압축 불가}
\end{itemize}
\keymsg{압축의 정직한 답은 대개 "범위"다 — 헌장 §13이 이미 그렇게 정했다. "네"가 아니라 선택지·권고·결정 기한}
\note{§3.4 "일정 압축"과 §3.7 "압축하면 된다". 대표이사의 "가능하면 더 앞당기라"에 "네"라고 답하고 crashing 비용과 재작업 리스크를 말하지 않는 것이 오류. 답은 선택지 — (a) 범위 이연(헌장 §13), (b) crashing 비용 N억(변경 대가), (c) 병행 시 재작업 리스크 — 와 권고와 결정 기한(Day1 3.4절 "결정을 결정권자에게 가져가기"). 생각해 볼 질문 2(정미란 "컷오버를 2월 첫 주로 당길 수 있는가")를 조별 과제로. 예상 소요 5분.}
\end{frame}

\begin{frame}{DCMA 14-Point — 일정의 위생, 정확히는 float 은닉 탐지 도구}
\begin{tbl}[\scriptsize]{cL{0.34\textwidth}L{0.12\textwidth}L{0.37\textwidth}}
\textbf{\#} & \textbf{지표} & \textbf{임계치} & \textbf{무엇을 잡는가} \\ \midrule
1 · 2 · 3 & Logic(선후행 없는 활동) · Leads · Lags & ≤5\% · 0\% · ≤5\% & 논리 없는 간트 · 숨겨진 fast tracking · 활동을 대기로 숨김 \\
4 · 5 & FS 관계 비율 · Hard constraints & ≥90\% · ≤5\% & SS·FF 남발 · 논리 대신 날짜 고정 \\
6 · 7 · 8 & High float(TF>44일) · Negative float · High duration(>44일) & ≤5\% · 0\% · ≤5\% & 논리가 끊겨 부풀려진 여유 · 제약과 논리의 충돌 · 진척이 안 보이는 활동 \\
9 · 10 · 11 & Invalid dates · Resources 미배정 · Missed tasks & 0\% · 0\% · ≤5\% & 갱신 안 된 일정 · 원가 미연결 · 기준선 대비 성과 \\
12 · 13 · 14 & Critical path test · CPLI · BEI & 통과 · ≥1.00 · ≥1.00 & 임계경로에 600일 주입 시 종료일이 정확히 600일 밀려야 · 임계경로 현실성 · 실행 효율 \\
\end{tbl}
\keymsg{SI에서 가장 자주 걸리는 것은 5번(임원이 말한 날짜를 제약으로)과 6번(논리를 안 그려 여유 200일인 활동이 절반)}
\note{§3.4 "DCMA 14-Point". 공식 문서로 현행 유지되는지는 확인되지 않았다(리딩리스트 [검증 필요])지만 14개 임계치는 업계 관행으로 살아 있고 대형 EPC·방산 발주처가 일정 제출 시 요구한다. 기준선에는 1\til{}9번이 적용된다(워크북 §3 항목 9). 워크북 예시본은 요약 네트워크라 4개 미달(래그 15\%, FS 81\%, 강제 제약 43\%, 과다 기간 40\%) — 요약 수준의 특성이며 상세 일정(약 180활동)에서 해소된다고 사유를 적었다. 오류: 도구의 점검표로만 돌리고 왜 걸렸는지 안 보기 / 기준을 맞추려고 제약을 지워 여유를 만들기. 예상 소요 5분.}
\end{frame}

\begin{frame}{몬테카를로와 P50 / P80 — 상관을 무시하면 P80은 거짓말이 된다}
\fig[0.55]{fig_3_4_monte_carlo}
\figcap{그림 3-4. 활동별 3점 추정을 독립으로 합산하면(회색) 분산이 줄어 P80이 P50에 붙는다. 같은 원인(R-01·R-10·R-05)이 여러 활동을 동시에 늦추는 risk driver 방식(남색)은 꼬리가 길다. 개념 시뮬레이션 [추가 설정]}
\keymsg{목표는 P50이고 버퍼가 P50\til{}P80을 덮는다 — P80은 버퍼의 크기이지 팀에 주는 목표가 아니다}
\note{§3.4 "몬테카를로와 P50/P80". 세 함정: 상관 무시(큰 수의 법칙으로 분산이 줄어 P80이 P50에 붙는다 — 실제로는 발주사 인력 부족·정본 지연이 여러 활동을 동시에 늦춘다; Hulett의 처방은 risk driver 방식 — R-06은 인터페이스 문서화·데이터 매핑·리허설 결함 수정에 동시에 걸린다), 병합 편향(입력이 독립이면 과소평가), P80을 기준선으로 약속하기. 예시 [추가 설정]: 컷오버 P50 = 02-26, P80 = 03-19라면 합류 버퍼 15영업일(03-12)은 약 P70 — P80까지 덮으려면 5영업일 더 필요하거나 R-01·R-10의 확률을 낮춰야 한다. 이 숫자를 스폰서에게 가져가는 것이 일정 리스크 분석의 목적. 예상 소요 6분.}
\end{frame}

\begin{frame}{제약 일자 — 외부 고정일에만, 게이트는 논리의 결과로}
\begin{itemize}
\item 종류 — SNET(이 날 이후 시작) / FNLT(이 날 이전 종료) / MSO·MFO(반드시 이 날). 하드 제약이 5\%를 넘으면 네트워크는 논리를 잃는다
\item \textbf{음의 여유} = "논리대로 하면 제약을 못 지킨다"는 신호 — 제약을 풀어 없애지 말고 리스크로 등록해 압축 또는 범위 결정으로
\item P1의 외부 제약 — 컷오버 창(A17 SNET 02-26) / 규제 ①(A07 FNLT 09-11) / 리허설 3회차(A13 FNLT 01-29) / T0(1.3.1 FNLT 04-24) / MC(A15 SNET 02-12)
\item \textbf{게이트는 논리의 결과여야 한다} — 예시본은 프로그램이 고정한 게이트에 FNLT를 걸되 출처를 적었다(하드 제약 9/21). 출처 없이 박으면 앞 활동이 음의 여유를 갖고 "문서 맞추기"가 된다
\item 워크북 예시본 — "논리가 만든 0"(후행을 당기면 풀린다)과 "제약이 만든 0"(A07 규제, A15 MC — 풀리지 않는다)을 구분해 표시
\end{itemize}
\keymsg{규제 기한을 제약 없이 두면 네트워크는 규제 기한을 넘기는 일정을 태연히 계산한다 — 반대 오류도 오류다}
\note{§3.4 "제약 일자"와 워크북 §3 항목 3·§3.6 해설. 규제 ① FNLT 09-11(A07 규제 검증 08-24\til{}09-11). 워크북 예시본은 제약 9개(FNLT 7, SNET 2)이고 외부 가정 날짜 3개(3PL 06-30, 부속서 08-14, 앱 벤더 12월)는 제약이 아니라 리스크 트리거로 등록했다 — "제약인 척하는 희망 날짜"를 걸러 내는 방법은 출처를 적는 것. 워크북 §3.6 첫 해설: A07의 LF를 G4에서 역산하면 여유 200일이 나오는데 그 여유는 거짓 — 중간 제약(FNLT)이 있는 활동은 LF를 제약 날짜로 눌러야 한다. 예상 소요 5분.}
\end{frame}

\begin{frame}{하이브리드 달력 — 36 스프린트 위에 게이트 · 외부 고정일 · 성수기 창 · 릴리스}
\fig[0.7]{fig_3_5_hybrid_calendar}
\figcap{그림 3-5. 스프린트 36회 · 동결 12회 / 릴리스 R0\til{}R5 / 게이트 / 외부 고정일 / 성수기 창 [추가 설정]}
\keymsg{게이트는 스프린트 경계에 놓여야 한다 — 그렇지 않은 곳(G2)이 바로 보인다}
\note{§3.5 "한 달력에 놓기", 워크북 §3 항목 7. 2026-01\til{}03 S1\til{}S6 설계 R0(설 전 2주 02-03\til{}16 매장 변경 없음, 동결 FZ-01\til{}03, 설계 종료 03-27) / 04 S7\til{}S8 R1(T0 04-24, 단말 조사 완료 04-30, G1 04-30) / 05\til{}06 S9\til{}S12(포털 프로토타입 05, R1 종료 06-19, 3PL 타결 목표 06-30) / 07\til{}08 S13\til{}S16 R2(부속서 08-14, 3PL API 종료 08-14, 보안진단 1회차, 여름 성수기 배포 금지) / 09 S17\til{}S19(규제 검증 08-24\til{}09-11, 규제 ① 09-11 = R2 종료, 추석 전 잠금 09-11\til{}24) / 10\til{}11 S20\til{}S24 R3(매핑 종료 10-23, FZ-10 10-28·FZ-11 11-25, G2 11-27, 포털 선행 릴리스, R3 종료 12-04) / 12 S25\til{}S26 R4(통합시험 착수 12-07, 리허설 1회차, FZ-12 12-30, 연말 성수기) / 2027-01 S27\til{}S28(통합시험 종료 01-15, UAT 착수 01-18, 리허설 3회차·전자금융업 신청 01-29, 설 전 잠금 01-24\til{}02-06) / 02 S29\til{}S30(MC 02-12, G3 02-19, 컷오버 02-26\til{}28 = R4 종료) / 03 S31\til{}S32 R5(버퍼 종료 03-12, 규제 ② 03-31) / 04\til{}05 S33\til{}S36(안정화 종료 04-23, 운영 이관, G4 05-28 — 헌장 05-29 토). 명절·성수기 창 날짜는 [추가 설정]. 이 그림은 6교시 실습 ②(워크북 §3 항목 7 스프린트 × 게이트 캘린더)의 참조본. 예상 소요 4분.}
\end{frame}

\begin{frame}{달력을 놓아야 보이는 네 가지}
\begin{itemize}
\item \textbf{G2는 스프린트 경계에 있지 않다} — 11-27은 S24(11-23\til{}12-06)의 첫 주 금요일. 옮길 수 없으니 S24를 "G2 스프린트"로 정의(첫 주 게이트 패키지, 둘째 주 판정 후 작업)
\item \textbf{가맹 포털 선행 릴리스는 12월을 피해야 한다} — O-02의 "2026-Q4"는 성수기. 11월 마지막 주(G2 준비와 충돌)거나 1월 첫 주(UAT와 충돌). 워크북 예시본은 S24(R3 종료 12-04, G2와 같은 스프린트)에 두었다 — 충돌을 알고 받아들인 결정
\item \textbf{2027년 설(2월 7일 전후)이 UAT 결함 수정에 걸린다} — 현업 UAT 참여자(매장·물류)가 성수기 업무로 빠진다. A16(UAT·리허설·롤백, 01-18\til{}02-18)이 임계경로다
\item \textbf{컷오버 주말 34시간의 앞뒤가 일정이다} — 금 20:00 전 패키지 동결·매장 통지·서비스데스크 증원(박준영)·go/no-go 판정자(정미란), 일 06:00 후 월요일 개점 대응(슈퍼유저 124명, P3)
\item 컷오버 계획은 Day4의 주제이지만 그 활동들은 Day2의 일정 기준선에 있어야 한다
\end{itemize}
\keymsg{"36회 × 2주 = 72주"라는 산수는 성수기·명절·G2가 스프린트 중간에 있다는 사실·설 연휴가 UAT에 걸린다는 사실을 모른다}
\note{§3.5 "이 달력에서 보이는 것 네 가지"와 §3.7 "달력 없이 스프린트 수만 세는 것". 계획에 적지 않으면 S24는 목표 없는 스프린트가 된다. 어느 쪽이든 자원 충돌이 있고, 그것이 릴리스 계획을 달력 위에 놓아야 하는 이유. 워크북 §3.6 해설 — 스프린트 캘린더에 동결일 12회와 접수 마감일(2주 전, R-07 대응)을 날짜로 넣어야 팀이 지키고 변경협의체가 "마감 뒤 접수분은 다음 달"이라 판정할 근거가 생긴다; 2027-01 이후는 동결 없이 결함만(UAT 코드 프리즈). 예상 소요 5분.}
\end{frame}

\begin{frame}{임계경로는 세 구간을 지난다 — T0는 납품 기한, MC는 입력 기한}
\fig[0.57]{fig_3_6_critical_path_full}
\figcap{그림 3-6. 착수 \ra{} G4 요약 네트워크 20활동(워크북 §3 값). 붉은 = 임계경로(365영업일), A15·A07 = 제약이 만든 0, 초록 = 근접 임계}
\keymsg{T0와 MC 사이 41주는 P1이 통제하지 못하는 시간 — 할 수 있는 것은 T0 전 규격 인도와 분리 컷오버 설계뿐}
\note{§3.6. 구간 1(2026-01\til{}04): 요구 확정(SyRS 396·SRS 570)과 인터페이스 13본 문서화 \ra{} G1 — 리스크 R-05·R-06, 둘 다 발주사 내부; G1 지연 = G2 지연이거나 구현 스프린트 수 감소(범위 축소). 구간 2(04\til{}11): 구현의 임계는 활동이 아니라 \textbf{마스터 데이터 6종의 확정}(상품·고객·매장·거래처·가격·재고) — 상품 코드 요구 접수 마감 2026-10(1.5.8), 3PL 협상 결과, 별도 법인 데이터 주체(R-09)에 걸림; PM이 볼 것은 velocity가 아니라 마스터 확정 진도. 구간 3(11\til{}2027-02): §3.2의 계산. T0는 P1의 납품 기한(1.3.1 재고 원장 규격을 T0 04-24 FNLT로 — 설계 스프린트 안에 재고 원장 데이터 구조 확정), MC는 입력 기한(MC 전엔 물류 인터페이스를 실제 설비와 시험 못 함). 워크북 예시본은 공휴일 미반영 — 반영 시 임계경로 약 20영업일 증가, G4 앞 여유 음수(§3.6 해설: 안정화 40\ra{}30일 / 이관과 SS 겹침 / 허용범위를 06-18로 이해 — 결정은 스폰서). 예상 소요 6분.}
\end{frame}

\begin{frame}{2교시 핵심 정리}
\begin{enumerate}
\item 하이브리드에서는 통합·전환 구간만 활동 수준 CPM, 구현 구간은 스프린트 백로그. PDM 4종 · lead/lag · \textbf{의존의 근거}를 적어야 압축 시 무엇을 풀 수 있는지 안다
\item CPM은 손으로 한 번. 총여유는 경로가 공유하는 자원, 자유여유는 활동의 것. P1 G2\ra{}G3의 임계경로는 UAT 사슬(A14\ra{}A16)과 P2 사슬(A15), 리허설 사슬 TF 5, G3 뒤 창 대기 4일(여유 아닌 간격), 프로그램 버퍼 15
\item Float는 별도 약정이 없으면 공용 자원·선착순(SCL CP8) — 기준선에 여유의 위치와 목적을 명시하고, 여유와 버퍼를 구분하고, 협의체 의결서에 소진을 적는다
\item CCPM은 계약 구조 때문에 SI에서 방법론으로 실패하지만 병합 편향·버퍼 소진율은 가져온다. 압축의 답은 대개 범위. DCMA는 float 은닉 탐지. 몬테카를로는 상관 반영, P80은 목표가 아니다. 제약은 외부 고정일에만
\item 달력을 겹치면 보이지 않던 것이 보인다. 임계경로는 요구 확정(G1) \ra{} 마스터 6종(G2) \ra{} UAT·P2 합류(G3)의 세 구간, T0는 납품 기한·MC는 입력 기한
\end{enumerate}
\keymsg{3교시 — 활동의 자원 × 기간이 원가의 한 칸이 된다. 그 칸들을 시간축에 놓으면 PV 곡선이다}
\note{제3장 핵심 정리 1\til{}6. 오늘 일정에서 가장 중요한 한 문장은 "P2 사슬이 임계경로에 있다"이다 — P1 PM이 통제하지 못하는 활동이 P1의 임계경로를 결정하며, 그것을 기준선에 보이게 하는 것이 분리 컷오버 결정을 이사회에 올리는 근거다. 생각해 볼 질문 4·5(자기 회사 일정을 DCMA로 점검하면 몇 개 걸리겠는가 / 수행사가 활동마다 안전여유를 넣었다는 것을 어떻게 알 수 있는가)를 던지고 3교시로. 예상 소요 3분.}
\end{frame}

% =====================================================================
\section{3교시 (90분) — 원가와 조달: AACE · 예산의 3층 · PV 곡선 · 계약 유형}
% =====================================================================

\begin{frame}{이 교시에서 배우는 것}
\begin{itemize}
\item 추정 방법과 정확도 — AACE Class는 정의도로 결정되고 정확도 범위는 자동으로 따라오지 않는다 [§4.1]
\item 예산의 3층 — 컨틴전시와 관리예비비는 \textbf{다른 사람의 돈}이다 · 컨틴전시의 산정 근거 셋 [§4.2]
\item 시간대별 배분과 S-curve — PV·약정·현금 세 곡선, 자금 조달 한도, 배분 방법 [§4.3]
\item 조달 — 계약 유형과 리스크 배분, FP 고정가의 두 부분, 견적 검산 3단계, 하도급, 자원 계획 [§4.4]
\item 온담 P1 — 168.0의 근거 등급, 157.8\ra{}145.8의 산수, 월별 PV, 컨틴전시 배정 원칙, 관리예비비 절차 [§4.5]
\end{itemize}
\keymsg{"예산 168억"이라는 한 숫자로 말하는 순간 156과 145.8이 어디 갔는지, 10.2를 누가 쓰는지 아무도 모른다}
\note{제4장 전체를 90분에. 시간 배분: §4.1 12분, §4.2 20분, §4.3 12분, §4.4 18분, §4.5 20분, 정리 5분, 오류 3분. §4.5의 12.0 산수와 월별 PV는 6교시 실습 ②(워크북 §4)의 직접 입력이고, 컨틴전시 배정 원칙은 5교시 리스크 정량화 뒤에 완성된다. 점심은 이 교시 뒤. 예상 소요 2분.}
\end{frame}

\begin{frame}{AACE 추정 분류 — Class는 정의도로 결정되고, 정확도는 리스크 분석으로 산정한다}
\fig[0.55]{fig_4_1_aace_cone}
\figcap{그림 4-1. 남색 띠 = 18R-97(프로세스 산업), 붉은 파선 = 56R-08(건축·일반 건설). 같은 정의도에서 프로세스 플랜트가 건축보다 불확실성이 크다는 것 자체가 요점. 범위 숫자는 인쇄 전 원문 대조}
\keymsg{"Class 3이면 ±20\%"는 규칙이 아니다 — "Class 3 정의도에서 리스크 분석을 제대로 하면 대개 이 밴드가 나온다"는 경험적 관찰이다}
\note{§4.1 개념. Class 5(정의도 0\til{}2\%) \ra{} 4(1\til{}15) \ra{} 3(10\til{}40, 예산 승인·통제 기준선) \ra{} 2(30\til{}75, 통제·입찰) \ra{} 1(65\til{}100). 두 문서 모두 명시 — 정확도 범위는 자동 적용값이 아니며 리스크 분석으로 산정해야 하고, 컨틴전시가 적정하다면 약 80\%의 프로젝트가 이 밴드에 든다. 건축용 문서 번호는 56R-08이며 56R-11이 아니다 — 오인이 매우 흔하다. SI에는 대응 표준이 없으므로 이 언어를 빌려 쓴다. 예상 소요 5분.}
\end{frame}

\begin{frame}{이 체계가 SI에 주는 것 둘 · 상향식 추정의 함정 둘}
\begin{itemize}
\item ① \textbf{"우리는 어느 Class로 견적했는가"} — 입찰 시점 FP 12,400은 RFP 기준 개략 계수 = Class 4\til{}3, G1 확인서는 설계 6회 후 = Class 2에 가깝다. 그 차이가 ±3\% 재확인 조항의 이유
\item 168.0억이라는 승인 예산이 Class 3 수준의 값이었음을 \textbf{스폰서가 알아야 한다}
\item ② \textbf{정의도를 올리기 전에 총액을 확정하지 않는다} — SI에서는 "설계 완료 전 FP 확정 불가". 그래서 G1이 착수 4개월 뒤
\item 상향식(작업패키지별 합)의 함정 ① \textbf{누락} — WBS 100\% 규칙이 깨지면 추정도 깨진다
\item 함정 ② \textbf{낙관의 합산} — P50 100개를 더하면 P50이 아니다(비대칭 분포 + 상관). 반드시 하향식(유사·모수·RCF)과 대조 — 2014 SAP 74억 vs P1 168억을 규모로 정규화
\end{itemize}
\keymsg{흔한 오류 — 56R-11이라고 쓰기 / "SI에는 Class가 없으니 우리는 그냥 견적"(그러면 정확도를 말할 언어가 없다) / 입찰 숫자와 설계 후 숫자가 같아야 한다고 믿기}
\note{§4.1 "이 체계가 SI에 주는 것"과 "상향식 추정과 그 함정". 낙관의 합산 — 비대칭 분포(초과는 크고 절감은 작다)에서는 합계의 평균이 P50 합보다 크고, 상관이 있으면 분산이 커진다. 하향식 대조의 최소 형태가 2014년 SAP 도입 74억과 P1 168억을 FP·테이블 수·매장 수로 정규화해 비교하는 것이다. 워크북 §4 항목 1 — 고정가 68.2는 등급 대상이 아니고, 추정부 77.6은 설계 스프린트 6회 종료 기준 "Class 3 상당" [추가 설정], 정확도 범위는 과거 실적 부재로 미산정 — 컨틴전시 배정 논리로 대체. 예상 소요 5분.}
\end{frame}

\begin{frame}{예산의 층 — 계획 원가 145.8 / 컨틴전시 10.2 / BAC 156.0 / 관리예비비 12.0 / 승인 168.0}
\fig[0.55]{fig_4_2_budget_layers}
\figcap{그림 4-2. 온담의 "집행 한도" 156.0이 PMBOK의 원가 기준선(BAC). 컨틴전시 10.2는 기준선 안(PM의 돈, 리스크 ID를 달고 집행), 관리예비비 12.0은 기준선 밖(스폰서의 돈, 예외 보고·결재·기준선 갱신). PV 곡선은 145.8로 그린다}
\keymsg{10.2가 승인 예산의 6.07\%라는 것은 결과이지 근거가 아니다}
\note{§4.2 개념. 원가 기준선 = 시간대별로 배분된 승인 예산에서 관리예비비를 제외한 것 — 정의의 세 부분(시간대별 / 승인 / 관리예비비 제외)이 각각 중요. 층: 계획 원가(작업패키지 상향식 합, PM, 알려진 작업, 145.8) + 컨틴전시(식별된 리스크 known unknowns, PM이 허용범위 안에서 집행·기록·보고, 10.2) = 원가 기준선/BAC(156.0 = 집행 한도); + 관리예비비(unknown unknowns·범위 밖 사건, 스폰서 결재, 12.0 = 프로그램 24억 중 P1 유보분) = 승인 예산 168.0. Day3 EVM: PV는 145.8로, BAC는 156.0, 컨틴전시 집행은 PV에 없던 원가가 AC에 나타나는 형태. 예상 소요 6분.}
\end{frame}

\begin{frame}{컨틴전시와 관리예비비 — 크기가 아니라 누가 어떤 근거로 쓰는가}
\begin{columns}[T]
\begin{column}{0.5\textwidth}
\subhead{컨틴전시 — 등록부에 있는 리스크의 대응 비용}
\begin{itemize}
\item 근거가 있다 — R-08(앱 벤더 계약 외 주장, P1측 재작업 배정 0.60 — 벤더 견적 4억은 앱 유지보수 계약) 현실화 시 집행하고 등록부에 "집행 N, 잔여 M" 기록
\item 크기는 등록부의 정량화(EMV·몬테카를로)에서, 집행은 \textbf{리스크 ID를 달고}
\item RACI 19행: A = P1 PM, CCB 월 보고, 재경 C — 2020년 감사 지적 "집행 승인 근거 부재"의 답
\end{itemize}
\end{column}
\begin{column}{0.48\textwidth}
\subhead{관리예비비 — 등록부에 없는 것}
\begin{itemize}
\item 식별 못 한 리스크, 범위 밖 사건(규제 변경, 프로그램 재정의, 별도 법인 계약 변경)
\item 근거가 없으므로 PM이 못 쓴다 — \textbf{스폰서가 기준선 밖의 사건임을 판단해 결재}
\item 집행하면 대개 \textbf{기준선이 바뀐다}. RACI 20행: A = 정미란
\end{itemize}
\end{column}
\end{columns}
\keymsg{섞으면 — 컨틴전시에 스폰서 결재를 요구하면 "권한 없는 PM", 관리예비비를 PM 재량으로 쓰면 이사회 조건 2항 위반이고 미지의 사건에 돈이 없다}
\note{§4.2 "컨틴전시와 관리예비비의 구분". 컨틴전시 산정의 세 근거(다음 슬라이드). 흔한 오류: 예비비를 일반 원가에 섞기(작업패키지마다 10\%씩 얹으면 컨틴전시는 보이지 않고 파킨슨 법칙으로 사라진다) / 백분율로 정하고 근거를 안 묻기 / 리스크 ID 없이 집행 / 컨틴전시가 남으면 "절감"으로 보고 — 남은 컨틴전시는 리스크가 실현되지 않은 것이지 PM이 잘한 것이 아니며, 성과로 보상하면 다음 프로젝트의 컨틴전시는 부풀려진다(Day1 1.2절 Eveleens \& Verhoef). 예상 소요 5분.}
\end{frame}

\begin{frame}{컨틴전시의 크기 — 백분율에서 EMV · 몬테카를로 · RCF로}
\begin{itemize}
\item \textbf{EMV의 합} — 리스크마다 확률 × 영향을 더한다. 투명하지만 \textbf{평균} — "평균적으로 충분"이지 "80\% 확률로 충분"이 아니다
\item \textbf{몬테카를로의 P-값} — 등록부와 추정 불확실성을 시뮬레이션해 P70\til{}P80에서 계획 원가와의 차이. AACE가 "리스크 분석으로 산정하라"고 한 것의 실체
\item \textbf{RCF의 uplift} — reference class의 초과율 분포에서 읽는다. 외부 시각, 내부 등록부의 낙관을 교정. 온담에는 reference class가 없다 — 그것이 첫 리스크
\item 세 방법이 다른 숫자를 주면 \textbf{그 차이가 정보다} — 잔여 EMV 합 6.8억, P80 10억이면 3억의 간격이 꼬리
\item 컨틴전시 10.2는 잔여 노출의 약 P80 [추가 설정, §5.7]이고, 그 위 꼬리는 관리예비비 12.0과 구조(분리·폴백)가 덮는다는 것을 스폰서가 알아야 한다
\end{itemize}
\keymsg{어느 표준도 컨틴전시의 크기를 정해 주지 않는다 — 구분하라고 하지 얼마인지는 말하지 않는다}
\note{§4.2 "컨틴전시의 산정". 이 슬라이드의 세 방법이 5교시 §5.3(EMV·민감도·몬테카를로·멱함수 꼬리)과 마무리 6.1(RCF 절차)에서 각각 깊어진다. 컨틴전시 = EMV 합은 5.7절과 워크북 §4 항목 6에서 "컨틴전시 대상 잔여 EMV 4.58 \ra{} 배정 4.60 + 미배정 풀 5.60(스프린트 4회분) = 10.2"로 구체화된다 — R-01·R-04의 잔여 2.20은 관리예비비 대상. 예상 소요 4분.}
\end{frame}

\begin{frame}{시간대별 배분과 S-curve — 같은 145.8이 세 개의 다른 곡선이 된다}
\begin{tbl}[\scriptsize]{L{0.15\textwidth}L{0.14\textwidth}L{0.55\textwidth}}
\textbf{곡선} & \textbf{기준} & \textbf{누가 보는가} \\ \midrule
\textbf{PV (계획가치)} & \textbf{작업} 시점 & PM — Day3 EVM의 입력 \\
\textbf{약정} & 계약·발주 시점 & 구매(강현도), 재경 — 계약·발주 한도 168.0 \\
\textbf{현금} & 지급 시점 & 재경 — 자금 계획, 연차 예산 176/198/54 \\
\end{tbl}
\begin{itemize}
\item 배분 원칙은 \textbf{작업의 시점} — SI 68.2가 선급·중도·잔금이어도 PV는 스프린트 진척(설계 0.7 / 구현 1.1 / 통합 1.2 가중 [추가 설정])을 따른다
\item CFO의 "1차연도 148억"은 현금의 문제, PM의 PV는 작업의 문제 — 지급 조건 협상으로 현금은 바뀌고 PV는 안 바뀐다
\item \textbf{자금 조달 한도}를 문서에 적고, 넘는 구간은 작업을 미루거나 한도를 올린다. 배분 방법(마일스톤·고정 공식·완료 비율·배분 노력·노력 수준)은 계획 시점에
\end{itemize}
\keymsg{PV를 지급 일정으로 그리면 Day3의 EVM은 "돈을 냈는가"를 재지 "일이 되었는가"를 재지 않는다}
\note{§4.3. 이것이 PMBOK 8판이 Cost를 Finance로 바꾼 이유의 실체. 그림 4-3(fig\_4\_3\_three\_curves — PV는 실값, 약정·현금·한도는 개념)은 시간이 남으면 띄운다. FP 기반 작업패키지는 FP 완료 비율, 리허설은 0/100(회차별), 라이선스는 납품 마일스톤, PMO는 노력 수준이 자연스럽다 — 배분 방법을 정하지 않으면 실행 시점에 "대충 50\%"로 EV를 잡고 다툰다. 연차 예산(176/198/54)과 PV 곡선의 연차 합계를 대조하지 않는 것도 오류. 예상 소요 6분.}
\end{frame}

\begin{frame}{계약 유형은 리스크 배분이다 — 고정가가 "안전"하다는 믿음은 절반만 맞다}
\begin{tbl}[\scriptsize]{L{0.2\textwidth}L{0.13\textwidth}L{0.24\textwidth}L{0.35\textwidth}}
\textbf{유형} & \textbf{원가 리스크} & \textbf{언제 맞는가} & \textbf{온담 P1} \\ \midrule
\textbf{고정가 (FP / LSTK)} & 공급자 & 범위가 정의되어 있을 때 (Class 2 이상) & SI 용역 68.2, 데이터 전환 11.4, 감리 4.2 \\
실비정산 (CPFF / CPIF) & 발주사 & 범위 불확실, 연구·개발 & (없음) \\
T\&M & 공유 & 소규모·단기·범위 미정 & 앱 벤더 API 연동 추가 시 후보 (R-08) \\
GMP & 최대 보증가까지 발주사 & 설계와 시공이 겹칠 때 & (없음, EPC에서 흔함) \\
단가 계약 & 공유 (단가 고정, 수량 실측) & 수량 불확실한 반복 작업 & 클라우드 사용료, 단말 148만 원 × 477대 \\
\end{tbl}
\begin{itemize}
\item 공급자는 리스크를 \textbf{입찰가의 프리미엄} 또는 \textbf{변경 대가의 공격성}으로 돌려준다 — 응찰 198 \ra{} 계약 168, 손익률 6.5\%. Merrow(2023)의 조건(설계 전 계약 + 감액 압박)이 P1에 그대로
\end{itemize}
\keymsg{FIDIC Emerald Book의 "기준선 대비 편차" 방식을 데이터 전환 용역에 이식 — 정본 결함률 기준선 N\%를 두고 초과분만 대가 조정}
\note{§4.4 "계약 유형과 리스크 배분". PMBOK 8판이 조달을 부록으로 내렸으므로 출처는 계약 실무. 한빛디지털은 198억으로 응찰해 168억으로 계약했다 — 30억 어딘가에 프리미엄이 있었거나, 없었다면 손익률을 지키기 위한 변경 대가 청구가 온다(이재현 PM은 이미 그렇게 말했다). Emerald Book(GBR — 발주자가 지반 기준선을 제공하고 실제가 다를 때만 보상)의 패턴을 SI에 이식하면 발주사는 데이터 정제 책임을, 용역사는 전환 기술 책임을 진다 — 워크북 §4 심화 과제. 그림 4-4(계약 스펙트럼)는 시간이 남으면. 예상 소요 6분.}
\end{frame}

\begin{frame}{FP 고정가의 두 부분 · 견적 검산 3단계 · 하도급에서 확인할 셋}
\begin{itemize}
\item FP 고정가의 원가 기준선 — \textbf{고정 부분} 68.2(지급 조건에 따라 현금, FP 진척에 따라 PV) + \textbf{변동 부분} 변경 대가(재원은 컨틴전시 또는 관리예비비). 산정 기준(FP × 단가? 인월? 보정계수?)이 계약서에 없으면 협의체는 매번 처음부터
\item 견적 검산은 \textbf{협상 전에} — 협상 중 발견하면 상대는 "실수"라 하고 그때는 이미 앵커링(6.2절)
\item ① \textbf{산수}(합계·단가·수량·세금) ② \textbf{정의}(FP인가 인월인가, 코어인가 사용자인가, "3년 유지보수 포함"이 3년치인가 1년 × 3인가, 설치·교육 포함?) ③ \textbf{비교}(다른 견적·과거 계약·공개 단가)
\item 상용SW 23.4(API GW 6.8 / MDM 9.2 / APM 4.1 / SAP 애드온 3.3)는 네 벤더 — 강현도 구매팀장과 검산하지 않으면 그 차이는 2027년 유지보수 갱신 때 나타난다
\item 하도급 — 구조를 계약서에서 확인할 수 있는가 / 산출물을 누가 형상관리하는가(2023 감사 지적) / 인력 이탈이 어떻게 통지되는가
\end{itemize}
\keymsg{조달에서 PM이 하는 가장 실용적인 일은 남의 숫자를 검산하는 것이다}
\note{§4.4 "한국 SI의 FP 고정가", "견적 검산", "컨소시엄·하도급". 워크북 §4 항목 8의 두 방향 검산: 12,400 × 55만 원 = 68.2 vs 스프린트 36회 × 1.408억 = 50.7(평균 44명) — 차이 17.5억 = 피크 61명 증분·간접비·이윤(목표 손익률 6.5\% = 4.4억을 빼면 약 13억). 응찰가 198 \ra{} 168의 30억 인하가 수행사 여유를 얇게 만든 구조를 이해해야 변경 협상에서 수행사가 어디에 서 있는지 안다. 61명 중 하도급이 누구인지 모르는 채 프로젝트가 돈다. 예상 소요 5분.}
\end{frame}

\begin{frame}{RACI에서 자원 계획으로 — 소프트웨어 인력은 평준화할 수 없다}
\begin{itemize}
\item 자원 계획 = RACI의 R을 인력 × 기간으로. 작업패키지마다 붙이면 기간별 수요 막대 = \textbf{자원 히스토그램}, 가용을 넘는 구간이 보인다
\item P1의 피크 82명은 구현 스프린트 중반(2026-09\til{}11 [추가 설정]) — 그때 발주사 14명(대부분 겸임)의 실투입이 50\% 아래로 떨어지면(R-05) 결정 지연 시작
\item \textbf{자원 평준화} = 초과 구간의 활동을 여유 안에서 미루거나 기간을 늘리기. 건설의 노무는 대체 가능하므로 작동한다
\item 소프트웨어 엔지니어는 \textbf{대체 불가} — OD-POS를 아는 사람은 한 명, 그 한 명이 인터페이스 문서화·데이터 매핑·리허설 결함 수정에 동시에 필요
\item 그 사람의 히스토그램은 평준화할 수 없으므로 \textbf{그 사람의 시간이 곧 임계 체인}이다(§3.4 CCPM의 자원 제약)
\end{itemize}
\keymsg{흔한 오류 — 히스토그램 없이 "피크 82명"만 말하기 / 소프트웨어 인력을 평준화할 수 있다고 믿기}
\note{§4.4 "RACI에서 자원 계획으로". Day1 1.5절의 경고 반복 — 건설의 자원 평준화는 노무가 대체 가능하기 때문에 작동하지만, 소프트웨어에서 평준화는 곧 지연이다. R-06(OD-POS 지식 보유자 1명)의 대응이 문서화 S2\til{}S6 전진 배치(1.5.5)인 이유가 여기 있다 — 그 사람의 시간이 임계 체인이면 그 시간을 가장 이른 구간에 집중시켜야 한다. 예상 소요 4분.}
\end{frame}

\begin{frame}{온담 P1 — 168.0의 여덟 내역은 근거의 등급이 균질하지 않다}
\begin{tbl}[\scriptsize]{L{0.24\textwidth}L{0.07\textwidth}L{0.31\textwidth}L{0.27\textwidth}}
\textbf{내역} & \textbf{억} & \textbf{근거의 형태} & \textbf{원가 리스크의 위치} \\ \midrule
SI 개발 용역 & 68.2 & FP 12,400 × 55만 원 — \textbf{계약} (FP는 Class 3 \ra{} G1에서 2) & 변경 대가 (R-05·R-07) \\
상용SW + 3년 유지보수 · 인프라 & 23.4 · 23.9 & 벤더 \textbf{견적} 4건(검산 필요) · 클라우드 14.6 + DR 4.2 + 망분리 5.1 & 수량·정의 오류, 갱신 단가 · 사용량, 규제(R-09) \\
POS 단말 · 데이터 전환 & 8.6 · 11.4 & \textbf{산수}(477 × 148만 + 설치 1.54) · 별도 계약(1,842 테이블) & 비호환 단말(R-13) · 정본 후 변경(R-10) \\
시험·감리 · 발주사 PMO & 9.7 · 12.6 & 부하 3.1 + 보안 2.4 + 감리 4.2 · 3명 × 24개월 + 도구 & 재시험 횟수 · 기간 연장(R-11) \\
\textbf{P1 컨틴전시} & 10.2 & 승인 예산의 6.07\% — \textbf{근거 없음, 5.7절에서 만든다} & — \\
\end{tbl}
\keymsg{68.2는 계약, 8.6은 산수, 23.4는 검산 안 된 견적, 10.2는 백분율 — 원가 기준선을 만드는 일은 이 근거의 등급을 맞추는 일이다}
\note{§4.5 "168.0의 여덟 내역과 그 근거의 등급". 견적은 검산하고, 백분율은 리스크 분석으로 바꾸고, 계약은 지급 조건과 변경 대가 기준을 확인한다. 워크북 §4 항목 1이 이 표를 "근거 등급(계약/산수/견적/백분율)" 열로 받는다. 예상 소요 4분.}
\end{frame}

\begin{frame}{계획 원가 145.8의 산수 — 일곱 내역의 합은 157.8이고 12.0은 어디서 오는가}
\begin{itemize}
\item 컨틴전시를 뺀 일곱의 합 = \textbf{157.8}(68.2 + 23.4 + 23.9 + 8.6 + 11.4 + 9.7 + 12.6). 계획 원가는 145.8. 차이 12.0 = 관리예비비 유보액
\item 이사회는 168.0을 여덟 내역으로 정당화했고, 프로그램 규정이 12.0을 원천 유보했으며, PM은 \textbf{일곱 내역을 12.0만큼 줄여서} 145.8을 만들어야 한다
\item (a) \textbf{재추정·이관·유보} — Class 2 정의도로. 상용SW 23.4\ra{}19.4(2·3년차 유지보수 4.0 P5 이관), 인프라 23.9\ra{}20.9(클라우드 사이징 3.0 유보), PMO 12.6\ra{}9.6(17개월 정합 3.0 유보), 데이터 전환 11.4\ra{}9.4(5년 초과 이력 아카이브 2.0 = 범위 결정)
\item (b) \textbf{범위 이연} — 헌장 §13의 이연 가능 범위를 컷오버 후 릴리스로. 단, 고정가에서 범위 축소는 감액 협상 = 수행사 동의 필요
\item (c) \textbf{스폰서에게 재원을 묻는다} — 이관한 약정분이 실제로 P5 예산(13억 + 운영비 27억)에 있는가. 없다면 12.0은 장부상 유보이지 별도 재원이 아니다
\end{itemize}
\keymsg{157.8에서 145.8로 가는 과정이 문서에 없으면 Day3에서 AC가 145.8을 넘는 순간 "원래 157.8이었다"와 "기준선은 145.8이다"가 충돌한다}
\note{§4.5 "계획 원가 145.8의 산수" [추가 설정]. 예시 기준선(워크북 §4 항목 3 조정표)은 (a)의 조합 — 상용SW 유지보수 4.0 P5 이관 / 클라우드 사이징 3.0 유보 / PMO 3.0 하자보수기 유보 / 5년 초과 이력 아카이브 2.0 범위 결정 = 12.0 — 에 조정 성격·승인자·되돌아올 조건을 붙였고, (c)의 질문은 "되돌아올 조건" 열(P5 예산 거부 시 관리예비비 요청, 사이징 초과 시 복귀)로 남는다 — Day1 3.2절 "헌장을 결정 요청서로 쓴다"가 원가 기준선에도 적용 — "12.0을 어디서 뺐는가"가 아니라 "누가 언제 내는 돈으로 다시 정했는가". 이 질문을 G1 전에 하지 않으면 G3에서 한다. 예상 소요 6분.}
\end{frame}

\begin{frame}{145.8의 월별 PV — 봉우리는 사실이고, 봉우리를 남겨야 Day3에서 읽을 수 있다}
\fig[0.55]{fig_4_5_pv_scurve}
\figcap{그림 4-5. 막대 = 월 PV, 굵은 선 = 누적. 봉우리 2026-03\til{}04(라이선스 14.0·선약정 6.0, 누적 25\%)와 2027-01\til{}02(단말·DR·시험, 누적 70\ra{}89\%). 워크북 §4 항목 5의 배분 [추가 설정]}
\keymsg{2026년 103.0 / 2027년 42.8 — CFO의 148억 삭감(R-04)이 오면 첫 타격 지점은 3\til{}4월 봉우리(라이선스 14.0·선약정 6.0), 프리모템 원문 4번이 정확히 그 시나리오}
\note{§4.5 "145.8의 월별 PV" — 워크북 §4 항목 5 [추가 설정 — 배분 규칙: SI는 스프린트 시작 월 배정(설계 0.7 / 구현 1.1 / 통합 1.2 가중 \ra{} 스프린트당 1.25 / 1.96 / 2.14), 상용SW는 라이선스 인도 시점(3월 8.0 API GW·MDM 1차, 4월 6.0, 5월 3.4 SAP 애드온) + 1년차 유지보수 월 0.2, 인프라는 선약정 6.0(4월) + 클라우드 월 0.4 + 망분리·보안 3.0·2.1(6\til{}7월) + DR 2.1·2.1(12\til{}1월), 단말 2026-12\til{}2027-02, 데이터 전환 10개월 0.6\ra{}1.2\ra{}0.8, 감리 게이트별 0.8/1.4/1.4/0.6 + 보안진단 1.0·1.4 + 부하시험 1.5·1.6, PMO 월 0.56]. 누적: 04 36.83(25.3\%), 06 54.61(37.5\%), 09 77.83(53.4\%), 11 91.81(63.0\%, G2), 12 103.00(70.6\%), 2027-02 129.19(88.6\%, G3·컷오버), 05 145.80. 2026년 누적 103.00 vs BC 연차 집행 계획 102.0 — 차이 1.0은 PV(작업 예정)와 현금 집행(검수 후 지급)의 시점 차이로 설명 가능, 정합(헌장 v1.1 수정 요청 10번). 봉우리를 남겨야 4월 SPI 하락이 개발 지연인지 조달 지연인지 구분된다(워크북 §4.6 해설). 예상 소요 5분.}
\end{frame}

\begin{frame}{컨틴전시 10.2의 배정 원칙 · 관리예비비 12.0의 집행 절차}
\begin{columns}[T]
\begin{column}{0.48\textwidth}
\subhead{컨틴전시 — 총액으로 보유하되 리스크별 근거를 적는다}
\begin{itemize}
\item 리스크별로 칸을 나누면 실현되지 않은 몫이 사장되고, 총액만 두면 근거가 없다
\item 등록부에 "R-08 EMV 2.00 \ra{} 잔여 0.60" / 배정표에 "10.2 = 잔여 4.58 \ra{} 배정 4.60 + 미배정 풀 5.60(스프린트 4회분)" [§5.7, 워크북 §4 항목 6]
\item 집행 시 등록부와 컨틴전시 문서에 \textbf{동시에} 기록
\end{itemize}
\end{column}
\begin{column}{0.5\textwidth}
\subhead{관리예비비 — 예외 보고 \ra{} 결재 \ra{} 재기준선}
\begin{enumerate}
\item PM이 사건이 \textbf{기준선 밖}임을 서술(등록부에 없었고, 제외·경계 밖에서 왔고, 컨틴전시로 불가)
\item 영향 분석 + 대안 3 + 권고를 예외 보고 형식으로(헌장 §12, 5영업일)
\item 정미란 10영업일 내 결재 또는 이사회 부의. 5억 초과는 CFO 병행
\item 집행과 동시에 \textbf{기준선 갱신}(새 작업패키지)
\item CCB 월 보고, 내부감사 분기 보고
\end{enumerate}
\end{column}
\end{columns}
\keymsg{2020년 감사 지적 "컨틴전시 집행 승인 근거 부재"에 대한 두 답 — 리스크 ID 기록, 그리고 이 절차}
\note{§4.5 "컨틴전시 10.2를 어느 리스크에 배정하는가"와 "관리예비비 12.0의 집행 절차". 5,000만 원 이상 발주는 내부통제 지침(2026-02-16). 관리예비비 집행은 거의 항상 재기준선을 동반한다. 워크북 §4 항목 6의 배정 논리 — 컨틴전시 대상 잔여 EMV 4.58 \ra{} 배정 4.60 + 미배정 풀 5.60(스프린트 4회분 5.64 상당, 하이브리드의 되돌림 단위) = 10.20, R-01·R-04는 관리예비비 대상으로 분리 — 가 "리스크 ID와 근거가 붙은 돈"의 실체다. 예상 소요 5분.}
\end{frame}

\begin{frame}{3교시 핵심 정리}
\begin{enumerate}
\item AACE Class는 \textbf{정의도}로 결정되고 정확도는 리스크 분석으로 — 입찰 FP는 Class 3\til{}4, G1 확인서는 Class 2. 건축용은 56R-08
\item 예산의 층 — 계획 원가 145.8 + 컨틴전시 10.2(PM의 돈, 리스크 ID) = BAC 156.0; + 관리예비비 12.0(스폰서의 돈, 재기준선) = 168.0. 크기는 EMV·몬테카를로·RCF에서
\item 원가 기준선은 \textbf{작업 시점} 기준의 PV 곡선. PV·약정·현금은 다르고, 자금 조달 한도가 문서에 있어야 CFO의 삭감에 답한다
\item 고정가는 프리미엄 또는 변경 대가로 돌아온다. 견적은 협상 전에 산수·정의·비교. 소프트웨어 인력은 평준화 불가
\item 온담: 근거 등급이 균질하지 않은 여덟 내역, 157.8\ra{}145.8의 12.0 산수(세 길), 2026-04와 2027-01\til{}02의 봉우리, 컨틴전시는 총액 + 리스크별 근거, 관리예비비는 예외 보고 \ra{} 10영업일 \ra{} 재기준선
\end{enumerate}
\keymsg{점심 후 4교시는 바로 실습 ① — 워크북 §1·§2, 헌장 §4·§5·§6 완성 예시본, 템플릿 PDF를 준비한다}
\note{제4장 핵심 정리 1\til{}5. 생각해 볼 질문 2·3(157.8\ra{}145.8의 세 길 중 무엇을 택하고 정미란에게 어떻게 설명할 것인가 / 컨틴전시 10.2가 잔여 노출의 P80 수준이고 그 너머 꼬리(F3 결합)는 덮지 못함을 알았다면 증액을 요청할 것인가, 구조(분리·폴백)로 꼬리를 자를 것인가, 그대로 두고 알리기만 할 것인가)을 점심 화제로 던진다. 점심 후 4교시는 실습 ①이므로 워크북 §1·§2와 빈 템플릿(06·07), 완성 예시본을 준비시킨다. 예상 소요 3분.}
\end{frame}

% =====================================================================
\section{4교시 (90분) — 실습 ①: 범위 기술서 · WBS(+사전) · RTM 골격}
% =====================================================================

\begin{frame}{이 교시에서 하는 것}
\begin{itemize}
\item \textbf{산출물 ⑥ 범위 기술서와 WBS, WBS 사전} — 워크북 §1, 약 50분. 헌장 §5·§6을 받아 인수 기준·경계 조건·L2 배분·WBS 사전 카드
\item \textbf{산출물 ⑦ 요구사항 추적표 골격} — 워크북 §2, 약 35분. 헌장 §4 상위 요구 8개 \ra{} StRS 샘플 25건, MoSCoW·릴리스·시험·검수 열
\item 진행: 과제 지시 \ra{} 작성 요령 \ra{} 흔한 오류 점검 \ra{} 완성 예시본과 대조
\item 준비물: 워크북 §0.3 한 페이지 요약·§0.4 "Day1에서 가져오는 것", §1.4·§2.4 빈 템플릿, §1.5·§2.5 완성 예시본
\item 3\til{}7년차: 각 절의 \textbf{심화 과제}(자기 회사 WBS·요구 목록으로 같은 것) 병행
\end{itemize}
\keymsg{오늘의 워크북에는 계산이 있다 — 옮겨 적기 전에 검산식(X.6)을 한 번 돌려 보라}
\note{워크북 §0.1 사용법. Day1 §0.1의 열 가지 규칙(값·단위·출처, 직책·실명, 흔한 오류 칸을 체크리스트로)은 그대로. 빈 템플릿 인쇄본: 템플릿/PM\_Day2/06\_범위기술서WBS\_빈템플릿.pdf, 07\_요구사항추적표\_빈템플릿.pdf(통합본 00\_Day2\_템플릿팩\_빈양식.pdf, A4 42쪽 — WBS 트리·RTM은 가로). 완성 예시 통합본 00\_Day2\_템플릿팩\_완성예시.pdf(32쪽). "Day1 값을 바꾸지 말고 인용하라" — 계산하다 Day1 값이 틀렸음을 발견하면 고쳐 적지 말고 §6 "Day1 문서로 되돌아가는 것"에 기록해 헌장 v1.1 수정 요청으로 만든다. 예상 소요 3분.}
\end{frame}

\begin{frame}{산출물 ⑥ 범위 기술서와 WBS — 과제 지시}
\begin{itemize}
\item \textbf{과제}: 워크북 §1.4 빈 템플릿 10개 항목 중 \textbf{3(제외·경계 조건) · 4(인수 기준) · 8(WBS 트리 L2·L3) · 9(L2 원가·FP 배분) · 10(WBS 사전 카드 1장)}. 나머지는 §1.5 예시본으로 확인
\item 세 문서는 한 세트 — 범위 기술서 없이 WBS를 그리면 조각의 경계가 사람마다 다르고, WBS 없이 기술서만 있으면 원가·일정을 붙일 곳이 없고, 사전이 없으면 이름표만 붙은 트리
\item 헌장의 한 줄 "3PL WMS 연동을 준실시간 API로 전환(명세 작성 및 P1측 구현)"을 수행사는 FP 420점으로, 물류본부는 "거부하면 P1이 어떻게든"으로, 박준영은 "폴백이 운영 문서에 있는가"로 읽는다
\item 2022년 WMS 고도화가 63\% 축소된 이유 — "무엇을 인도해야 끝인지 아무도 적지 않았다"
\end{itemize}
\keymsg{시간이 부족하면 항목 4(인수 기준)와 9(L2 배분)만 — 이 둘이 RTM의 검수 열과 원가 기준선의 입력이다}
\note{워크북 §1.1\til{}1.2. 온담 예시: L1 = P1, L2 = 헌장 포함 5개 + 프로젝트 관리·통합 검증(여섯), L3 = 작업패키지 43개(스프린트 1\til{}6회 크기 — 관리 단위가 스프린트 1.41억·2주이기 때문). 인프라를 별도 L2로 빼거나 통합 검증을 각 산출물에 분산하는 답도 가능하며, 어느 쪽이든 근거를 적는 것이 학습. 분야별: SI는 기능 중심 WBS라 인프라·데이터 전환·형상 인수가 빠져 계약 분쟁(2023 감사 지적 "위탁 산출물 소스 인수 절차 미비") / EPC는 공종·구역 중심, 원가 코드와 1:1 / NPD는 BOM(제품 구조)과 WBS(작업 구조) 혼동. 표준 위치: PMBOK 8판 Scope 성과영역(적응형에서는 백로그가 범위 기준선 역할 — P1은 둘 다) / 6판 5.3·5.4 / PMI Practice Standard for WBS 3판 / PRINCE2 PBP / ISO 21502 §7.4. 예상 소요 5분.}
\end{frame}

\begin{frame}{작성 요령 ① — 제외·경계 조건 · 인수 기준 (워크북 §1.3 항목 3\til{}4)}
\begin{itemize}
\item \textbf{항목 3} — 제외 7개를 담당(프로젝트·부서·실명)과 함께 옮기고, 경계 조건 4개는 \textbf{P1이 인도하는 것 / 상대가 인도하는 것 / 폴백} 세 칸으로. 폴백이 빈 행은 §5 등록부에 위협으로
\item 온담 폴백: 3PL "야간 배치 2회(21:00·03:00)" / P2 "물류 2차 분리(이사회 결정)" / P4 "규제 버퍼 소진 시 P1 흡수 — 결정은 프로그램" / P5 "임시 운영 절차 부록"
\item \textbf{항목 4} — 헌장 §6 산출물 8개 × 인수 기준 2\til{}4개(참·거짓 판정 가능) + \textbf{판정 방법}(시험·검수·감리·문서) + 인수자·시점. 숫자는 헌장 §3·§4에서(99.5\%, 15분, 12시간)
\item 통합 재고: ① 단일 원장 조회·반영 15분 이내(TC-INV-001\til{}040) ② 3PL 준실시간 또는 폴백 자동 전환 시험 통과 ③ 재고 정확도 분모 100\% SKU 리포트(R-12)
\item 소스·형상 인수: 발주사 저장소 클린 빌드 재현 100\%, SBOM 제출 — 빠뜨리기 쉬운 기준
\end{itemize}
\keymsg{오류 — "P1은 3PL 연동을 안 한다"와 "P1은 3PL 연동 명세를 쓴다"가 같은 문서 안에서 충돌하는 것}
\note{워크북 §1.3 항목 3·4. 항목 1(문서 정보 — v0.9, 2026-03-27 작성, G1 승인 예정, RTM v0.9·FP 기준선 확인서 BL-01)과 항목 2(포함 범위 — 산출물마다 4\til{}6줄, 총 27줄, 헌장 5줄과 RTM 1,180건 사이의 층, "이연 가능(R5)" 표시), 항목 5(가정·제약 — [신규] 가정 둘: 파일럿 12개점 임시 연동 WBS 1.1.7, 5년 초과 이력 아카이브 WBS 1.5.7), 항목 6(변경 규칙 — 동결 2주 전 접수 마감, Must ≤ 60\%, 허용범위 내 PM 승인, FP ±3\% 재확인 3인 서명, R5 이연은 CCB + §13)은 구두로 한 줄씩. 인수 기준의 흔한 오류 — "정상 동작 확인" / 판정 방법 없음 / 소스·형상 인수 기준 누락. 예상 소요 6분.}
\end{frame}

\begin{frame}{작성 요령 ② — WBS 트리 · L2 원가·FP 배분 · WBS 사전 (항목 7\til{}10)}
\begin{itemize}
\item \textbf{항목 7\til{}8} — L1 P1 / L2 산출물 / L3 작업패키지(명사 — "\til{}규격서", "\til{}모듈", "\til{}패키지"). 담당에 \textbf{발주/수행} 구분: 발주사 담당 10개(단말 조달 1.1.5, 단말 조사 1.1.6, 규격·협상 지원 1.3.1, 인프라 1.5.6, 정본 판정 1.5.8, 기준선·RTM·감리 1.6.1\til{}3·8, 소스 인수 1.6.10)
\item 온담 L2: 1.1 POS·매장 / 1.2 주문 허브 / 1.3 통합 재고 / 1.4 가맹 포털 / 1.5 통합 계층·전환·\textbf{공통 기반(인프라)} / 1.6 프로젝트 관리·통합 검증. WP 43 = 8 + 5 + 5 + 4 + 10 + 11
\item \textbf{항목 9} — 교차표(행 L2 × 열 원가 항목). FP 3,100 / 2,900 / 2,300 / 1,500 / 2,600 = 12,400. 원가 25.65 / 15.95 / 12.65 / 8.25 / 64.00 / 19.30 = \textbf{145.80}
\item \textbf{항목 10} — 카드 1장: 산출물 / 인수 기준·판정 방법 / 담당·인수자 / 선행·후행·외부 의존(상대 실명·기한) / 추정(값 + 근거) / RTM ID·R-ID / 가정. 우선순위: \textbf{외부 의존·임계경로·발주 담당}부터
\end{itemize}
\keymsg{L2 원가 합이 145.8이 아니면 WBS 또는 원가 기준선 중 하나가 100\% 규칙을 어긴 것이다}
\note{워크북 §1.3 항목 7\til{}10과 §1.6 검산. SI 배분 × 0.55억/100FP = 17.05 + 15.95 + 12.65 + 8.25 + 14.30 = 68.20. 원가 항목 열 합 68.20 + 19.40 + 20.90 + 8.60 + 9.40 + 9.70 + 9.60 = 145.80(승인 내역 157.8과의 12.0 차이는 §4 조정표). 인프라를 별도 L2가 아니라 1.5에 둔 이유(§1.6 해설) — 헌장·범위 기술서·WBS·FP 기준선이 같은 상위 구조를 가져야 하고, 인프라는 API GW·MDM·SAP 애드온과 같은 층. 그 대신 1.5가 64.0억(44\%)이 되는데 이것은 문제가 아니라 사실 — 옴니채널 화면 몇 개를 빼도 원가가 크게 안 주는 이유를 스폰서가 WBS에서 본다. 예시본 카드 3장: 1.3.3(3PL 연동, 외부 의존), 1.5.8(마스터 정본 판정, 비가역 — A가 이사회이고 박세연이 R인 이유는 이해 충돌), 1.1.6(가맹 단말 조사·배포 패키지, 발주 담당·R-02·R-13 트리거 "비호환 10\%"가 인수 기준 안에). 예상 소요 8분.}
\end{frame}

\begin{frame}{산출물 ⑥ 흔한 오류 점검}
\begin{itemize}
\item 범위 기술서 버전과 FP 기준선 버전이 따로 논다 — 수행사는 계약 부속서의 12,400을, 발주사는 기술서를 범위로 안다
\item 헌장 문장을 그대로 복사 / 반대로 SRS 570건을 늘어놓음(그것은 RTM의 일) / 인수 기준이 "정상 동작 확인"
\item 제외만 적고 경계에서 P1이 하는 일을 안 적음 / 폴백이 비어 상대가 인도하지 않을 때 무한정 대기
\item L2가 "설계하기·개발하기·시험하기" / 인프라·형상 인수·교육 자료가 빠져 상위 합 ≠ 하위 합 / 발주사 담당 WP가 없어 R-05가 일정에 안 보임
\item SI만 배분하고 라이선스·인프라·PMO는 "공통" — L2 합이 68.2 / 사전을 43장 다 쓰려다 아무것도 안 씀 / 추정 근거가 없어 변경 시 "원래 얼마였는지" 다툼
\end{itemize}
\keymsg{점검 질문 — 모든 WP명이 명사인가 / L2 원가 합이 계획 원가와 소수 첫째 자리까지 같은가 / 폴백이 빈 행이 등록부에 있는가}
\note{워크북 §1.3 각 항목의 "흔한 오류" 칸과 §1.7 점검 질문. 20분 작성 후 이 목록으로 자기 초안을 점검하게 한다. 발주사 담당 작업패키지가 없으면 이재현 PM의 "대기는 변경대가"에 반박할 유일한 근거 — "그 결정은 WBS 1.5.8이고 기한은 2026-11-20이며 우리는 지켰다" — 가 없다(§1.6 해설). 예상 소요 3분 (설명) + 작성 시간.}
\end{frame}

\begin{frame}{산출물 ⑥ 예시본 참조 — 워크북 §1.5\til{}1.6}
\begin{itemize}
\item 왜 인프라를 별도 L2가 아니라 \textbf{1.5 안에} 두었는가 — 그 대신 1.5가 64.0억(44\%)이 된 것은 문제가 아니라 사실
\item 왜 발주사 담당 작업패키지를 \textbf{10개}나 두었는가 — 수행사가 대신할 수 없는 발주사의 인도물, R-05의 일정상 근거
\item 왜 옴니채널 고도화를 별도 WP가 아니라 1.2.5 안의 "이연 가능(R5)" 부분으로 — 별도로 빼면 첫 예산 압력에서 사라진다
\item 왜 [신규] 가정 두 개를 헌장 수정 요청으로 — 파일럿 12개점 임시 연동(FP 약 90)은 지금 넣지 않으면 2027-01에 변경대가 협상
\item 왜 WBS 사전 3건을 1.3.3 · 1.5.8 · 1.1.6으로 골랐는가 — 외부 의존·비가역·발주 담당
\end{itemize}
\keymsg{완성 예시 PDF: 템플릿/PM\_Day2/06\_범위기술서WBS\_완성예시.pdf}
\note{5분 예시본 대조. 조의 L2가 예시본(여섯)과 다르다면 그 이유를 말하게 한다 — 정답이 아니라 근거를 적어야 하는 결정. §1.7 기본 과제(2026-06-30 한동석이 "준실시간 불가, 야간 배치 2회 절충, 2027-06 재협상 조항"을 가져왔을 때 §3 3PL 행·§4 통합 재고 인수 기준·사전 1.3.3 카드 재작성 — 인도물 ③ P1측 준실시간 API는 "구현 완료 후 보관"인가 "범위 제외"인가, FP 420점과 헌장 §5 문장이 어떻게 달라지는가, 승인은 RACI 10·17행)는 시간이 남는 조에만. 예상 소요 5분.}
\end{frame}

\begin{frame}{산출물 ⑦ 요구사항 추적표 골격 — 과제 지시}
\begin{itemize}
\item \textbf{과제}: 워크북 §2.4 빈 템플릿 10개 항목 중 \textbf{2(HR 8개 노드) · 3(StRS 행 — 상위 요구 3 가맹 포털의 StRS 5건 이상) · 4(MoSCoW·릴리스) · 8(검수 항목·인수자)}. 나머지는 §2.5 예시본으로
\item RTM은 명세서가 아니다 — 명세서는 요구의 \textbf{내용}, RTM은 요구의 \textbf{행방}(누가 요구 \ra{} 어느 설계 \ra{} 어느 시험 \ra{} 어느 검수)
\item 왜 골격부터 — \textbf{열의 구조와 계층 규칙은 첫 요구가 접수되기 전에} 있어야 한다. HR에 매달리지 않는 요구는 "누가 왜"를 답해야 접수 = R-07의 구조적 대응
\item 2019 스마트오더 종료 보고서의 중단 사유가 세 줄뿐이었던 것은 "어느 요구가 왜 빠졌는가"를 적을 표가 없었기 때문
\item 감리(정하늘)는 G2·G3에서 커버리지를 본다 — 프리모템 원문 15번은 감리 직전에 RTM을 만들 때 일어난다
\end{itemize}
\keymsg{SI에서 RTM은 검수의 법적 근거다 — "요구했다"와 "명세에 없다"가 충돌할 때 출처·변경 이력 열이 판정한다}
\note{워크북 §2.1\til{}2.2. 분야별: 규제 산업(제약·의료기기·금융)에서는 RTM이 규제 제출 문서(GAMP 5, IEC 62304), NPD에서는 QFD(VOC \ra{} 기능 요구 \ra{} 설계 파라미터)가 같은 역할. 공통점 — 추적표의 열이 초기에 정의되지 않으면 나중에 채울 수 없다. 표준 위치: PMBOK 8판 Scope(요구사항 관리는 Scope 안, RTM은 validate scope의 입력) / 6판 5.2·5.5 / ISO/IEC/IEEE 29148(214/396/570의 명칭 출처, 요구의 속성) / PRINCE2에는 RTM이라는 이름이 없고 Product Description의 품질 기준 + Quality register가 그 역할 / Must ≤ 60\%는 DSDM 권고를 P1 계약 구조에 맞춘 것. 예상 소요 5분.}
\end{frame}

\begin{frame}{작성 요령 ③ — HR 노드 · StRS 행 · MoSCoW와 릴리스 (워크북 §2.3 항목 2\til{}4)}
\begin{itemize}
\item \textbf{항목 2 HR} — 헌장 §4 문장 그대로, ID·제공자·유형·\textbf{연결 편익(B-ID)}. HR-01 재고 원장(B-02·B-04) / HR-02 주문·옴니(B-01·B-05) / HR-03 포털(B-08·B-03) / HR-04 배포 창 / HR-05 가용성 / HR-06 규제 / HR-07 자체 운영(감사 5항) / HR-08 SAP
\item \textbf{항목 3 StRS} — "\til{}할 수 있다"의 판정 가능한 문장, 제공자는 \textbf{사람}(S-03 오정근, S-10 서정필), 접수 일자·경로. 해결책("Kafka로 구현")은 요구가 아니다
\item \textbf{항목 4 MoSCoW} — 규칙을 문서 앞에: Must = 규제·가용성·헌장 고정·편익 직접 / Should = 편익 간접 / Could = 편의 / Won't = 헌장 제외 또는 P2·P3·P5 소관. \textbf{Must ≤ 60\%(FP 가중)}. 릴리스 R1\til{}R5는 §3 캘린더와 같은 번호
\item 우선순위는 PM이 아니라 \textbf{Senior User가} 정한다(RACI 5행 A = 오정근). PM은 규칙을 지킨다
\item 예시본 25건: Must 14(56\%) / Should 7 / Could 3 / Won't 1(StRS-054 배달 주문 가맹 배분 — 부속서 협상 P3 소관)
\end{itemize}
\keymsg{HR에 매달리지 않는 요구를 "기타"로 받지 말라 — 기타는 범위 밖이거나 헌장 수정 요청이다}
\note{워크북 §2.3 항목 2\til{}4. 항목 1(문서 정보 — v0.9, 베이스라인 v1.0 2026-03-25 동결분, 4층 HR 8 / StRS 214 / SyRS 396 / SRS 570, ID 규칙 HR-nn·StRS-nnn)은 구두로. 예시본 제공자 분포 — 오정근 7, 박세연 7, 한동석 3, 최영주 3, 박준영 3, 서정필 3, 나영선 1. "현업"이라고 적힌 제공자 없는 요구가 동결 직전에 들어온 요구의 대부분(R-07)이고, 그것을 걸러내는 것이 이 열의 목적. 예상 소요 6분.}
\end{frame}

\begin{frame}{작성 요령 ④ — 하위 추적 · 설계·WBS · 시험 · 검수 · 이력 · 집계 (항목 5\til{}10)}
\begin{itemize}
\item \textbf{5 하위 추적} — "SyRS 4 / SRS 9 (SRS-101\til{}109)". 골격 단계엔 건수만. 0건 = 아직 설계 안 된 요구 \ra{} 집계 미달
\item \textbf{6 설계·WBS} — 문서명·절 번호까지, WBS 코드는 §1과 같은 번호(StRS-031 \ra{} 아키텍처 정의서 §4.2 / 1.3.1, 1.3.2)
\item \textbf{7 시험} — 종류 + ID. 비기능은 다르다: HR-05 \ra{} 장애 전환 TC-DR-001\til{}012, HR-06 \ra{} TC-REG-001\til{}142(항목 1개 = 케이스 1개). \textbf{"검토"는 시험이 아니다}
\item \textbf{8 검수} — "인수 기준 ③ / G3 / 한동석". 시험은 수행사, 검수는 발주사. 검수 열이 비면 G3 커버리지 100\% 불가. \textbf{인수자에 PM 없음}
\item \textbf{9 이력} — CR-nn·일자·승인자·ΔFP 누적, 동결 회차 FZ-nn. 기각도 이력 / \textbf{10 집계} — 연결률 4종 + Must 비율(FP 가중), 동결마다 갱신
\end{itemize}
\keymsg{검산 — 214 + 396 + 570 = 1,180 / ±3\% = ±35건 / FP ±372점 / 샘플 25 = 4+4+3+2+3+3+3+3}
\note{워크북 §2.3 항목 5\til{}10과 §2.6 검산. 예시본 샘플 25건의 SyRS 합 71, SRS 합 125. StRS-034(재고 정확도 리포트 두 분모)는 SyRS 2 / SRS 3으로 작지만 R-12와 직결 — 요구의 크기와 무게는 다르다. 이력 예: CR-03(StRS-072 엑셀 업로드 확장, 2026-03-11, ΔFP +22, 허용범위 내 PM 승인), CR-05(StRS-054 Won't 판정, 2026-03-25, 오정근). 집계: 샘플 기준 연결률 100\%, 전체 214건 기준 하위 분해 71\% [추가 설정] — G1까지 100\%가 통과 조건. 감리 G2 기준 "설계·시험 연결 100\%", G3 기준 "검수 연결 100\%". 예상 소요 6분.}
\end{frame}

\begin{frame}{산출물 ⑦ 흔한 오류 점검}
\begin{itemize}
\item 층을 섞는다 — StRS("한 화면에서 발주할 수 있다")와 SRS("수량 입력은 정수만")가 같은 층에. 그러면 커버리지 100\%가 무의미
\item HR을 새로 쓴다(헌장과 다른 문장) / 매달리지 않는 StRS를 "기타"로 받는다
\item 제공자 칸에 "현업" / 해결책을 요구로("Kafka로") / 전부 Must / PM이 우선순위를 정함
\item 하위 요구가 어느 StRS에서 왔는지 없어 역추적 불가 / "설계서 참조" / WBS 코드 없음
\item 비기능에 "검토" / 검수 열 비움 / 인수자가 PM / 변경 이력을 별도 문서에 / 기각 요구를 지움 / 집계를 감리 직전에 처음 / 건수 비율만 보고 FP 가중 안 봄
\end{itemize}
\keymsg{점검 질문 — 모든 StRS에 사람 이름이 있는가 / 비기능의 시험 열에 시험 ID가 있는가 / Must의 FP 가중 비율을 한 숫자로 말할 수 있는가}
\note{워크북 §2.3 각 항목의 "흔한 오류"와 §2.7 점검 질문. 15분 작성 후 점검. Must의 FP 가중 비율을 말할 수 없다면 헌장 §13의 조정 변수는 아직 작동하지 않는다. 예상 소요 3분 (설명) + 작성 시간.}
\end{frame}

\begin{frame}{산출물 ⑦ 예시본 참조 — 워크북 §2.5\til{}2.6}
\begin{itemize}
\item 왜 Won't 행(StRS-054)을 지우지 않고 남겼는가 — P1이 만들 수 없는 것이 아니라 \textbf{결정할 수 없는 것}. 지우는 RTM은 같은 요구를 세 번 받는다
\item 왜 규제 요구(StRS-121)의 시험 케이스가 142개인가 — 최영주 CPO는 "검토 의견을 제공하지, 결정하지 않는다". 완료 기준이 "142/142"라는 숫자가 된다
\item 왜 Must가 \textbf{61\%}인 상태를 그대로 적었는가 — "G1까지 오정근이 조정"이 헌장 §13을 실제로 작동시키는 방법
\item 왜 StRS-034(두 분모 리포트)를 Must로 — R-12의 대응이 여기 걸려 있다. MoSCoW는 크기가 아니라 무게로
\item 왜 검수 열의 인수자에 PM이 한 명도 없는가 — 검수는 받아서 쓸 사람이 한다. PM이 나타나면 그 요구는 아직 주인을 못 찾은 것
\end{itemize}
\keymsg{완성 예시 PDF: 템플릿/PM\_Day2/07\_요구사항추적표\_완성예시.pdf}
\note{5분 예시본 대조. 조의 MoSCoW 비율이 예시본과 다르다면 그 근거(편익 연결, 규제, 헌장 §13)를 말하게 한다 — 우선순위는 PM이 아니라 Senior User의 결정이므로 "오정근이라면 무엇을 Should로 내릴 것인가"를 조별로 한 항목씩 발표시키면 헌장 §13이 실제로 작동하는 모습이 보인다. §2.7 기본 과제(2026-05-27 동결 FZ-05 직전 주에 영업본부 요구 38건 접수, 그중 22건이 어느 HR에도 매달리지 않음 — 접수인가 기각인가 헌장 수정 요청인가, 각각 누가 결정(RACI 5·6·16행), ±35건·FP ±372점에 어떻게 계산되는가, 접수 마감 규칙이 몇 건에 적용되는가)는 시간이 남는 조에만. 예상 소요 5분.}
\end{frame}

\begin{frame}{4교시 정리 — 두 문서에서 반복된 규칙}
\begin{enumerate}
\item \textbf{같은 것은 같은 번호를 갖는다} — WBS 코드가 RTM의 WBS 열에, 사전 1.3.3 카드의 시험 ID(TC-INV-021\til{}040)가 RTM의 시험 열에, 릴리스 R1\til{}R5가 §3 캘린더에
\item \textbf{합계는 검산된다} — L2 원가 25.65 + 15.95 + 12.65 + 8.25 + 64.00 + 19.30 = 145.80 = 계획 원가, FP 배분 합 12,400, 214 + 396 + 570 = 1,180
\item \textbf{사람은 실명, PM은 인수자가 아니다} — 요구 제공자가 곧 인수자(서정필 \ra{} 1.4 검수, 박준영 \ra{} SAC)
\item \textbf{WBS를 그리다 헌장에 없던 것이 보인다} — 파일럿 12개점 임시 연동, 5년 초과 이력 아카이브 \ra{} 헌장 v1.1 수정 요청(§6)
\item 발주사 담당 작업패키지 10개와 Won't 행 1건 — 안 하는 것과 못 하는 것과 결정할 수 없는 것을 문서로 구분한 것
\end{enumerate}
\keymsg{5교시가 설명하는 것 — 방금 WBS에 넣은 폴백·분리 컷오버·전수 조사가 왜 "리스크 대응 비용은 계획 원가에 있다"의 실체인가}
\note{워크북 §6 "공통 규칙" 셋(합계 검산 / 같은 번호 / 여유·예비·허용범위의 소유자와 정의)을 앞당겨 말한다. 완성하지 못한 항목은 예시본과 대조해 사후에 채우되, L2 배분표(항목 9)만은 6교시 원가 기준선의 입력이므로 조별로 지금 확정시킨다. 5교시는 제5장 리스크 — 대응 조치 13건 중 11건이 방금 그린 WBS 작업패키지 안에 있다는 것(1.3.3 폴백, 1.4.3 선행 릴리스, 1.1.6 조사, 1.5.8 동결 규칙…)이 세 문서가 하나라는 증거이며, 그것이 5.4절 "계획된 대응의 비용은 계획 원가에, 현실화 비용은 컨틴전시에"의 실체다. 예상 소요 3분.}
\end{frame}

% =====================================================================
\section{5교시 (75분) — 리스크: 정량화 · 대응 · 허용범위}
% =====================================================================

\begin{frame}{이 교시에서 배우는 것}
\begin{itemize}
\item 계획 단계의 식별 — RBS · 가정 분석 · 프리모템 원문 회수 · base rate 질문 [§5.1]
\item 정성 분석 — 확률-영향 매트릭스의 네 함정 · 정량 분석 — EMV · 민감도 · 몬테카를로 · 멱함수 꼬리 [§5.2\til{}5.3]
\item 대응 — 위협 5 · 기회 5 · 소유자 · 트리거 · 2차 리스크 · 리스크 예산은 어디에 있는가 [§5.4]
\item 허용범위는 노출의 함수다 · 등록부 vs 이슈 로그 · 리뷰 주기 [§5.5\til{}5.6]
\item 온담 P1 — 13+2건 정량화, EMV 17.44 \ra{} 잔여 6.78 vs 허용범위 8억, 컨틴전시 10.2 = 4.60 + 5.60, 허용범위 7축 확정 [§5.7]
\end{itemize}
\keymsg{정량화의 목적은 정확한 숫자가 아니다 — 순위, 합계, 대응의 가치, 그리고 스폰서에게 결정을 요청할 근거다}
\note{제5장 전체를 75분에. 시간 배분: §5.1 8분, §5.2 8분, §5.3 12분, §5.4 10분, §5.5\til{}5.6 8분, §5.7 20분, 정리 4분, 오류 3분. §5.7의 정량화 표와 허용범위 격자는 6교시 실습 ②(워크북 §5)의 직접 입력이다 — 확률·영향 값은 워크북 §5 예시본과 같은 [추가 설정](대응 전 17.44 \ra{} 잔여 6.78)이며 실습에서 그대로 쓴다. 예상 소요 2분.}
\end{frame}

\begin{frame}{식별 — 착수 때 없던 것이 계획 단계에 생겼다: 작업패키지와 기준선}
\begin{itemize}
\item WBS가 생겼으므로 \textbf{작업패키지 단위}의 리스크가 보이고, 기준선이 생겼으므로 영향을 \textbf{날짜와 금액}으로 적을 수 있다
\item \textbf{RBS} — 점검표("외부 의존" 가지에 P2·3PL·앱 벤더·규제·별도 법인이 각각 있는가) + 집중의 진단(P1은 "외부 의존"과 "조직"에 몰려 있다)
\item \textbf{가정 분석} — 범위 기술서의 가정을 "이 가정이 틀리면"으로 뒤집기. 새 가정(생산성 계수, 지급 조건, 단말 호환 기준, 리허설 3회로 충분, 설 연휴에 UAT 참여자)마다 새 리스크 후보
\item \textbf{프리모템 회수} — 원문 61건 중 15건만 발췌, 나머지 46건은? "설 연휴에 UAT 참여자가 빠진다"는 달력이 생기자 구체적 리스크가 됐다
\item \textbf{Base rate} — "정본 판정 후 마스터 변경이 일어난 비율", "FP 고정가에서 변경 대가가 계약의 몇 \%인 비율" — 답이 없다는 사실을 기록
\end{itemize}
\keymsg{서술 규칙은 같다 — 원인(현재 사실) \ra{} 사건 \ra{} 영향. 계획 단계에서 달라지는 것은 영향의 단위: "컷오버 지연"이 아니라 "A15가 G3 뒤로 밀려 창 이탈, 다음 창 하반기, 2차 컷오버 N억"}
\note{§5.1. Day1 3.5절(가정·구조·base rate, 프리모템) 위에 둘을 더한다. Day1 초기 리스크 목록의 "범주" 열이 RBS의 레벨 1. Flyvbjerg(2021)가 base rate neglect를 저성과의 1차 원인으로 꼽은 것을 기억하라. 원인은 트리거를, 사건은 확률을, 영향은 금액과 일수를 준다. 흔한 오류: RBS 없이 "외부 의존"이 다섯인데 하나로 세기 / 가정 목록 미갱신 / 프리모템 원문을 옮긴 뒤 버리기 / "일정 지연 리스크"처럼 결과로 서술해 원인이 없음. 예상 소요 5분.}
\end{frame}

\begin{frame}{정성 분석 — 확률-영향 매트릭스의 네 함정}
\fig[0.54]{fig_5_2_pi_matrix}
\figcap{그림 5-1. 가로 = 대응 전 확률, 세로 = P1 원가 영향(억 — 등급에 실제 범위를 붙였다). 붉은 파선 = 한동석 축(R-01·R-03·R-12), 초록 파선 = 박세연 축(R-05·R-06·R-10). 확률·영향은 5.7절 예시 [추가 설정]}
\keymsg{점수를 더하지 말고 위치로 읽는다 — 확률 4 × 영향 2와 확률 2 × 영향 4는 같은 8점이지만 전혀 다른 리스크다}
\note{§5.2. 정성 분석의 가치는 정확성이 아니라 속도와 합의(13건을 한 회의에서 상·중·하로). 네 함정: ① 5×5의 착시(서열척도를 비율척도처럼; "매우 큼"은 "큼"의 두 배가 아니라 열 배일 수 있다 — 처방은 위치로 읽고 영향 등급에 억 원·영업일 정의를 붙이는 것: "영향 상 = 원가 2억 초과 또는 일정 10영업일 초과") ② 상관 무시(R-01·R-03·R-12는 한동석의 결정과 이천센터 협조에, R-05·R-06·R-10은 박세연의 시간에 걸려 있다 — 등록부에 "연관 리스크" 열을 두면 진짜 상위 리스크가 보인다; Day1 프리모템 군집화가 이것) ③ 확률의 의미("중 = 30\til{}50\%"는 어디서 왔는가 — 6.5절 Hubbard의 보정, base rate 우선) ④ 근접성·긴급성의 부재(Day1 "근접성"·"다음 검토" 열이 보완). 흔한 오류: "총 리스크 점수 142"를 보고하기 / "상 리스크 3건"이 사실은 원인 하나. 예상 소요 6분.}
\end{frame}

\begin{frame}{EMV — 평균이고 선형이지만 투명하다 · 의사결정 나무 · 민감도}
\begin{itemize}
\item \textbf{EMV = 확률 × 영향}. 합하면 리스크 노출의 평균 = 컨틴전시의 첫 근거(§4.2)
\item 한계 ① \textbf{평균} — 위협은 오른쪽 꼬리가 길다. EMV 합만큼의 컨틴전시는 "평균적으로 충분"이지 특정 확률로 충분한 것이 아니다
\item 한계 ② \textbf{선형} — 10\% × 100억과 100\% × 10억은 같은 EMV지만 전자는 회사를 흔든다. 리스크 허용도가 다른 둘을 같은 숫자로 만든다
\item 그래도 쓰는 이유 — \textbf{투명성}. 스폰서가 "그 확률은 왜 60\%인가"를 물을 수 있다. 몬테카를로는 그 질문을 시뮬레이션 안에 숨긴다
\item \textbf{의사결정 나무} — R-01 "분리 컷오버 설계를 지금 한다(확정 비용 N) vs 안 한다(2차 컷오버 M × 확률)" = G1까지 이사회에 가져갈 결정의 형식 / \textbf{토네이도} — 가장 긴 막대(원가: 변경 대가 R-05·R-07·R-08, 일정: R-01·R-10)의 원인을 줄이는 것이 짧은 막대 열 개보다 낫다
\end{itemize}
\keymsg{흔한 오류 — EMV 합을 컨틴전시로 놓고 "80\% 안전"이라 믿기 / 토네이도 없이 열 개 리스크에 자원을 고르게 쓰기}
\note{§5.3 "EMV", "민감도". 관례: 위협의 EMV를 양의 비용, 기회를 음으로. 민감도의 용도는 "어디에 리스크 대응 자원을 쓸 것인가". 몬테카를로(원가)는 3.4절 일정과 같은 방법 — 상관을 무시하면 P80 과소평가, Hulett의 risk driver 방식이 처방. 다음 슬라이드에서 멱함수 꼬리로. 예상 소요 5분.}
\end{frame}

\begin{frame}{멱함수 꼬리 재방문 — 컨틴전시는 몸통을 덮고, 구조가 꼬리를 자른다}
\fig[0.54]{fig_5_3_exposure_dist}
\figcap{그림 5-2. 잔여 EMV 합 6.78은 평균이며 P50(약 6.4)보다 크다. 컨틴전시 10.2 ≈ P80이 몸통을 덮고, 그 너머는 관리예비비 12.0, 결합 사건(프리모템 F3)의 꼬리는 어떤 유한한 컨틴전시도 못 덮는다. 개념 시뮬레이션 [추가 설정]}
\keymsg{P1의 꼬리 = 1차연도 삭감 + 3PL 결렬 + P2 지연의 결합(F3) — 셋은 독립이 아니고, 결합하면 영향은 합이 아니라 곱에 가깝다}
\note{§5.3 "멱함수 꼬리 재방문". Flyvbjerg 등(2022): IT 원가 초과는 멱함수 분포 — 평균은 꼬리에 끌려 올라가고 분산은 무한할 수 있다. 그런 분포에서 "EMV 합만큼의 컨틴전시"는 대부분의 프로젝트에서 남고 소수에서 턱없이 모자란다; P80을 덮으면 P90이 있고 그 간격은 정규분포보다 훨씬 크다. 처방은 증액이 아니라 상호의존을 끊는 분할(modularity) — P1에서는 이미 나와 있다: 매장·주문/물류 분리 컷오버(R-01), 가맹 포털 선행 릴리스(O-02), 야간 배치 폴백(R-03). "한 부분이 실패해도 다른 부분이 산다"는 구조가 꼬리를 자르는 유일한 방법. 워크북에서 컨틴전시 10.2가 어디까지 덮고 그 너머는 무엇이 덮는지(분리 설계, 관리예비비, 이사회의 중단 결정 — 헌장 §14)를 한 장에 적는다. 예상 소요 5분.}
\end{frame}

\begin{frame}{대응 — 위협 5 · 기회 5, 그리고 P1의 배정}
\begin{tbl}[\scriptsize]{L{0.13\textwidth}L{0.36\textwidth}L{0.13\textwidth}L{0.3\textwidth}}
\textbf{위협} & \textbf{P1의 예} & \textbf{기회} & \textbf{P1의 예} \\ \midrule
\textbf{회피} & R-10: G2 전 마스터 6종 변경 동결을 CCB 의결 — 사건 자체를 막는다 & \textbf{활용} & O-02: 포털을 이연 불가 항목으로, 2026-Q4 선행 릴리스 \\
\textbf{전가} & R-08: 앱 유지보수 계약에 소스 인수·API 연동 조항 & \textbf{공유} & 수행사와 조기 완료 인센티브 \\
\textbf{완화} & R-06: 13본 문서화를 S2\til{}S6에 전진 배치(1.5.5) & \textbf{증대} & O-01: 문서화 전진을 R-06 완화와 통합 \\
\textbf{수용} & R-03: 폴백 설계(계획 원가) + 잔여 수용 & \textbf{수용} & — \\
\textbf{상신} & R-01: P2 리드타임은 프로그램의 것 — PMO장 소유, P1은 분리 설계로 자기 몫만 & \textbf{상신} & — \\
\end{tbl}
\keymsg{대응 전략이 전부 "완화"인 등록부 — 원인을 제거하거나(회피) 밖으로 넘길(전가·상신) 생각을 안 했다는 뜻. P1은 회피 1(R-10) · 전가+완화 1(R-08) · 완화+수용 1(R-03), R-01·R-04는 완화하되 자금원·결정을 프로그램·스폰서로(사실상 상신)}
\note{§5.4 "대응 전략". Hillson(2002)이 기회를 리스크 프로세스에 정식으로 포함시키자고 제안한 이래 표준은 위협과 기회를 대칭으로 다룬다(상신은 6판에서 추가). 기회 칸이 비어 있는 등록부는 프로젝트를 방어적으로 만든다 — Day1에서 O-02가 "이 프리모템에서 나온 가장 실용적인 발견"이었던 것을 기억. R-04도 상신(스폰서 결정) + 완화(이연 시나리오 선제). 그림 5-3(fig\_5\_4\_responses)은 시간이 남으면. 예상 소요 5분.}
\end{frame}

\begin{frame}{소유자 · 트리거 · 2차 리스크 · 리스크 예산은 어디에 있는가}
\begin{itemize}
\item \textbf{소유자} = 대응 계획에 \textbf{서명}하고, 트리거를 감시하고, 대응을 집행할 권한을 가진 사람. Day1 13건 중 PM이 소유자인 것은 R-07·R-13뿐
\item 한동석이 R-03의 소유자라는 것 = 3PL 협상 착수(03)와 폴백 결정(06-30)을 하고 협상 미착수라는 트리거를 보고할 의무 — RACI의 A와 일치해야
\item \textbf{트리거}에 측정 방법과 주기 — "참석률 70\% 미만 2주 연속"(R-05)은 누가 매주 세는가. 감시되지 않으면 이슈가 된 뒤에야 발견
\item \textbf{2차 리스크} — R-02 대응(분리 릴리스)은 두 번 배포 = 매장 부하·통합 결함, R-08 대응(계약 조항)은 벤더가 조항을 거부할 리스크. 별도 행으로, 원 리스크와 연결. \textbf{잔여} EMV 합이 컨틴전시의 근거
\item \textbf{리스크 예산} — 계획된 대응의 비용(문서화 전진, 전수 조사, 폴백 설계, 분리 컷오버 옵션)은 \textbf{계획 원가(작업패키지)}에, 현실화되었을 때의 비용은 \textbf{컨틴전시}에
\end{itemize}
\keymsg{이 구분이 없으면 대응 비용이 컨틴전시에서 나가 컨틴전시가 계획 단계에 이미 반쯤 소진된다}
\note{§5.4 "소유자·트리거·2차 리스크", "리스크 예산". 분리 컷오버 설계는 WBS 1.3.4·1.6.9의 작업패키지(계획 원가)이고, P2가 실제로 늦어 2차 컷오버를 실시하는 비용은 컨틴전시(또는 프로그램의 관리예비비). 워크북 §5.6 해설 — 13건 중 11건의 조치가 WBS 작업패키지 안에 있다(1.3.3 폴백, 1.4.3 선행 릴리스, 1.5.5 복원, 1.1.6 조사, 1.5.8 동결 규칙, 1.6.10 기술 문서…); Day1 프리모템의 즉시 조치가 WBS를 그릴 때 작업패키지가 되었기 때문이며 우연이 아니다. 조건부 대응 6건(트리거 발동 시)만 컨틴전시에. 흔한 오류: 소유자 칸이 전부 PM / 소유자가 대응 계획에 서명하지 않음 / 트리거에 주기 없음 / 2차 리스크를 안 적음 / 기회 칸이 빔. 예상 소요 5분.}
\end{frame}

\begin{frame}{허용범위는 그 계층이 흡수하기로 한 리스크 노출의 크기다}
\begin{itemize}
\item 원가 +6\%(10.08억) = "이 정도 편차는 PM이 리스크 대응으로 흡수한다" \ra{} PM의 컨틴전시(10.2)와 같은 크기여야 논리가 맞는다 — 거의 같지만 145.8 + 10.08 = 155.88 ≠ 156.0, 그래서 확정본은 \textbf{허용범위 = 컨틴전시 전액 10.2(+7.0\%)}
\item 일정 +15영업일 = 프로그램 합류 버퍼 / 범위 ±3\% = FP 재산정 트리거 — 온담의 초안은 우연이 아니라 설계된 것
\item \textbf{리스크 축 "EMV 8억"}은 결과의 편차가 아니라 \textbf{노출 자체의 한계선} — 잔여 EMV 합이 8억을 넘으면 사건이 일어나지 않았어도 예외 보고(PRINCE2 risk tolerance)
\item \textbf{편익 $-$5\%p}는 손익분기 66\%에서 역산 — 100\%에서 66\%까지 34\%p의 여지를 프로젝트($-$5) / 프로그램($-$10) / 이사회(그 너머)로 나눈 것. 5\%p ≈ 5.9억/년
\item 사내 규정 대조 — 리스크 ID를 단 컨틴전시 집행 건이 5억을 넘으면(R-06 3.5 + R-10 3.8이 한 달에 겹치면) CFO 결재 10영업일이 대응을 늦춘다 — 그것 자체가 2차 리스크
\end{itemize}
\keymsg{허용범위 EMV 8억이 있는데 등록부에 EMV가 없으면 그 허용범위는 장식이다 — 리스크 등록부를 정량화해야 하는 거버넌스적 이유}
\note{§5.5. Day1 3.4절의 격자(7축 × 3계층)는 거버넌스 문서였고, 계획 단계에서는 리스크 노출과 대조된다. P1 직접 편익 118억의 실현률이 66\% 아래로 가면 5년 NPV가 음. Day1에서 "이 폭이 적절한지는 스폰서가 결정해야 한다"고 한 것의 근거가 이 역산. 헌장 §12의 "단일 건 5억 초과는 허용범위 내라도 CFO 결재 병행"이 컨틴전시 집행 절차(4.5절)와 맞는지 확인 — R-08 4억, R-01 2차 컷오버 4.5억은 안 넘지만 겹치면 넘는다. 예상 소요 5분.}
\end{frame}

\begin{frame}{리스크 등록부 vs 이슈 로그 · 리뷰 주기 세 층}
\begin{columns}[T]
\begin{column}{0.48\textwidth}
\subhead{리스크는 아직, 이슈는 이미}
\begin{itemize}
\item 현실화되면 등록부 상태를 "발생"으로 바꾸고 이슈 로그에 행을 만든다. 컨틴전시 집행은 \textbf{이슈에 대해}
\item 섞으면 확률 100\%인 "리스크"가 등록부에 남고 아무도 이슈로 관리하지 않는다
\item 이슈 로그는 Day3의 주제
\end{itemize}
\end{column}
\begin{column}{0.5\textwidth}
\subhead{리뷰 주기}
\begin{itemize}
\item \textbf{스프린트마다(2주)} — 트리거 관측값, 근접한 리스크의 상태
\item \textbf{월 1회(동결일 다음 주 PM 회의, CCB 보고)} — 등록부 전체, 컨틴전시 잔액·집행 기록, 잔여 EMV 합 vs 8억
\item \textbf{게이트마다(G1·G2·G3)} — 프리모템 재실시, 등록부 재기준선, 격자 재확인
\end{itemize}
\end{column}
\end{columns}
\keymsg{착수 때 한 번 쓰고 잠자는 등록부를 막는 유일한 방법 — 등록부를 월간 성과 보고의 고정 칸(잔여 EMV / 8억 / 컨틴전시 잔액 / 발생 / 트리거 관측값)으로 만드는 것}
\note{§5.6. 등록부의 "다음 검토" 열이 이 주기를 리스크마다 정한 것. 워크북 §5 항목 9 — 월 1회(동결일 다음 주 PM 회의) + 게이트, 트리거 관측 담당 = 소유자, 잔여 EMV 합 8억 초과 예상 또는 단일 리스크 EMV 2억 초과 시 5영업일 내 스폰서 보고, 신규 리스크는 동결일 접수분과 함께 등록. 워크북 §5.6 해설 — R-04는 판정 시점에 확정 통보가 있었다면 리스크가 아니라 이슈(확률 100\%)이며 등록부에서 이슈 로그로 옮기고 관리예비비 집행 요청이 된다. 예상 소요 4분.}
\end{frame}

\begin{frame}{온담 P1 — 13+2건의 정량화, 대응 전과 후 (워크북 §5 예시본)}
\begin{tbl}[\scriptsize]{L{0.06\textwidth}L{0.34\textwidth}L{0.1\textwidth}ccL{0.16\textwidth}c}
\textbf{ID} & \textbf{원인 \ra{} 사건} & \textbf{소유자} & \textbf{P·영향} & \textbf{EMV} & \textbf{대응} & \textbf{잔여} \\ \midrule
R-01 & P2 리드타임 41주 \ra{} MC 04-02 \ra{} 물류 시험 G3 후 & 프로그램 PMO장 & 60\% · 4.5 & 2.70 & 완화(분리 설계) + \textbf{상신} & 1.20 \\
R-04 & 1차연도 148억 \ra{} 배분 축소 \ra{} 라이선스·선약정 지연 & 정미란 & 60\% · 5.0 & 3.00 & \textbf{상신} + 이연 시나리오 & 1.00 \\
R-05 & 겸임 참석률 \ra{} 대기 \ra{} 변경대가 & 박세연 & 70\% · 2.4 & 1.68 & 완화(SLA 5일, 대체 결정권자) & 0.60 \\
R-08 & 앱 벤더 계약 외 주장 \ra{} 별도 견적 4억 & 박세연 & 50\% · 4.0 & 2.00 & \textbf{전가} + 완화 & 0.60 \\
R-10 & 정본 후 마스터 변경 \ra{} 재매핑 3.8억·6주 & 박세연 & 25\% · 3.8 & 0.95 & \textbf{회피}(동결 의결) & 0.38 \\
\textbf{위협 13건 합} & & & & \textbf{17.44} & & \textbf{6.78} \\
\end{tbl}
\begin{itemize}
\item 기회 O-01 $-$0.40, O-02 $-$0.60(대응 전후 같음). 순 노출 16.44 \ra{} 5.78 — 그러나 \textbf{위협 합과 기회 합은 따로 보고}(O-02와 R-02는 같은 원인)
\end{itemize}
\keymsg{대응 전 17.44는 허용범위 8억의 두 배 — 대응하라는 신호이지 돈을 더 달라는 근거가 아니다}
\note{§5.7 "정량화 — 대응 전과 후" [추가 설정 — 확률·영향은 워크북 §5 예시본과 같다]. 표에 없는 8건: R-02(오정근, 50\%·3.0, 1.50 \ra{} 0.60), R-03(한동석, 40\%·2.0, 0.80 \ra{} 0.48), R-06(박세연, 30\%·3.5, 1.05 \ra{} 0.30), R-07(PM, 60\%·1.5, 0.90 \ra{} 0.40), R-09(최영주, 40\%·1.8, 0.72 \ra{} 0.30), R-11(박준영, 40\%·2.5, 1.00 \ra{} 0.40), R-12(오정근·한동석, 40\%·1.0, 0.40 \ra{} 0.20), R-13(PM, 35\%·2.1, 0.74 \ra{} 0.32). 영향은 P1 원가에 미치는 것만, 일정(97.1 \ra{} 33.0영업일)·편익 영향은 별도 열. 이것이 워크북 §7이 예고한 "EMV 합계가 8억을 넘는 경우" — 대응 전 17.44. 예상 소요 6분.}
\end{frame}

\begin{frame}{잔여 6.78 vs 허용범위 8억 — 여유는 1.22뿐, 그리고 컨틴전시 10.2 = 4.60 + 5.60}
\begin{itemize}
\item 잔여 6.78은 8억 이내 — 그러나 \textbf{여유 1.22는 R-01(1.20) 하나의 잔여}. 확률 하나가 조금만 올라가도 넘고, R-04\ra{}R-05\ra{}R-07의 결합 경로가 있다
\item 자금원으로 가른다 — \textbf{컨틴전시 대상 11건 4.58}(원인이 P1 계층) / \textbf{관리예비비 대상 R-01 1.20 + R-04 1.00 = 2.20}(프로그램·스폰서 계층). PM 컨틴전시로 막으면 나머지 11건이 무방비
\item 리스크 축 확정 — \textbf{잔여} 위협 EMV 합 8억(대응 전 17.44는 어느 프로젝트나 넘는다) + \textbf{단일 잔여 2억 초과 시 스폰서 보고}(합계 아래에 R-04가 숨지 않게)
\item \textbf{10.2 = 배정 4.60 + 미배정 풀 5.60}(스프린트 4회분). 몬테카를로 예시(상관 반영, 평균 6.78): P50 약 6.4 / P80 약 10.1 / P90 약 12.7 \ra{} 10.2 ≈ \textbf{P80}
\item 스폰서의 선택지 — (a) P90까지 +2.5억(관리예비비 = 이사회 조건 2항 충돌) (b) 구조 변경(분리·폴백·선행 릴리스, 이미 대응에) (c) 유지하고 P80 기록. 권고 (b) + (c)
\end{itemize}
\keymsg{P1의 컨틴전시는 다섯 중 네 번은 충분하고 한 번은 모자란다 — 그 한 번이 결합 사건(F3)임을 스폰서에게 보이는 것이 정량화의 목적이다}
\note{§5.7 "세 가지를 읽어 낸다", "컨틴전시 10.2와의 대조" [추가 설정 — 몬테카를로 P-값은 워크북 잔여 EMV 6.78을 평균으로 한 개념값]. R-04의 재원 문제 — CFO의 1차연도 삭감은 P1 원가를 늘리는 리스크가 아니라 현금 곡선을 바꾸는 리스크(4.3절)이므로 PV·현금 곡선의 이연 시나리오로 대응하며, 소유자가 정미란이고 PM은 대응안을 만든다; 잔여 1.00은 이연으로 생기는 종량제 프리미엄·대기·압축 비용. 워크북 §4·§5 예시본은 R-01·R-04를 관리예비비 대상으로 분리한다(잔여 2.20) — 스폰서가 자기 결정의 원가를 자기 예비비에서 보게 하기 위해서. 생각해 볼 질문 2 — "컨틴전시 10.2가 P80이고 그 너머는 구조와 관리예비비가 덮는다는 것을 이사회에 보고해야 하는가. 보고하면 무엇이 일어나고, 안 하면 무엇이 일어나는가". 예상 소요 6분.}
\end{frame}

\begin{frame}{허용범위 7축 × 3계층의 확정 — 값보다 측정 정의, 변경 6건}
\fig[0.54]{fig_5_7_tolerance_final}
\figcap{그림 5-4. 원가 +10.2 = 컨틴전시 전액(+7.0\%) / 일정 +15, 창·규제 0 / 범위 ±35건·FP ±372 / 품질 0.4건/FP + 심각도 1·2 미해결 0 / 리스크 \textbf{잔여} 8억 + 단일 2억 / 편익 $-$5\%p / 지속가능성 0}
\keymsg{일정 +15영업일은 컷오버 외 마일스톤에만 — 컷오버 창과 규제 기한은 절대 제약이므로 허용범위 0}
\note{§5.7 "허용범위 7축 × 3계층의 확정 논리"와 워크북 §5 항목 7\til{}8. 세 계층: 작업패키지(수행사 위임) 발주 집행 패키지 +2\% / SI 패키지 FP ±1\% / +5영업일 / 0.13건/FP / 잔여 EMV 2.7억 / 0; 프로젝트 P1 +10.2억(= 컨틴전시 = +7.0\%) / +15(컷오버 창·규제 0, G1·G2·리허설 +5 보고) / ±3\%(±35건·FP ±372) / 0.4 + 심각도 1·2 미해결 0 / 잔여 8억 + 단일 2억 / $-$5\%p / 0; 프로그램 +12.0억(= 관리예비비 = 승인 예산 168.0) / +30(합류 버퍼 15 포함) / ±6\% / 0.8 / 16억 / $-$10\%p / 0. Day1 초안 대비 변경 6건과 이유(항목 8) — 원가 분모(145.8 + 10.08 = 155.88 ≠ 156.0 \ra{} 컨틴전시 전액), 일정 마일스톤 예외(창 이탈은 15일이 아니라 6개월), 범위 FP 병기(계약 재확인 트리거와 정합), 품질 측정 정의(4,960건이 미해결이면 인수 불가, 발견 총수면 흔한 수준), 리스크 "잔여 + 단일 2억"(대응 전 17.44는 어느 프로젝트나 넘고, 합계 아래 R-04가 숨는다), 지속가능성 정의(개인정보 유출 0·High 미조치 0). 변경 요청서가 G1 기준선 패키지의 일부(헌장 v1.1 수정 요청 4·5·6번)이며 스폰서가 서명하면 격자는 초안에서 확정본이 된다. 예상 소요 5분.}
\end{frame}

\begin{frame}{5교시 핵심 정리}
\begin{enumerate}
\item 계획 단계의 식별은 작업패키지 단위와 기준선 단위(날짜·금액)를 더한다 — RBS, 가정 분석, 프리모템 원문 회수, base rate 질문. 서술은 원인 \ra{} 사건 \ra{} 영향
\item 5×5 매트릭스는 위치로 읽고 점수를 더하지 않는다. 영향 등급에 금액·일수, 상관은 연관 리스크 열, 근접성은 별도 열
\item EMV는 평균·선형·투명. 토네이도가 자원 배분의 근거, 몬테카를로는 상관 반영. 멱함수 꼬리는 컨틴전시가 아니라 \textbf{분할 구조}가 자른다
\item 위협 5·기회 5. 소유자는 서명·감시·집행 권한. 2차 리스크는 별도 행, 잔여 EMV가 컨틴전시 근거. 계획된 대응 비용은 계획 원가에
\item 허용범위 = 그 계층이 흡수하기로 한 노출. 온담: 17.44 \ra{} 6.78(8억 이내, 여유 1.22 — 잔여 기준 + 단일 2억 보고로 확정), 10.2 = 배정 4.60 + 풀 5.60 ≈ P80, 7축 확정에 변경 6건
\end{enumerate}
\keymsg{다음 교시는 바로 실습 ② — 워크북 §3·§4·§5, 빈 템플릿 08·09·10, 그리고 어제의 리스크 목록 13+2건을 준비한다}
\note{제5장 핵심 정리 1\til{}6. 리스크 장에서 가장 자주 틀리는 것 일곱(§5.8) — 등록부가 착수 때 한 번 쓰고 잠자는 것 / 소유자 칸이 전부 PM / 전략이 전부 완화 / EMV 합으로 컨틴전시 증액 요구 / 기회로 위협을 상쇄해 보고 / 허용범위와 노출을 대조하지 않음 / 꼬리를 컨틴전시로 덮으려는 것 — 을 한 줄씩 읽어 준다. 생각해 볼 질문 3(R-01의 소유자가 프로그램 PMO장인데 영향은 P1 컷오버 — 프로그램이 대응하지 않으면 P1 PM은 무엇을 할 수 있는가)과 5(F3 결합 실패의 확률은 개별 확률의 곱인가 그보다 큰가)를 6교시 실습 중 조별 토론 재료로. 예상 소요 3분.}
\end{frame}

% =====================================================================
\section{6교시 (75분) — 실습 ②: 일정 기준선 · 원가 기준선 · 리스크 등록부(정량화)}
% =====================================================================

\begin{frame}{이 교시에서 하는 것}
\begin{itemize}
\item \textbf{산출물 ⑧ 일정 기준선} — 워크북 §3, 약 25분. 제약 일자 · CPM 계산표의 빈칸 · 임계경로와 버퍼 · float 문안
\item \textbf{산출물 ⑨ 원가 기준선과 예비비 구조} — 워크북 §4, 약 25분. 3층 · 157.8\ra{}145.8 조정표 · 월별 PV · 컨틴전시 배정표
\item \textbf{산출물 ⑩ 리스크 등록부 정량화 + 허용범위 격자} — 워크북 §5, 약 20분. 확률·영향·EMV · 대응·잔여 · 격자 확정
\item 셋은 서로를 검산한다 — 활동은 WBS에서, PV 월 배분은 활동 시점에서, 컨틴전시 배정은 잔여 EMV에서, 격자의 원가 축은 3층에서
\item 준비물: §3.4·§4.4·§5.4 빈 템플릿(08·09·10), 완성 예시본, 어제의 리스크 목록 13+2건, 오늘 오전의 WBS L2 배분
\end{itemize}
\keymsg{워크북 예시본의 값은 [추가 설정]이다 — 실습은 예시본 값을 그대로 쓰되, 왜 그 값인지(근거 열)를 말할 수 있으면 된다}
\note{워크북 §3\til{}§5. 빈 템플릿: 템플릿/PM\_Day2/08\_일정기준선\_빈템플릿.pdf, 09\_원가기준선\_빈템플릿.pdf, 10\_리스크등록부\_빈템플릿.pdf(마일스톤·네트워크·CPM·스프린트 캘린더·조정표·PV·컨틴전시·등록부·격자는 가로). 교재·워크북·그림의 [추가 설정] 값은 워크북 §1·§3·§4·§5 예시본(WBS L2 6·WP 43, 네트워크 A01\til{}A20, PV 배분 규칙, 리스크 17.44\ra{}6.78, 컨틴전시 4.60 + 5.60)을 정본으로 한다. 시간이 부족하면 각 산출물의 "시간이 부족하면" 항목만. 예상 소요 3분.}
\end{frame}

\begin{frame}{산출물 ⑧ 일정 기준선 — 과제 지시}
\begin{itemize}
\item \textbf{과제}: 워크북 §3.4 빈 템플릿 10개 항목 중 \textbf{3(제약 일자) · 5(CPM 계산표 — A15·MG3·A17·A20 행을 손으로 검산) · 6(임계경로·근접 임계·버퍼) · 8(float 문안)}. 나머지는 §3.5 예시본으로
\item 일정 기준선은 간트 한 장이 아니다 — 간트는 결과의 그림, 기준선은 그 그림이 왜 그렇게 그려졌는지의 \textbf{네트워크와 계산}
\item 마일스톤 16개(명절 잠금 2 제외) 중 9개가 외부 고정인데도 기준선이 필요한 이유 — 첫 지연이 왔을 때 "컷오버가 위험한가"에 답하려면
\item \textbf{세 층이 같은 날짜를 가리켜야 한다} — 마일스톤(스폰서) / 활동 네트워크 20개(PM) / 스프린트 캘린더 36회·릴리스 6·동결 12(팀). 통합시험 A14 = 12-07 = S25 시작 = R3 종료 다음 날
\item \textbf{후진 계산은 고정된 날짜에서 시작한다} — G4·G3·규제 ①. 임계경로는 뒤에서부터 결정된다
\end{itemize}
\keymsg{발주사 결정 지연이 어느 경로를 얼마나 밀었는지 증명하지 못하면 변경대가 협상에서 진다 — 둘 다 네트워크가 없으면 목소리 큰 쪽이 이긴다}
\note{워크북 §3.1\til{}3.2. 표준 위치: PMBOK 8판 Schedule 성과영역(예측형 CPM과 적응형 릴리스·이터레이션 계획을 한 도메인에, 하이브리드 연결은 테일러링 판단) / 6판 6.2\til{}6.5 / PMI Practice Standard for Scheduling 3판 / DCMA 14-Point(기준선에는 1\til{}9번) / PRINCE2 Project Plan·Stage Plan(게이트 사이가 Stage) / ISO 21502 §7.6 / SCL Protocol 2판 CP8·CP3. 시간이 부족하면 항목 5의 검산 3개 활동과 항목 8 문안만. 예상 소요 4분.}
\end{frame}

\begin{frame}{작성 요령 ⑤ — 마일스톤 · 제약 일자 · 네트워크 · CPM 계산표 (워크북 §3.3 항목 2\til{}5)}
\begin{itemize}
\item \textbf{항목 1\til{}2} — 달력 규칙(영업일, 작업일 0 = 2026-01-05, 스프린트 1회 = 14달력일 = 10영업일, 공휴일은 예시본 미반영 — 실제는 약 30일). 마일스톤 16개 + 명절 잠금 2개에 제약 유형·연결 활동
\item \textbf{항목 3} — 제약 9개(FNLT 7: G1·규제 ①·G2·리허설 3회차·G3·G4… / SNET 2: P2 MC·컷오버 창) + \textbf{출처와 어기면}. 외부 가정 날짜 3개(3PL 06-30, 부속서 08-14, 앱 벤더 12월)는 제약이 아니라 리스크 트리거
\item \textbf{항목 4} — 요약 활동 20개 + MG3. 의존 27 = FS 22 / SS 4(래그 10·25·30·10) / FF 1. 리드 금지, SF 안 씀. 활동명에 스프린트 범위("구현 R1 S7\til{}S12")
\item \textbf{항목 5} — ES/EF/LS/LF/TF/FF를 작업일 번호로. 검산: A05 ES = max(A01 EF 60, A03 ES 25 + 30) = 60 / MG3 ES = max(294, 294, 275, 180) = 294 = 02-19 / A20 EF = 365 = G4 05-28 / A15 LF = 294, LS = 289 = ES \ra{} TF 0
\end{itemize}
\keymsg{TF 0인 행 12개 = 논리가 만든 임계 10 + 제약이 만든 0(A07 규제, A15 MC) 2 — 둘의 대응은 다르다}
\note{워크북 §3.3 항목 1\til{}5와 §3.6 검산. 최대 TF 20(A08 3PL API), 음수 여유 0. 자유 여유 예: A03 FF = A05 ES 60 $-$ (25 + 30) = 5. 강제 제약은 DCMA 기준 5\% 이하 권고이지만 게이트·규제로 움직이는 프로그램의 요약 네트워크에서는 그 이상이 불가피 — 대신 출처를 적어 "제약인 척하는 희망 날짜"를 걸러 낸다. 흔한 오류: 달력 없이 "일" / 마일스톤을 네트워크에 연결하지 않아 지나가는 날짜 / 내부 목표를 제약으로 / 논리 없는 활동 / 래그로 기간 숨기기 / 전진만 하고 여유 미계산 / 제약이 만든 음수 여유를 못 봄. 예상 소요 6분.}
\end{frame}

\begin{frame}{작성 요령 ⑥ — 임계경로·버퍼 · 캘린더 · float 문안 · DCMA 점검 · 변경 규칙 (항목 6\til{}10)}
\begin{itemize}
\item \textbf{6} — 임계경로 A01\ra{}A05\ra{}A06\ra{}A10\ra{}A14\ra{}A16\ra{}MG3\ra{}A17\ra{}A18\ra{}A20, 360 + 창 대기 5 = 365. 근접 임계: 데이터 경로(TF 5), G1 경로(A04, TF 4). 버퍼는 \textbf{별도 행}: 합류 15(프로그램), 규제 22(P4), G3\ra{}컷오버 4일(P1)
\item \textbf{7} — 36행 캘린더. 릴리스 R0\ra{}G1 / R2\ra{}규제 ① / R3(포털 선행)\ra{}G2 / R4(UAT)\ra{}G3 / R5\ra{}G4. 동결 FZ-01\til{}12 + 접수 마감 2주 전
\item \textbf{8} — 표준 문안(SCL CP8) + 예외 2(합류 버퍼 = 프로그램 PMO장, 규제 버퍼 = P4). 변경협의체가 대기 비용 판정 시 인용
\item \textbf{9} — DCMA 1\til{}9번 값/기준/판정/사유. 예시본은 4개 미달(래그·FS 비율·제약·과다 기간) — 요약 수준의 특성, 상세 일정에서 해소
\item \textbf{10} — G4 +15영업일 / 컷오버 창·규제 기한 0 / G1·G2 +5 초과 예상 시 스폰서 보고 / 재기준선 CCB + 이사회, 최대 2회
\end{itemize}
\keymsg{G3\ra{}컷오버 4일은 여유가 아니라 간격이다 — 여유는 소비할 수 있고, 간격은 소비하면 매장 통보와 패키지 확정이 없어진다}
\note{워크북 §3.3 항목 6\til{}10과 §3.6 해설. 스프린트 검산: S36 종료 2027-05-23 = 2026-01-05 + 14 × 36 $-$ 1일, 36 = 6 + 20 + 10, 동결 12회 = 2026년 매월 마지막 수요일. 임계경로 길이 검산 360 + 5(창 대기) = 365 = END. "왜 G4 앞 여유가 0인 기준선을 그대로 승인받으려 하는가" — 공휴일 반영 시 약 20영업일 증가로 G4 앞 여유가 음수가 되는데 예시본은 숨기지 않고 점검표 7번에 적었다; 실제 v1.0에서는 안정화 A18 40\ra{}30일(박준영 합의) / 이관 A20과 SS 겹침 / 허용범위 +15를 06-18로 이해 중 하나 — 결정은 스폰서, 재료는 계산표. "여유가 0인 계산표를 보여 주지 않으면 스폰서는 결정할 기회를 갖지 못한다." 예상 소요 6분.}
\end{frame}

\begin{frame}{산출물 ⑧ 흔한 오류 점검 · 예시본 참조 — 워크북 §3.5\til{}3.6}
\begin{columns}[T]
\begin{column}{0.5\textwidth}
\subhead{흔한 오류}
\begin{itemize}
\item 캘린더와 네트워크가 따로 — S25 시작일 ≠ 통합시험 시작일
\item 임계경로가 하나라고 가정 — UAT 경로와 P2 경로가 G3에서 만나며 둘 다 0
\item 버퍼를 활동에 조금씩 숨김 / float 문안이 없어 양쪽이 모든 여유를 자기 것으로
\item 점검 기준을 맞추려고 제약을 지워 여유를 만듦
\end{itemize}
\end{column}
\begin{column}{0.48\textwidth}
\subhead{예시본에서 볼 것}
\begin{itemize}
\item 왜 후진 계산이 G4만이 아니라 G3·규제 ①에서도 시작하는가
\item 왜 A15가 5일뿐인가 — 충분하지 않다는 것을 보여 주는 것이 존재 이유
\item 왜 캘린더에 동결일과 접수 마감일을 넣었는가 — 날짜로 들어가야 팀이 지킨다
\end{itemize}
\end{column}
\end{columns}
\keymsg{점검 질문 — 세 층이 같은 날짜를 가리키는가 / "논리가 만든 0"과 "제약이 만든 0"을 구분하는가 / 여유를 누가 쓸 수 있는지 한 문장으로}
\note{워크북 §3.3 흔한 오류, §3.6 해설, §3.7 점검 질문. 완성 예시 PDF: 템플릿/PM\_Day2/08\_일정기준선\_완성예시.pdf. §3.7 기본 과제(T0에서 41주 확정, MC 2027-04-02 — A15 ES는 몇 번 작업일이 되고 MG3와 컷오버 창의 관계는; (a) 5월 창으로 이동 vs (b) 매장·주문 1차 02-26 유지 + 물류 2차 04-23 분리 — 각각 추가·분할할 활동, G4 유지 여부, 결정자 RACI 24행·R-01 소유자)는 시간이 남는 조에만. 예상 소요 4분 + 작성 시간.}
\end{frame}

\begin{frame}{산출물 ⑨ 원가 기준선과 예비비 구조 — 과제 지시}
\begin{itemize}
\item \textbf{과제}: 워크북 §4.4 빈 템플릿 9개 항목 중 \textbf{2(3층 표) · 3(승인 내역 \ra{} 기준선 조정표 — 12.0을 어디서, 누가) · 5(월별 PV — 예시본 17행 중 2026-04·2027-01 봉우리 검산) · 6(컨틴전시 배정표)}. 나머지는 §4.5 예시본으로
\item 원가 기준선은 예산이 아니다 — 예산은 쓸 수 있는 돈의 상한, 기준선은 일정에 따라 인도될 작업의 \textbf{계획 가치(PV)의 누적}
\item "우리 예산이 얼마인가"에 답이 셋(168.0 / 156.0 / 145.8)이고, 헌장 §8이 적은 3층이 실제로 맞물리는지 확인한다 — \textbf{그리고 맞물리지 않는다}: 내역 7항목 합 157.8 ≠ 145.8
\item 12.0의 차이는 이사회 승인 조건 2항 "관리예비비 원천 유보"에서 온다 — 발견하는 것이 첫 번째 일, 어느 항목에서 조정할지 결정·문서화하는 것이 두 번째 일
\item 조정 없이 145.8을 선언하면 첫 달부터 기준선과 계약 금액이 안 맞고, 문서화하지 않으면 윤태호 감사팀장의 질의에 답할 수 없다
\end{itemize}
\keymsg{시간이 부족하면 항목 3(조정표)과 6(컨틴전시 배정표)만 — 이 둘이 §5 등록부의 자금원 열과 격자의 원가 축에 직결된다}
\note{워크북 §4.1\til{}4.2. 컨틴전시 배정 논리(워크북): 잔여 위협 EMV 6.78 = 컨틴전시 대상 4.58(원인이 P1 계층) + 관리예비비 대상 2.20(R-01·R-04 — 원인이 프로그램·스폰서 계층). 4.58 \ra{} 0.1 단위 반올림 4.60을 리스크별 배정, 남는 5.60은 스프린트 4회분(1.41 × 4 = 5.64)의 미배정 풀 — 하이브리드에서 되돌림의 단위가 스프린트이기 때문. 표준 위치: PMBOK 8판 Finance 성과영역(신설 — 원가 기준선과 자금 조달·재무 정당성을 한 도메인에) / 6판 7.2\til{}7.3 / AACE 18R-97·56R-08 / PRINCE2 risk budget·change budget(P1은 배정분 4.60·미배정 풀 5.60으로 흉내) / ISO 21502 §7.7 / Fleming \& Koppelman. 예상 소요 4분.}
\end{frame}

\begin{frame}{작성 요령 ⑦ — 견적 등급 · 3층 표 · 조정표 · L2 배분 (워크북 §4.3 항목 1\til{}4)}
\begin{itemize}
\item \textbf{1} — 고정가 68.2는 등급 대상 아님, 추정부 77.6은 "Class 3 상당"(정확도 범위는 별도 산정 — 실적 부재로 미산정, 배정 논리로 대체)
\item \textbf{2} — 145.8 + 10.2 = 156.0, + 12.0 = 168.0 + 각 층의 \textbf{변경 절차}. 컨틴전시 = PM 승인·재경 확인·CCB 월 보고(RACI 19행), 관리예비비 = 정미란(20행)
\item \textbf{3 조정표} — 승인 내역 / 조정 / 기준선 / 성격 / 승인자 / \textbf{되돌아올 조건}. 상용SW 유지보수 4.0 \ra{} P5 이관 / 클라우드 사이징 3.0 \ra{} 유보 / PMO 17개월 정합 3.0 \ra{} 유보 / 5년 초과 이력 2.0 \ra{} 범위 결정 = 12.0
\item 조정은 숫자를 줄이는 것이 아니라 \textbf{"누가 언제 내는 돈인지를 다시 정하는 것"} — 68.2를 건드리거나(불가) 인프라를 몰래 줄이거나(99.5\% 위반) 감리를 줄이면(감사 지적) 오류
\item \textbf{4} — §1 항목 9와 같은 교차표(25.65 / 15.95 / 12.65 / 8.25 / 64.00 / 19.30 = 145.80). 두 표가 다르면 어느 쪽이 틀렸는지 찾는다
\end{itemize}
\keymsg{검산 — 4.00 + 3.00 + 2.00 + 3.00 = 12.00 / 157.80 $-$ 12.00 = 145.80 / 68.20 + 19.40 + 20.90 + 8.60 + 9.40 + 9.70 + 9.60 = 145.80}
\note{워크북 §4.3 항목 1\til{}4, §4.6 해설 "왜 12.0의 차이를 발견이라고 부르는가"와 검산. 상용SW 조정 내역 6.8$-$5.6 = 1.2, 9.2$-$7.6 = 1.6, 4.1$-$3.4 = 0.7, 3.3$-$2.8 = 0.5 \ra{} 4.0 [추가 설정]. 교재 4.5절 (a)와 같은 조합 — 어느 조합이든 "12.0을 누가 언제 내는 돈으로 다시 정했는가"가 표에 보여야 한다. 흔한 오류: 한 층만 적기 / 컨틴전시를 PV에 섞기 / 조정을 PM의 머릿속에 두기 / 등급에서 정확도를 자동으로 끌어오기. 예상 소요 6분.}
\end{frame}

\begin{frame}{작성 요령 ⑧ — 월별 PV · 컨틴전시 배정표 · 관리예비비 절차 · 견적 검산 · 변경 규칙 (항목 5\til{}9)}
\begin{itemize}
\item \textbf{5 PV} — 17행 × 항목 열 + 월 PV + 누적. 배분 규칙 한 줄씩(SI 스프린트 가중, 라이선스 인도 시점, PMO 균등). 봉우리 2026-04(17.69)·2027-01(16.20). 2026년 누적 103.00 ≈ BC 102.0
\item \textbf{6 배정표} — 리스크 ID / 잔여 EMV / 배정액 / 트리거 / 결재 / 자금원. 11개 리스크 4.60 + 미배정 풀 5.60(스프린트 4회분) = 10.20. R-01·R-04 잔여 2.20은 \textbf{관리예비비 대상}
\item \textbf{7} — 요청(예외 보고 형식) \ra{} 재경팀장 검토 \ra{} 정미란 결재(10영업일) \ra{} CCB 보고 \ra{} 기준선 갱신(BL-0n) \ra{} 감사 사본. 유보 항목의 복귀도 같은 절차
\item \textbf{8 검산} — 12,400 × 55만 원 = 68.2 vs 36회 × 1.408억 = 50.7(평균 44명). 차이 17.5 = 피크 증분·간접비·이윤(손익률 4.4 제외 시 약 13억)
\item \textbf{9} — 원가 허용범위 = \textbf{컨틴전시 전액 10.2(계획 원가의 +7.0\%)}, Day1 "+6\%(10.08)"에서 변경. 컨틴전시 집행은 기준선 변경이 아니고 관리예비비 집행·범위 변경은 기준선 변경
\end{itemize}
\keymsg{PV의 봉우리를 남겨야 R-04가 오면 어느 봉우리가 밀리는지, 4월 SPI 하락이 개발 지연인지 조달 지연인지 보인다}
\note{워크북 §4.3 항목 5\til{}9와 §4.6 해설. "왜 R-01·R-04를 컨틴전시에서 뺐는가" — PM 컨틴전시로 막으면 실현될 때 나머지 11개 리스크가 무방비가 되고, 스폰서는 자기 결정의 원가를 자기 예비비에서 보지 못한다; Day1 RACI 19·20행과 소유자에 원가를 정합시키는 것. "왜 원가 허용범위를 컨틴전시 전액으로" — 10.08과 10.2는 거의 같지만 분모가 다르고 145.8 + 10.08 = 155.88 ≠ 156.0; 0.12의 틈이 "허용범위 내인데 집행 한도 초과"를 논리적으로 허용한다 \ra{} 헌장 v1.1 항목 12 수정 요청(§6). 흔한 오류: 지급 일정을 PV로 / 균등 배분 / "총액 × n\%" / 관리예비비 대상 리스크를 컨틴전시로 / 한 방향만 검산 / 허용범위 분모 불명확. 예상 소요 6분.}
\end{frame}

\begin{frame}{산출물 ⑨ 흔한 오류 점검 · 예시본 참조 — 워크북 §4.5\til{}4.6}
\begin{columns}[T]
\begin{column}{0.5\textwidth}
\subhead{흔한 오류}
\begin{itemize}
\item 조정 없이 145.8을 선언 / 조정을 PM 머릿속에
\item 지급 일정을 PV로 / 145.8 ÷ 17 균등 배분
\item 컨틴전시 = "총액 × n\%" / 등록부와 따로 놂 / 관리예비비 대상을 컨틴전시로
\item 관리예비비 집행 후 기준선 미갱신 / 검산 차이를 "이상하다"로 끝냄 / 허용범위 분모 불명확
\end{itemize}
\end{column}
\begin{column}{0.48\textwidth}
\subhead{예시본에서 볼 것}
\begin{itemize}
\item 왜 12.0의 차이를 "발견"이라고 부르는가
\item 왜 PV 곡선에 봉우리를 남겼는가
\item 왜 R-01·R-04를 컨틴전시에서 뺐는가
\item 왜 미배정 풀을 스프린트 4회분으로 설명했는가
\item 왜 원가 허용범위를 컨틴전시 전액으로 바꾸자고 하는가
\end{itemize}
\end{column}
\end{columns}
\keymsg{점검 질문 — 승인 내역과 기준선의 차이가 조정표로 설명되는가 / 배정액마다 리스크 ID와 트리거가 있는가 / 봉우리가 어느 인도 시점 때문인지 말할 수 있는가}
\note{워크북 §4.3 흔한 오류, §4.6 해설, §4.7 점검 질문. 완성 예시 PDF: 템플릿/PM\_Day2/09\_원가기준선\_완성예시.pdf. §4.7 기본 과제(2026-02-20 CFO 148억 확정, P1 2026년 102\ra{}86억, 삭감 16억 이월 — 항목 5 PV 행·항목 6 관리예비비 대상 행·항목 3 조정표 재작성; 라이선스 2차 6.0·선약정 6.0·DR 4.2 이월 시 A03·A05가 밀리고 임계경로는; 총액 삭감 168\ra{}152이면 §13의 어느 항목이 R5에서 제외되는가)는 시간이 남는 조에만. 예상 소요 4분 + 작성 시간.}
\end{frame}

\begin{frame}{산출물 ⑩ 리스크 등록부 정량화 + 허용범위 격자 — 과제 지시}
\begin{itemize}
\item \textbf{과제}: 워크북 §5.4 빈 템플릿 9개 항목 중 \textbf{2(정량화 — 13+2건 중 6건 이상에 확률·영향·EMV) · 3(대응·소유자·트리거) · 5(EMV 합계 vs 컨틴전시·허용범위 대조) · 7(격자 확정본 — 원가·일정·리스크 축)}. 나머지는 §5.5 예시본으로
\item 어제의 목록에 \textbf{얼마나}(확률·영향·EMV) · \textbf{무엇을 할 것이며}(전략·조치·비용) · \textbf{하고 나면 얼마가 남는가}(잔여) · \textbf{그 돈은 어디서}(컨틴전시 / 관리예비비)를 붙인다
\item 정량화는 정확한 숫자를 만드는 일이 아니다 — 근거는 셋: \textbf{순위}(R-04 3.00 vs R-13 0.74의 네 배) / \textbf{합계}(대응 전 17.44 = 컨틴전시의 1.7배 노출) / \textbf{대응의 가치}(잔여 6.78 = 10.66억의 기대 손실 감소)
\item 격자 확정 = 숫자를 다시 적는 것이 아니라 \textbf{각 축의 측정 정의}(분모·시점·집계)를 붙이고 Day1 이후 발견한 불일치를 고치는 것. 무엇이 왜 바뀌었는지 항목 8에
\end{itemize}
\keymsg{Hubbard — 측정은 불확실성을 없애는 것이 아니라 줄이는 것이다. 확률 60\%와 50\%의 차이를 증명할 수 없어도 숫자를 붙이는 이유가 그것이다}
\note{워크북 §5.1\til{}5.2. 온담에는 과거 IT 투자의 리스크 실현 기록이 없고(그 자체가 감사 지적) 확률의 차이를 증명할 방법도 없다. 분야별: EPC는 QRA(몬테카를로 P50·P80, 컨틴전시 = P80 $-$ 기본 추정)가 표준, SI는 EMV 합계조차 안 하는 경우가 대부분이라 "총액의 n\%"가 되고 감사가 근거를 묻는다, NPD는 FMEA의 RPN. 이 워크북은 몬테카를로 없이 EMV 합계로 시작하되 꼬리의 한계를 항목 5에 명시. 표준 위치: PMBOK 8판 Risk 성과영역(Uncertainty의 후신) / 6판 11.3\til{}11.5 / PRINCE2 Risk practice(owner와 actionee 구분 — "소유자 (조치: 실명)", 7축 tolerance) / ISO 21502 §7.8 / ISO 31000(윤태호가 인용하는 상위 표준) / Hillson \& Simon ATOM, Hubbard. 예상 소요 4분.}
\end{frame}

\begin{frame}{작성 요령 ⑨ — 척도 · 정량화 · 대응·소유자·트리거 · 2차·잔여 (워크북 §5.3 항목 1\til{}4)}
\begin{itemize}
\item \textbf{1 척도} — 확률 \%(5구간 중앙값, 근거 명시), 영향은 원가(억)와 일정(영업일) \textbf{두 축을 따로}, 일정 1일 ≈ 0.14억. 편익은 별도 열(억/년). 1\til{}5점 곱 "25점"은 순위 외 정보가 없다
\item \textbf{2 정량화} — 영향 금액에 출처(롤백 2.1 헌장 §7, 재매핑 3.8 R-10, 앱 벤더 4.0 프리모템 8). 100\%(이슈)·0\%는 리스크가 아니다. 위협 13건 17.44억 / 97.1영업일, 기회 $-$1.00
\item \textbf{3 대응} — 전략 + \textbf{조치를 WBS 코드로}("1.3.3 폴백 모듈, 기준선 포함"). 기준선에 없는 조치 = 조건부 대응 = 컨틴전시 트리거. 소유자 + 조치 실행자. PM 소유 3개 이하. 트리거는 관측 가능 + 시점
\item \textbf{4 2차·잔여} — 2차 리스크 한 줄 + 등록 여부(R-14\til{}R-18). 잔여를 0으로 적지 말 것. 자금원은 \textbf{원인·결정 권한의 계층}으로 — 잔여 6.78 = 컨틴전시 대상 4.58 + 관리예비비 대상 2.20
\end{itemize}
\keymsg{온담 13건 — 회피 1(R-10) · 전가+완화 1(R-08) · 완화 10 · 완화+수용 1(R-03) / 기준선 포함 조치 11, 조건부 6 / 소유자 PM 2건}
\note{워크북 §5.3 항목 1\til{}4와 §5.6 검산. 위협 EMV(대응 전): 2.70 + 1.50 + 0.80 + 3.00 + 1.68 + 1.05 + 0.90 + 2.00 + 0.72 + 0.95 + 1.00 + 0.40 + 0.74 = 17.44. 잔여: 1.20 + 0.60 + 0.48 + 1.00 + 0.60 + 0.30 + 0.40 + 0.60 + 0.30 + 0.38 + 0.40 + 0.20 + 0.32 = 6.78. 교재 5.7절 표와 같은 [추가 설정]임을 말한다 — 예시일 뿐이고 워크북에서 여러분은 스스로 정하고 근거를 적는다. 흔한 오류: 영향을 "상"으로만 / 원가와 일정을 뭉침 / 전략만 적고 조치 없음("완화") / 조치가 기준선에 없어 아무도 안 함 / 트리거가 "상황 악화 시" / 잔여 0 / 컷오버 분리는 "두 번 컷오버" 2차 리스크를 만든다는 것을 무시. 예상 소요 6분.}
\end{frame}

\begin{frame}{작성 요령 ⑩ — EMV 대조 · 일정 노출 · 격자 확정 · Day1 대비 변경 · 검토 주기 (항목 5\til{}9)}
\begin{itemize}
\item \textbf{5 대조} — "컨틴전시 대상 잔여 ≤ 배정 ≤ 컨틴전시"와 "잔여 위협 합 ≤ 8억". 온담: 17.44 \ra{} 6.78(8억 이내, \textbf{단 여유 1.22 = R-01 하나의 잔여}). 4.58 ≤ 4.60 ≤ 10.20
\item \textbf{6 일정 노출} — 단순 합이 아니라 \textbf{임계경로 여부}로. 임계 직접(R-01·R-04·R-05·R-10·R-11) 약 19일 vs 허용범위 15 + 버퍼 15 = 30. 그러나 컷오버 창은 0이고 R-01 하나가 20일
\item \textbf{7 격자} — 7축 × 3계층 + \textbf{측정 정의 열}("+10.2억 = 계획 원가의 +7.0\%, EAC 기준 예상 초과 시") + 마일스톤 예외 + 변경 표시
\item \textbf{8 변경 6건} — 원가 분모 / 일정 마일스톤 예외 / 범위 FP ±372 병기 / 품질(발견 밀도 + 심각도 1·2 미해결 0건) / 리스크 "잔여 8억 + 단일 2억 보고" / 지속가능성 정의
\item \textbf{9} — 월 1회(동결일 다음 주) + 게이트, 트리거 관측 = 소유자, 8억 초과 예상 또는 단일 2억 초과 시 5영업일 내 스폰서 보고
\end{itemize}
\keymsg{원가는 돈으로 흡수되지만 일정은 창으로 흡수되지 않는다 — 이 비대칭이 헌장 §13(일정·원가 고정, 범위 조정)을 숫자로 정당화한다}
\note{워크북 §5.3 항목 5\til{}9와 §5.6 해설. "왜 6.78이 8억 이내인데도 안심하지 말라고 적었는가" — 어느 리스크 하나의 확률이 조금만 올라가도 넘고, 6.78은 평균이며 R-04\ra{}R-05\ra{}R-07의 결합 경로가 있다; "단일 리스크 2억 초과 보고"를 추가한 이유는 합계가 8 아래라는 사실이 R-04(대응 전 3.00)를 숨기지 않게 하기 위해서 — 허용범위는 안심의 기준이 아니라 보고의 기준. "왜 품질 허용범위의 정의를 두 개로 쪼갰는가" — 12,400 × 0.4 = 4,960건이 UAT 종료 시 미해결이면 인수 불가, 시험 중 발견 총 결함이면 흔한 수준; 정의 없이 숫자만 있으면 G3에서 "0.4 이내"와 "인수 불가"가 동시에 참(Day1 편익 분모 문제와 같은 구조). 격자 검산: 145.8 + 10.2 = 156.0 / 1,180 × 3\% = 35.4 / 12,400 × 3\% = 372 / 118 × 5\% = 5.9억/년. 예상 소요 6분.}
\end{frame}

\begin{frame}{산출물 ⑩ 흔한 오류 점검 · 예시본 참조 — 워크북 §5.5\til{}5.6}
\begin{columns}[T]
\begin{column}{0.5\textwidth}
\subhead{흔한 오류}
\begin{itemize}
\item 1\til{}5점 × 1\til{}5점 = "25점 만점" / 척도 정의 없이 상·중·하를 숫자로
\item 대조를 안 함 / 넘는데 그대로 둠 / 맞추기 위해 확률을 낮춤
\item 일정은 "지연 가능"으로만 / 일정 EMV를 합쳐 "97일 지연"(병렬 경로는 더해지지 않는다)
\item 격자에 숫자만 있고 정의가 없음 / 초안을 그대로 확정 / G1에 만들고 G4까지 안 고침
\end{itemize}
\end{column}
\begin{column}{0.48\textwidth}
\subhead{예시본에서 볼 것}
\begin{itemize}
\item 왜 R-04를 60\%로 두고 관리예비비 대상으로 뺐는가 — 확정 통보가 있었다면 이슈다
\item 왜 대응 조치 11/13이 "기준선 포함"인가 — 프리모템의 즉시 조치가 작업패키지가 되었다
\item 왜 일정 노출을 원가와 따로 계산했는가 — 등록부가 헌장 §13을 검산하는 자리
\end{itemize}
\end{column}
\end{columns}
\keymsg{점검 질문 — 잔여 EMV마다 자금원이 있고 합이 §4 배정표와 같은가 / 기준선에 없는 조치는 누가 언제 하는가 / 7축 각각 "무엇을 언제 어떻게 재는가"}
\note{워크북 §5.3 흔한 오류, §5.6 해설, §5.7 점검 질문. 완성 예시 PDF: 템플릿/PM\_Day2/10\_리스크등록부\_완성예시.pdf. §5.7 기본 과제(T0에서 41주 확정, R-01 확률 100\% = 이슈, 이사회가 G1에서 "매장·주문 1차 02-26 유지, 물류 2차 04-23" 의결 — R-01 행을 이슈로 전환하고 2차 리스크 R-14(두 번 컷오버)를 정식 정량화; 잔여 합 vs 8억, 임계 직접 영향, 격자 일정 축의 2차 컷오버 창 허용범위, 관리예비비 집행 요청서의 대안 3개)는 시간이 남는 조에만. 예상 소요 4분 + 작성 시간.}
\end{frame}

\begin{frame}{다섯 문서의 상호 관계 — 서로를 검산한다, 그리고 Day1 문서로 되돌아간다}
\begin{tbl}[\scriptsize]{L{0.12\textwidth}L{0.5\textwidth}L{0.3\textwidth}}
\textbf{관계} & \textbf{무엇이 무엇이 되었는가} & \textbf{같아야 하는 숫자} \\ \midrule
범위 \ra{} 일정 \ra{} 원가 & WP 43개 \ra{} 요약 활동 20개(A05 = 1.5.2·1.5.3·1.1.2·1.4.2), 활동 시점 \ra{} PV 월 배분(A14 12-07 \ra{} 부하시험 3.1은 2027-01\til{}02) & L2 원가 합 = 계획 원가 = PV 누적 = 145.8 \\
RTM ↔ 범위·일정·품질 & WBS 열 = §1 코드, TC 번호 = 사전 1.3.3 인수 기준, 릴리스 = 캘린더 R1\til{}R5, Won't = 제외 항목의 다른 표현 & Must 60\%(FP), ±35건, ±372점 \\
리스크 ↔ 원가·일정·범위 & 잔여 4.58 \ra{} 배정 4.60, R-01·R-04 \ra{} 관리예비비 절차, R-01 근접성 \ra{} A15 TF 0, R-10 \ra{} 데이터 경로 TF 5 + 동결 규칙 & 대응 13건 중 11건이 WBS 안 \\
Day1로 되돌아감 & 헌장 v1.1 수정 요청 \textbf{10건} — 파일럿 연동, 5년 이력, 12.0 조정표 첨부, 원가 허용범위 분모, 격자 예외·정의, 품질 정의, R-01·R-04 관리예비비 대상, RACI 3행 추가, 등록부 갱신, BC 연차 정합 & 결정자: 정미란(CCB)·박세연·재경팀장·정하늘 \\
\end{tbl}
\keymsg{열 개 중 어느 것도 Day1 문서가 틀렸기 때문에 생긴 것이 아니다 — 결정을 분해하고 계산하자 결정이 미처 다루지 않은 곳이 드러난 것이다}
\note{워크북 §6. Day1의 다섯 문서가 서로를 고쳤다면 Day2의 다섯 문서는 서로를 검산한다 — 어느 하나가 어긋나면 나머지 넷 중 하나가 틀린 것. 세 숫자(L2 원가 합, 계획 원가, PV 누적)가 같은 것은 우연이 아니라 100\% 규칙의 결과. 헌장 v1.1 수정 요청 10건의 목록(워크북 §6 표)을 읽어 준다 — 1 헌장 §5·§9 파일럿 12개점 임시 연동(WBS 1.1.7, FP 약 90) / 2 5년 초과 이력 아카이브(조정 2.0의 근거) / 3 §8에 12.0 조정표 첨부 / 4 §12 원가 허용범위를 컨틴전시 전액 10.2로 / 5 §12 일정 예외·범위 FP 병기·리스크 단일 2억·품질·지속가능성 정의 / 6 §3 품질 측정 정의 + 심각도 1·2 미해결 0건 / 7 §10·§8 R-01·R-04 관리예비비 대상 명시 / 8 RACI 27\til{}29행(R5 이연 승인, 유보 복귀·관리예비비 집행, 마스터 동결 의결) / 9 등록부 S-16 정하늘·S-14 조성태 갱신 / 10 BC §5 연차 계획 102.0 vs PV 103.0 정합 기록. Day1 §7의 마지막 문장 — "Day1 문서는 오늘 완성되는 것이 아니라 G1에서 v1.0 \ra{} v1.1로 확정된다" — 가 이 표다. 예상 소요 5분.}
\end{frame}

\begin{frame}{6교시 정리 — 세 문서에서 반복된 규칙}
\begin{enumerate}
\item \textbf{합계는 검산된다} — WBS 원가 합, PV 누적, EMV 합, 임계경로 길이(360 + 5 = 365)는 손으로 다시 더할 수 있어야 한다
\item \textbf{같은 것은 같은 번호를 갖는다} — WBS 코드, TC 번호, 릴리스 번호, 리스크 ID가 다섯 문서에서 같아야 한다
\item \textbf{여유·예비·허용범위에는 소유자와 정의가 있다} — 누구의 여유인지(합류 버퍼 = 프로그램), 어느 리스크의 컨틴전시인지(4.60 + 풀 5.60), 무엇을 재는 허용범위인지(0.4건/FP의 분자·분모·시점). 적히지 않은 숫자는 첫 지연에서 다툼이 된다
\item \textbf{여유 0인 계산표, 12.0의 조정표, 8억의 두 배인 대응 전 EMV 17.44와 여유 1.22뿐인 잔여 6.78을 숨기지 않는다} — 보여 주지 않으면 스폰서는 결정할 기회를 갖지 못한다
\item 리스크 대응 13건 중 11건이 WBS 안에 있다는 것 — 어제의 프리모템 즉시 조치가 오늘 작업패키지가 되었다는 것 — 이 다섯 문서가 하나라는 증거다
\end{enumerate}
\keymsg{마무리 블록 — 왜 우리는 항상 낙관하는가(planning fallacy·앵커링), 왜 CCPM은 여기서 안 되는가(Raz), 왜 평균 컨틴전시는 실패하는가(멱함수)}
\note{워크북 §6 "공통 규칙" 셋 + 오늘의 정직함 셋. 완성하지 못한 항목은 예시본(§3.5·§4.5·§5.5, 통합본 00\_Day2\_템플릿팩\_완성예시.pdf)과 대조하며 사후에 채운다. 마무리 블록은 제6장 학술 배경 — 오늘 인용한 문헌 여덟 갈래를 "왜 그런가 / 언제 틀리는가"의 관점에서. 예상 소요 3분.}
\end{frame}

% =====================================================================
\section{마무리 (30분) — 학술 배경 깊이 읽기 · 읽을 자료 · Day3 예고}
% =====================================================================

\begin{frame}{이 블록에서 배우는 것}
\begin{itemize}
\item 오늘 인용한 학술 문헌 여덟 갈래를 \textbf{왜 그런가 / 언제 틀리는가}의 관점에서 다시 본다 [교재 제6장]
\item 추정에 관한 것 — planning fallacy와 RCF의 절차 · 앵커링 · Hubbard의 보정 · Jørgensen의 소프트웨어 추정 연구
\item 방법에 관한 것 — Goldratt과 Raz 등의 비판 · Reinertsen의 배치와 큐 · Flyvbjerg 2022의 모듈화 · Hillson의 대응 프레임
\item 읽을 자료 — 필수 · 30일 · 90일 · 참고 · 읽지 않아도 되는 것 [§6.9]
\item 오늘 만든 산출물 다섯 가지와 사후 읽기 과제, Day3 예고
\end{itemize}
\keymsg{계획 단계에서 학술 문헌이 답하는 두 질문 — 왜 우리는 항상 낙관하는가, 왜 이 방법은 여기서 안 되는가}
\note{제6장. Day1 제4장과 같은 이유 — 표준이 "이렇게 하라"고만 말하는 지점에서 학술 문헌은 "왜 그런가"와 "언제 틀리는가"를 말해 주며, 계획 단계에서 그 두 질문은 추정과 방법에 관한 것이다. 시간 배분: 6.1\til{}6.8 네 장에 각 4분, 읽을 자료 3분, 산출물·과제·예고 각 3분. 예상 소요 2분.}
\end{frame}

\begin{frame}{planning fallacy와 RCF의 실제 절차 · 앵커링 — 추정은 첫 숫자에 묶인다}
\begin{columns}[T]
\begin{column}{0.52\textwidth}
\subhead{Kahneman \& Tversky(1979) · Flyvbjerg(2006)}
\begin{itemize}
\item Buehler et al.(1994): \textbf{자기 과거 실적을 알려 주어도 예측이 개선되지 않는다} \ra{} outside view는 "참고"가 아니라 \textbf{절차}
\item RCF 3단계 — ① reference class(20\til{}30건, 없으면 없다고 적기) ② \textbf{분포}(P50·P80·꼬리) ③ uplift
\item \textbf{순서가 전부} — 내부 추정을 먼저 만들고 RCF로 검증하면 앵커링 때문에 형식이 된다
\end{itemize}
\end{column}
\begin{column}{0.46\textwidth}
\subhead{Tversky \& Kahneman(1974, Science)}
\begin{itemize}
\item 무관한 룰렛 숫자도 추정을 움직인다
\item 계획 단계의 앵커 — 입찰 198 \ra{} 168, 승인 168 \ra{} 145.8, 벤더 견적의 첫 숫자, \textbf{임원이 먼저 말한 날짜}
\item 처방 — 앵커 전에 자기 추정 먼저 / 첫 숫자를 말하는 순서 통제(Planning Poker의 동시 공개)
\end{itemize}
\end{column}
\end{columns}
\keymsg{오독 — RCF를 "과거 평균 초과율을 곱한다"로 / "우리는 특별해서 reference class가 없다"(uniqueness bias — 그 믿음이 planning fallacy의 작동 방식)}
\note{§6.1\til{}6.2. Kahneman \& Tversky 1979 「Intuitive prediction」의 수록지·페이지는 2차 출처로만 확인 — 인용 시 원문 대조. RCF의 reference class 폭이 관건: 온담 정보화 사업 4건은 너무 좁고, 모든 IT 5,392건은 너무 넓다 — "중견기업·수백억·기존 시스템 대체·복수 벤더·규제 기한" 같은 특성으로 정의. uplift가 나오면 컨틴전시의 세 번째 근거(4.2절)이며 EMV·몬테카를로와 다른 숫자를 주면 그 차이가 내부 등록부의 낙관을 재는 자. Jørgensen \& Sjøberg의 실험 — 고객이 "이 정도면 될 것 같다"고 말한 금액이 전문가 추정을 유의하게 움직였다(6.6절). 읽을 곳: Flyvbjerg 2006 전체(11쪽, arXiv), Buehler 실험 1\til{}2, 『생각에 관한 생각』 23장. 예상 소요 5분.}
\end{frame}

\begin{frame}{Goldratt의 CCPM과 Raz 등의 비판 · Reinertsen — 배치 크기와 큐}
\begin{columns}[T]
\begin{column}{0.5\textwidth}
\subhead{Raz, Barnes \& Dvir(2003, PMJ)의 다섯 비판}
\begin{itemize}
\item 구성 요소 대부분은 새롭지 않다 / "안전여유 50\%" 전제의 실증이 약하다 / 버퍼 크기 산정법이 통계적으로 정당화되지 않는다 / 임계 체인이 안 바뀐다는 가정 / 성공 사례는 "경영진 관심" 효과일 수 있다
\item 실무자에게 — 버리라는 것이 아니라 \textbf{전제가 자기 조직에서 성립하는지 확인하라}. SI·EPC의 다섯 번째 이유 = 여유 회수의 법적 근거 부재
\end{itemize}
\end{column}
\begin{column}{0.48\textwidth}
\subhead{Reinertsen(2009) — 계획에 닿는 셋}
\begin{itemize}
\item \textbf{배치 크기} — 월 1회 동결·2주 스프린트·릴리스 6개(R0\til{}R5)는 각각 배치 결정. 왜 월 1회이고 격주가 아닌가
\item \textbf{큐와 이용률} — 100\%에 가까우면 대기 시간 폭증(M/M/1의 ρ/(1$-$ρ)). R-05의 실체 — 처방은 이용률을 낮추는 것
\item \textbf{지연 비용} — 컷오버 창을 놓치면 편익 6개월 ≈ 59억 + IPO. 헌장 §13의 정량적 근거
\end{itemize}
\end{column}
\end{columns}
\keymsg{케이던스와 배치와 이용률은 "정해진 것"이 아니라 결정이다 — Day3의 흐름·성과 측정에서 본격적으로}
\note{§6.3\til{}6.4. Raz et al. 서지는 2차 출처로 확인 — 원문 대조 권장. CCPM 오독 둘: "여유를 없애는 방법"(위치를 옮기고 소유자를 정하는 것) / 버퍼 소유자를 안 정한 채 소진율 관리 도입(3.3절 float 문제가 이름만 바뀌어 돌아온다). Reinertsen 175개 원칙 중 계획 단계에 닿는 셋. 리틀의 법칙 WIP = 처리율 × 리드타임. SAFe의 WSJF가 지연 비용에서 나왔다. 예상 소요 5분.}
\end{frame}

\begin{frame}{Hubbard의 보정된 추정 · Jørgensen — 소프트웨어 추정 연구가 말하는 것}
\begin{columns}[T]
\begin{column}{0.48\textwidth}
\subhead{Hubbard — How to Measure Anything}
\begin{itemize}
\item "측정할 수 없다"의 대부분은 "측정을 정의하지 않았다"
\item \textbf{보정} — 90\% 구간이 열 번 중 아홉 번 맞는 사람. 대부분은 절반만(과신). 몇 시간 훈련으로 개선
\item 워크북에서 확률마다 "보정됨 / 보정 안 됨"을 표시하는 것이 실습
\item \textbf{정보의 가치(EVI)} — 단말 전수 조사는 정보를 사는 것. 결정을 바꿀 확률로 정당성 판단
\end{itemize}
\end{column}
\begin{column}{0.5\textwidth}
\subhead{Jørgensen(2004 이후 20년)}
\begin{itemize}
\item 전문가 판단이 모델을 이긴다 — \textbf{조건부로}. 처방은 모델 대체가 아니라 \textbf{판단의 구조화}(유사 참조, 앵커 차단, 복수 독립 추정)
\item \textbf{추정의 목적을 분리하라} — 예측 / 입찰 / 목표. 198 \ra{} 168 \ra{} 145.8이 세 답, 셋을 같은 숫자로 만들려는 압력이 planning fallacy의 조직적 형태
\item "최대 X"를 물으면 P70\til{}80을 말하며 P100이라 믿는다 — 확률을 명시해 묻는다
\end{itemize}
\end{column}
\end{columns}
\keymsg{"확률 중 = 30\til{}50\%는 어디서 왔는가"의 답 — 보정된 사람의 판단이거나, 아니면 base rate}
\note{§6.5\til{}6.6. Hubbard 판본(3판 2014)과 Jørgensen 2004 JSS 리뷰 서지는 [검증 필요]. 대부분의 프로젝트는 결정을 바꾸지 않을 정보를 측정하는 데 돈을 쓰고 바꿀 정보는 측정하지 않는다. 3점 추정의 비관값이 충분히 비관적이지 않은 이유(3.1절)가 Jørgensen의 마지막 발견. 예상 소요 4분.}
\end{frame}

\begin{frame}{Flyvbjerg 2022 — 멱함수와 모듈화 · Hillson — 리스크 대응의 프레임}
\begin{columns}[T]
\begin{column}{0.5\textwidth}
\subhead{Flyvbjerg et al.(2022, JMIS) — 함의 셋}
\begin{itemize}
\item IT 5,392건의 초과는 멱함수 — "평균 초과율"은 의미 없는 통계. 메커니즘은 컴포넌트 간 \textbf{상호의존의 연쇄}
\item ① 컨틴전시를 평균 초과율로 잡지 말라 ② 꼬리를 줄이는 것은 \textbf{모듈화} — 분리 컷오버, 폴백, 선행 릴리스, 인터페이스 명세화
\item ③ \textbf{WBS의 구조가 리스크의 구조다} — 1.5가 넷의 공통 의존 = 구조적 리스크. 처방은 1.5의 조기 안정화
\end{itemize}
\end{column}
\begin{column}{0.48\textwidth}
\subhead{Hillson(2002, IJPM)}
\begin{itemize}
\item 기회를 대칭으로 — 활용·공유·증대·수용. PMBOK 5판 이후 채택(상신은 6판)
\item \textbf{"리스크 = 불확실성 × 목표에의 영향"} — 목표가 없으면 리스크를 정의할 수 없다. 같은 불확실성이 어떤 목표엔 위협, 다른 목표엔 기회(O-02 = R-02의 역발상)
\item Risk attitude — 같은 EMV를 다르게 평가. 태도를 명시하지 않으면 격자는 합의가 아니라 우연
\end{itemize}
\end{column}
\end{columns}
\keymsg{"How to Solve Big Problems"(2022) — 왜 태양광·서버랙은 학습곡선을 타고 원전·올림픽은 못 타는가를 같은 논리로}
\note{§6.7\til{}6.8. Flyvbjerg 2022 JMIS 39(3), arXiv 2210.01573; Oxford Review of Economic Policy의 「Bespoke Versus Platform Strategies」가 모듈화 처방의 학술 원전. Hillson 2002 IJPM 20(3), 235-240 [검증 필요]. 5.3절 "확률 10\% × 100억과 100\% × 10억은 EMV가 같지만 다른 리스크"의 문제를 risk attitude가 다룬다. 예상 소요 5분.}
\end{frame}

\begin{frame}{읽을 자료 — 필수 · 30일 · 90일 · 참고 (제6장 6.9)}
\subhead{필수 (과정 전) — Day1 워크북 완성 예시본 5종 / AACE 18R-97 Table 1 / Flyvbjerg(2006) 11쪽 / SCL Core Principle 8}
\begin{tbl}[\scriptsize]{L{0.4\textwidth}L{0.52\textwidth}}
\textbf{30일 안에} & \textbf{왜} \\ \midrule
Flyvbjerg et al.(2022) 멱함수 논문, arXiv 무료 & 평균 컨틴전시가 왜 실패하는지의 실증과 모듈화 처방. 2시간 \\
Buehler, Griffin \& Ross(1994) 실험 1\til{}2 & 자기 실적을 알아도 안 고쳐진다 — outside view가 절차여야 하는 이유 \\
Hulett(2009) risk driver·correlation 장 / Raz et al.(2003) & 상관 무시의 P80 과소평가 / CCPM 전제 검증 \\
Hubbard(2014) calibration·EVI 장 / Hillson(2002) & 보정 훈련과 정보의 가치 / 위협 5·기회 5의 원전 \\
\end{tbl}
\keymsg{읽지 않아도 되는 것 — "6판 ITTO 암기집" / "Class별 ±N\%" 표만 떼어 놓은 2차 자료 / CCPM 성공 사례집 / "P80을 구했다"고만 말하는 도구 백서}
\note{§6.9. 과정 전 제출 과제 하나: 자기 회사 최근 프로젝트 예산이 몇 층이었는지(승인/기준선/컨틴전시/관리예비비), 각 층을 누가 갖고 있었는지 — 층이 하나뿐이었다면 그것이 Day2 오전의 첫 토론 재료. 90일 안에(자기 조직에 적용할 때): Reinertsen(2009) 배치·큐·케이던스 장 / Merrow(2023) 계약 전략 / Goldratt(1997) 원전과 Raz 대조 / PRINCE2 7 Plans·Risk practice / PMBOK 8판 Scope·Schedule·Finance·Risk / ISO/IEC/IEEE 29148 / Flyvbjerg(2021) 편향 10. 참고(워크북을 쓰다가 손을 뻗을 것): ISO 21511 / MIL-STD-881F / AACE 56R-08(56R-11 아님) / DCMA 14-Point [현행 여부 검증 필요] / SCL 2판 전체 / AACE 29R-03 / ISO/IEC 25010:2023 / SW사업 대가산정 가이드 / 소프트웨어 진흥법(과업심의위원회 — 공공만) / ANSI/PMI 19-006-2019 EVM / Kahneman \& Tversky 1979 / Tversky \& Kahneman 1974 / Jørgensen 2004 / Williams 2002. 서지 검증 표기는 리딩리스트를 따른다. 예상 소요 3분.}
\end{frame}

\begin{frame}{오늘 만든 산출물 다섯 가지}
\begin{tbl}[\scriptsize]{cL{0.26\textwidth}cL{0.26\textwidth}L{0.28\textwidth}}
\textbf{\#} & \textbf{산출물} & \textbf{워크북} & \textbf{PRINCE2 7 대응} & \textbf{상태} \\ \midrule
⑥ & 범위 기술서 + WBS + 사전 & §1 & PPD, PBS, Product Descriptions & v0.9 — G1 승인, FP 확인서와 함께 \\
⑦ & 요구사항 추적표 골격 & §2 & PD quality criteria + Quality register & v0.9 — HR 8 / StRS 25 샘플 \\
⑧ & 일정 기준선 & §3 & Project Plan(일정부), Stage Plan & v0.9 — G4 여유 결정 필요 \\
⑨ & 원가 기준선 + 예비비 구조 & §4 & Project Plan(원가부), risk budget & v0.9 — 12.0 조정표, 결정 요청 \\
⑩ & 리스크 등록부 + 격자 & §5 & Risk register, Tolerances & v1.0 — 격자 변경 6건, 헌장 v1.1 요청 \\
\end{tbl}
\begin{itemize}
\item 다섯 문서 모두 \textbf{"결정된 것을 분해하고 숫자로 만든 문서"} — 그 숫자가 어제의 문서에 수정 요청 10건을 만들었다
\end{itemize}
\keymsg{G1에 네 사람이 서명하면 P1은 "무엇에 대해 늦었다고 말할 수 있는 프로젝트"가 된다}
\note{워크북 §0.2 표의 Day2 부분. 완성하지 못한 항목은 예시본(§1.5·§2.5·§3.5·§4.5·§5.5, 통합본 템플릿/PM\_Day2/00\_Day2\_템플릿팩\_완성예시.pdf)과 대조하며 사후에 채운다. 교재 "Day2를 마치며" — 계획은 예측이 아니라 합의이며 합의의 형태는 측정 가능한 기준선: 범위 기준선(기술서 + WBS + 사전 + FP 확인서), 일정 기준선(임계경로·여유·버퍼·제약 일자 + 하이브리드 달력), 원가 기준선(작업 시점 PV, 예비비 구분, 자금 조달 한도), 그 셋을 흔드는 것의 목록(정량화된 등록부), 편차를 누가 흡수하는가의 합의(격자). 예상 소요 3분.}
\end{frame}

\begin{frame}{사후 읽기 과제 (Day3까지)}
\begin{columns}[T]
\begin{column}{0.48\textwidth}
\subhead{읽기 (총 1시간 안)}
\begin{itemize}
\item AACE 18R-97 Table 1 각주 — "정확도 범위는 자동 적용값이 아니다"를 원문에서
\item SCL Core Principle 8 — "별도 약정이 없으면 공용 자원"을 원문에서
\item 교재 §4.5·§5.7 — 오늘 쓴 조정표와 배정표를 다시 읽으며
\end{itemize}
\end{column}
\begin{column}{0.5\textwidth}
\subhead{쓰기 (심화 과제, 3\til{}7년차 — 하나만)}
\begin{itemize}
\item §1.7 자기 회사 WBS — 동사 노드 수 / 어느 노드에도 없는 금액 / 자사 담당 WP 수
\item §3.7 요약 활동 15개의 전진·후진 — TF 0인 활동 수, 외부 고정일 수, float 문장 유무
\item §4.7 자기 예산을 3층으로 — 예비비가 어디에 숨어 있는가, 12개월 PV의 봉우리
\item §5.7 상위 10건의 EMV 합 vs 컨틴전시 — 소유자가 PM인 비율, 자금원이 적힌 비율
\end{itemize}
\end{column}
\end{columns}
\keymsg{생각해 볼 질문 — 제3장 3 / 제4장 3 / 제5장 2 중 하나씩 골라 답을 적어 오라. Day3 첫 10분에 두세 명이 말한다}
\note{읽기 과제는 총 1시간 안. 심화 과제는 네 개 중 하나만. 기본 과제 변형(시간이 남는 조): §3.7 MC 04-02 시나리오 / §4.7 CFO 148억 확정·2026년 102\ra{}86억 / §5.7 R-01 이슈 전환 + R-14 정량화. 제3장 3(MC 04-02로 통지되면 임계경로와 2차 컷오버), 제4장 3(컨틴전시 10.2가 잔여 노출의 P80이고 꼬리는 못 덮음을 알았을 때 무엇을 할 것인가), 제5장 2(P80과 그 너머의 구조·관리예비비를 이사회에 보고해야 하는가). 예상 소요 2분.}
\end{frame}

\begin{frame}{Day3 예고 — 실행과 통제: 기준선이 생겼으므로 편차를 잴 수 있다}
\begin{tbl}[\scriptsize]{L{0.26\textwidth}L{0.66\textwidth}}
\textbf{Day3 산출물} & \textbf{Day2에서 받는 입력} \\ \midrule
⑪ 변경요청서와 영향 분석 & ⑥ WBS 사전(추정·의존) · ⑦ RTM(하위 추적 열 = "이 변경이 무엇을 건드리는가") · ⑩ 격자(±35건·FP ±372·컨틴전시 10.2 = 허용범위 내/초과 판정) \\
⑫ 성과보고 (EVM·이슈·리스크) & ⑨ PV 곡선 17개월(2026-04·2027-01 봉우리가 SPI 해석의 열쇠) · ⑧ CPM의 TF 소진 추적(A15의 0, 데이터 경로의 5) · ⑩ 등록부 월 갱신(잔여 EMV vs 8억) \\
⑬ 조달·계약 변경 기록 & ⑨ 견적 검산(55만 원/FP, 스프린트 비용, 손익률) · ⑥ L2별 FP · ⑧ 여유 소유권 문안 \\
(품질 통제) & ⑦ 시험·검수 열 · ⑩ 품질 정의(발견 밀도 ÷ 12,400, 심각도 1·2 미해결 0건) \\
\end{tbl}
\begin{itemize}
\item Day3에서 만나는 것 — EVM(PV·EV·AC, CPI·SPI, \textbf{Earned Schedule} — SPI가 종료 시점에 1.0으로 수렴하는 결함), 변경 통제(CCB와 협의체의 순서, 재기준선 최대 2회), QA/QC의 분리, 이슈 로그, Reinertsen의 흐름
\end{itemize}
\keymsg{2026년 7월 말, 누적 PV 62.7억 시점에 P1이 어디에 있는지를 여러분이 읽게 된다 — 계획 시점에 잘못 적힌 것은 실행 시점에 고쳐지지 않는다}
\note{워크북 §7 표와 교재 "Day3 예고". 첫 성과보고(2026-05, G1 직후)에서 4월 SPI가 0.8 아래로 나오면 그것이 조달 지연(R-04 실현)인지 설계 지연(R-05)인지를 PV 곡선의 열별로 구분해야 하고, 그 결과에 따라 일정 기준선의 A03·A05 관계나 원가 기준선의 봉우리가 재기준선 후보가 된다 — 오늘 적은 "재기준선 최대 2회, CCB + 이사회" 규칙이 Day3에서 처음 시험된다. EAC가 집행 한도 156.0 초과를 예상하는 순간이 원가 예외 보고 트리거. 조달 변경 기록의 첫 세 건: 상용SW 유지보수 4.0 P5 이관, 클라우드 재구성(2년 선약정 + 종량), 앱 벤더 계약 변경(R-08). Day4(종료·이관)까지 이어지는 것: ⑥ 인수 기준·⑦ 검수 열 \ra{} ⑭ SAC, ⑩ 편익 허용범위 \ra{} ⑯ 편익 실현 원장. 예상 소요 3분.}
\end{frame}
